% 本文件由 LaTeX 排版 Agent 根据有效 translation.json 自动生成。
% author=Sulpitius Severus (c. 363-c. 420); work_id=3505; title=圣史

\chapter{圣史}

% structure: p0001 source=h1 target=omitted reason=page-title-already-rendered-by-parent

卷一\newline{}卷二\par

\SourceRecord{Translated by Alexander Roberts.；Nicene and Post-Nicene Fathers, Second Series；Vol. 11.；Edited by Philip Schaff and Henry Wace.；Buffalo, NY: Christian Literature Publishing Co.,；1894.；<http://www.newadvent.org/fathers/3505.htm>.}

% structure: p0001 source=h1 target=section reason=page-title

\section{圣史（卷一）}

% structure: p0002 source=h2 target=subsection reason=relative-heading-rank; base=h2

\subsection{第一章}

\Emph{我}谨此着手，从世界之初，将圣经所记载之事，按年代顺序及其重要性，简洁地叙述下来，直至我们记忆所及的时代。许多渴望通过简明论著认识神圣事物的人，热切地恳请我承担此工作。我为了实现他们的愿望，不辞辛劳，最终成功地将散见于多卷巨著的内容，汇编成两本小书。同时，在追求简洁的过程中，我几乎没有遗漏任何事实。此外，我认为，在通览从基督受难至宗徒行事的圣史之后，再补充此后发生的事件，并非不妥。因此，我将讲述耶路撒冷的毁灭、基督徒所受的迫害、随后的和平时期，以及因教会内部危险而再次陷入混乱的一切。但我不会羞于承认，在需要之处，我借助了世俗历史学家来确定年代、保持事件序列的连贯，并从中提取完整了解事实所缺的内容，以便既教导无知者，又说服有识者。然而，对于我从圣经中缩编的内容，我并不希望自己以作者的身份出现在读者面前，以致他们忽视我材料的来源，而满足于我所写的内容。我的目标是，熟悉原著的人能在此认出他曾在彼处读到的东西；因为神圣事物的所有奥秘，只有从源头本身才能揭示出来。我现在就要开始叙述了。\par

% structure: p0004 source=h2 target=subsection reason=relative-heading-rank; base=h2

\subsection{第二章}

世界由天主创造，距今约六千年，正如我们将在本书中所述；尽管那些从事年代计算并公布结果的人彼此意见不一。然而，这种分歧或出于天主的旨意，或归咎于古代的谬误，不应成为指责的理由。世界形成后，人受造，男性名为亚当，女性名为厄娃。他们被安置在乐园中，却吃了禁果，因此被驱逐出乐园，流放到我们的地上。他们生了加音和亚伯尔；但加音是个不虔敬的人，杀死了自己的兄弟。加音生了一个儿子，名叫哈诺客，此人建造了第一座城，并以建造者的名字命名。从哈诺客生出依辣得，又从依辣得生出玛胡雅耳。玛胡雅耳生了一个儿子，名叫玛突舍拉；玛突舍拉又生下拉默客，据传拉默客杀死了一个青年，但被杀者的名字未被提及——智者认为这预示了未来的奥秘。亚当在小儿子死后，又生了一个儿子名叫舍特，当时他已二百三十岁；他总共活了八百三十岁。舍特生厄诺士，厄诺士生刻南，刻南生玛拉肋耳，玛拉肋耳生耶勒得，耶勒得生哈诺客，哈诺客因他的正义据说被天主接去。他的儿子名叫玛突舍拉，玛突舍拉生拉默客；从拉默客生出诺厄，他因正义而卓越，且比其他凡人更受天主的喜爱和悦纳。此时，人类已繁衍成众多，一些居于天上的天使，被一些美丽少女的容貌所迷，对她们产生了非分之想，以致堕落，背离了他们本性的原初位置，离开了他们所居住的高天，与人间女子结为婚姻。这些天使逐渐传播邪恶的风俗，败坏了人类家族，据说从他们的结合中生出了巨人，因为与不同性体的结合，自然会产生怪物。\par

% structure: p0006 source=h2 target=subsection reason=relative-heading-rank; base=h2

\subsection{第三章}

天主因这些事，尤其因人类的邪恶已到无可复加的地步而震怒，决意毁灭全人类。但他豁免了诺厄，一个正义且生活无瑕的人，免于注定的厄运。诺厄蒙天主警告洪水将临，便建造了一只巨大的木制方舟，用沥青涂抹使其不透水。他与妻子、三个儿子及三个儿媳一同进入方舟。成对的飞鸟和各种走兽也一同进入，其余一切都被洪水灭绝。诺厄见雨势已停，方舟平稳地漂浮在深渊上，认为水势（正如实际情况）正在消退，便先放出一只乌鸦去探查情况；乌鸦没有返回，我猜想它是落在尸体上；于是他又放出一只鸽子。鸽子找不到落脚之处，便飞回他那里；再次放出后，它带回一片橄榄叶，清楚地表明树梢已可看见。第三次放出后，它不再返回，由此可知洪水已退；诺厄便从方舟出来。据我推算，此事发生在大地肇始后二千二百四十二年。\par

% structure: p0008 source=h2 target=subsection reason=relative-heading-rank; base=h2

\subsection{第四章}

于是诺厄率先为天主筑了一座祭台，从鸟类中献上祭品。紧接着，他与儿子们一同蒙天主祝福，并领受命令：不可吃血，也不可流任何人的血，因为加音没有这样的诫命，却玷污了最初的世界。因此，诺厄的儿子们成了当时空荡荡的世界中仅存的人；他有三个儿子：闪、含和雅弗特。但含因父亲醉酒失态而嘲弄了他，招致了父亲的诅咒。他的儿子名叫雇士，生下了巨人尼默洛得，据说巴比伦城就是由他所建。据说当时还有许多其他城镇建立，此处不打算逐一列举。虽然人类如今繁衍生息，分散居住在不同地区和岛屿，但所有的人都使用同一种语言，直到那日后将散布全世界的群众仍聚居一体之时。这些人按人性的常情，想要在彼此分离之前，通过建造一项伟业来获得大名。于是他们企图建造一座通天之塔。但天主安排，为使工程中的人受阻，他们开始用与惯常言语截然不同的各种语言说话，彼此无法理解。这促使他们更容易分散，因为他们彼此视对方为外人，很容易分离。世界就这样分给了诺厄的儿子们：闪占据东方，雅弗特占据西方，含占据中间地带。此后直到亚巴郎的时代，他们的族谱没有出现什么特别突出或值得记载的事。\par

% structure: p0010 source=h2 target=subsection reason=relative-heading-rank; base=h2

\subsection{第五章}

\Person{亚巴郎}的父亲是特辣黑，生于洪水后第一千零十七年。他的妻子名叫撒辣，起初住在加色丁地区。随后他与父亲同住于哈兰。此时天主对他说话，他便离开故乡和父亲，带着侄儿罗特，来到客纳罕地，在名为舍根的地方定居。不久因粮荒，他下到埃及，后又返回。由于家口众多，罗特与叔父分离，以便在当时空旷的地区占取更广阔的土地，于是在索多玛定居。该城因其居民而臭名昭著，男子强行与男子同寝，据说因此为天主所憎恶。那时，邻近民族的国王们开始武装争斗，而在此之前人类从未有战争。索多玛、哥摩辣及邻近地区的国王们出兵攻打那些侵扰周边区域的敌人，但在首次交锋中败北，将胜利拱手让与对方。于是索多玛被战胜的敌人劫掠，罗特也被掳走。亚巴郎听说此事后，迅速武装了他的仆人，共计三百一十八人，将那些得意洋洋的国王们剥去战利品和武器，击溃了他们。然后，他受到司祭默基瑟德的祝福，并将战利品的十分之一献给了默基瑟德。他将剩余的归还给原主。\par

% structure: p0012 source=h2 target=subsection reason=relative-heading-rank; base=h2

\subsection{第六章}

同时，天主对亚巴郎说话，应许他的后裔将多如海沙；又说他预言的子孙将住在不属于自己的地方，他的后裔在异国将受奴役四百年，但之后将重获自由。接着，天主改变了他和妻子的名字，各加了一个字母：亚巴郎改名为亚巴辣罕，撒辣改名为撒辣。其中蕴含的奥秘绝非微不足道，但并非本作品所应论及。同时，天主规定了割损礼，亚巴郎由婢女生了一个儿子，名叫依市玛耳。此外，当他本人一百岁、妻子九十岁时，天主应许他们将生一个儿子依撒格——主与两位天使一同来到他那里。随后，天使被派往索多玛，看见罗特坐在城门口。罗特以为他们是凡人，便请他们分享款待，在家中设宴招待他们；但城中的恶少却要求将新来者交出，以满足淫秽的目的。罗特提出将女儿们代替客人，但他们不接受，反而更渴求被禁止之事，甚至对罗特本人动了恶念。但天使迅速解救了罗特，使这些不洁的罪人双眼失明。然后，罗特被告知该城将被毁灭，便与妻子和女儿离开了，但被告知不可回头观看。他的妻子却没有遵守这诫命（这是人性偏向邪恶，难以远离禁物），回头看了一眼，据说立刻变成了一根盐柱。至于索多玛，被天火烧成了灰烬。罗特的女儿们以为全人类都已灭亡，便趁父亲醉酒时与他同寝，由此产生了摩阿布和阿孟的种族。\par

% structure: p0014 source=h2 target=subsection reason=relative-heading-rank; base=h2

\subsection{第七章}

几乎同时，当\Person{亚巴郎}已满百岁时，他的儿子\Person{依撒格}诞生了。然后\Person{撒辣}赶走了那使\Person{亚巴郎}生了儿子的婢女；据说她与她的儿子住在旷野，并蒙\Person{天主}的帮助受到保护。此后不久，\Person{天主}试炼\Person{亚巴郎}的信德，要求他将自己的儿子\Person{依撒格}献为祭物。\Person{亚巴郎}毫不犹豫要献上他，已将少年放在祭台上，正拔刀要杀他时，有声音从天而来，命令他宽免那少年；且发现一只公绵羊在近处可作为牺牲。祭献献上后，\Person{天主}对\Person{亚巴郎}说话，向他应许了早已说过要赐予的那些事。但\Person{撒辣}在一百二十七岁时去世，她的遗体由她丈夫悉心安葬在客纳罕人的城镇\Place{赫贝龙}，因为\Person{亚巴郎}正暂居那地方。然后\Person{亚巴郎}见他的儿子\Person{依撒格}已届青年，因为他实际已四十岁，便吩咐他的仆人为他寻找妻子，但只从他本人所属的宗族和地域中找。他受命找到那女子后，要将她带入\Place{客纳罕}地，且不可以为\Person{依撒格}要返回父亲之地娶妻。为使仆人热心执行这些指示，\Person{亚巴郎}让他起誓，将手放在主人（\Person{亚巴郎}）的胯下。仆人于是起程往美索不达米亚，来到\Person{亚巴郎}的兄弟\Person{纳曷尔}的城镇。他进了\Person{纳曷尔}之子、叙利亚人\Person{贝突耳}的家；看见\Person{纳曷尔}的女儿\Person{黎贝加}，一位美丽的处女，便向她求婚，并将她带到主人那里。此后，\Person{亚巴郎}娶了一位名叫\Person{刻突辣}的妻子，在\WorkTitle{编年纪}中她被称为他的妾，并同她生了儿女。但他将财产留给了\Person{撒辣}的儿子\Person{依撒格}，同时将礼物分给了由妾所生的儿子们；这样他们便与\Person{依撒格}分开了。\Person{亚巴郎}享寿一百七十五岁而去世；他的遗体被安放在他妻子\Person{撒辣}的坟墓里。\par

% structure: p0016 source=h2 target=subsection reason=relative-heading-rank; base=h2

\subsection{第八章}

此时，\Person{黎贝加}长久不妊，终于因她丈夫\Person{依撒格}向\Person{主}不住的祈祷，在婚后约二十年产下双胞胎。据说他们在母腹中时常相争；\Person{主}以回答表明，这两个孩子预表两个民族，且将来大的要服事小的。首先出生的，满身是毛，名叫\Person{厄撒乌}；幼子取名\Person{雅各伯}。那时发生了严重饥荒。\Person{依撒格}因这急需，往\Place{革辣尔}去见\Person{阿彼默肋客}王，因受\Person{主}的警告不可下\Place{埃及}。在那里，\Person{主}应许他获得全地，并祝福他，使他牲畜和一切财物大大增多，却因嫉妒被当地居民驱逐。他被逐出那地后，寄居在称为\Quote{誓约井}的井旁。不久，\Person{依撒格}年迈，眼目昏花，准备祝福他的儿子\Person{厄撒乌}时，\Person{雅各伯}受母亲\Person{黎贝加}的指教，前来代替兄弟接受祝福。这样，\Person{雅各伯}在\Person{厄撒乌}之前被立为受诸王和万民尊敬的一位。\Person{厄撒乌}因这些事恼怒，图谋杀害他的兄弟。\Person{雅各伯}因恐惧，又受母亲劝告，逃往\Place{美索不达米亚}，他父亲也催他娶\Person{辣班}（\Person{黎贝加}的兄弟）家的女子为妻：他们寄居异国时，极为关心子女应与本族通婚。这样，\Person{雅各伯}起程往\Place{美索不达米亚}，据说他在梦中看见了\Person{主}的异象；因此将做梦的地方视为圣地，并取了一块石头立作柱子，且许愿说，如果他能平安归来，这柱子的名称该是\Quote{主的殿}，并将所有获得财产的十分之一奉献给\Person{天主}。然后他到母亲的兄弟\Person{辣班}那里，\Person{辣班}和善地接待他，视他为姐妹的儿子，与他同住共享款待。\par

% structure: p0018 source=h2 target=subsection reason=relative-heading-rank; base=h2

\subsection{第九章}

拉班有两个女儿：肋阿和辣黑耳；肋阿双目柔弱，辣黑耳则容貌美丽。雅各伯被她的美貌所迷，热恋这少女，向父亲求亲，甘愿服事七年。但期满后，肋阿被塞给了他，他又被迫再服事七年，之后辣黑耳才被许配给他。然而，据说辣黑耳长久不孕，肋阿却多产。雅各伯从肋阿所生的儿子名字如下：勒乌本、西默盎、肋未、犹大、依撒加尔、则步隆，还有一个女儿狄纳；从肋阿的婢女所生的是加得和阿协尔，从辣黑耳的婢女所生的是丹和纳斐塔里。但辣黑耳在绝望之后生了若瑟。于是雅各伯想回到父亲那里，他的岳父拉班给了他一部分羊群作为服事的报酬，而女婿雅各伯认为岳父在这事上不公义，且怀疑他有欺诈，就在他到达约三十年后悄悄离去了。辣黑耳背着丈夫偷了父亲的神像，因这损害，拉班追赶女婿，但没有找到神像，便和好返回，严令女婿不得再娶妻超过他的女儿。接着，雅各伯上路，据说看见了天使和主军旅的异象。当他行经他兄弟厄撒乌居住的厄东地区时，因疑虑厄撒乌的性情，先派使者并送礼物去试探他。然后他去迎接兄弟，但雅各伯小心不轻易信任他。在兄弟会面的前一天，天主取了人形，据说与雅各伯搏斗。当他赢了天主时，他仍知道对手并非凡人，因此请求祝福。随后，天主打发给他改名，从此雅各伯被称为以色列。但他反过来询问天主的名字时，被告知这名字因奇妙而不可追问。此外，因这场搏斗，雅各伯的大腿窝萎缩了。\par

% structure: p0020 source=h2 target=subsection reason=relative-heading-rank; base=h2

\subsection{第十章}

因此，\Emph{以色列}避开他兄弟的家，打发队伍先到舍根人的城舍根，在那里他买了一块地搭起帐棚。那城的统治者是曷黎人王子厄摩尔。他的儿子舍根玷污了雅各伯从肋阿所生的女儿狄纳。狄纳的兄弟西默盎和肋未发现后，用计谋杀了城中所有男丁，残酷地为妹妹报了仇。城被雅各伯的儿子们洗劫一空，所有战利品被抢走。据说雅各伯对这些事非常不满。随后，受天主指示，他去了贝特耳，在那里为天主立了一座祭台。然后他在加德尔塔附近的地界支搭帐棚。辣黑耳难产而死，所生的男孩名叫本雅明。以色列活到一百八十岁去世。厄撒乌财富众多，并娶了客纳罕民族的女子为妻。在如此简略的著作中，我认为不必提及他的后裔；若有人对此好奇，可查阅原著。父亲去世后，雅各伯继续住在依撒格曾居住的地方。他的其他儿子不时带着羊群外出放牧，但若瑟和小本雅明留在家中。若瑟深得父亲宠爱，因此被兄弟们嫉恨。还有一个原因使他们厌恶：若瑟多次做梦，似乎预示他将比他们都大。于是，当父亲派他去探望羊群和看望兄弟们时，他们觉得是加害他的好机会。一看见弟弟，他们就商议要杀他。但勒乌本心中厌恶这样的罪行，反对他们的计划，于是若瑟被丢进一口井里。后来，因犹大的劝说，他们缓和了主意，把他卖给前往埃及的商人。商人又把他交给法老的宰相普提斐尔。\par

% structure: p0022 source=h2 target=subsection reason=relative-heading-rank; base=h2

\subsection{第十一章}

大约在同一时期，\Person{雅各伯}的儿子\Person{犹大}娶了一位客纳罕女子\Person{萨瓦}为妻。她为他生了三个儿子：\Person{厄尔}、\Person{敖难}和\Person{舍拉}。\Person{厄尔}与\Person{塔玛尔}有婚约。\Person{厄尔}死后，\Person{敖难}娶了哥哥的妻子；据说他因为将精液遗在地上而被天主消灭。于是\Person{塔玛尔}装扮成妓女，与她的兄弟同寝，并为他生了两个儿子。当她分娩时，发生了一件奇特的事：一个男孩出生时，接生婆用一根朱红线绑住他的手，以指明谁是先出生的，但他又缩回母腹，结果成了最后出生的那个男孩。这两个孩子分别命名为\Person{培勒兹}和\Person{则辣黑}。\newline{}至于\Person{若瑟}，由于那位用银钱买下他的国王管家对他很好，并让他管理自己的家业，他因出众的容貌引起了主母的注意。她因那卑劣的情欲而疯狂，多次向他求欢，当他拒绝她的欲望时，她便用虚假的罪名诬陷他，向丈夫抱怨说\Person{若瑟}试图侵犯她的贞洁。因此，\Person{若瑟}被投入监牢。与他一同被关押的有两位国王的仆役，他们向\Person{若瑟}讲述了自己的梦，\Person{若瑟}解释这些梦预示未来，宣布其中一人将被处死，另一人将获赦免。事情果然如此。两年后，国王也做了一个梦。埃及的智者无法解释这梦，那位被释放的国王仆役便告诉国王，\Person{若瑟}是解梦的奇才。于是\Person{若瑟}被带出监牢，为国王解释了他的梦，大意是：未来七年土地将大丰收，但随后七年将出现饥荒。国王被这恐惧所警醒，又见\Person{若瑟}身上有天主的神，便任命他掌管粮食供应，使他在治理上与国王平起平坐。\Person{若瑟}在全埃及谷物丰饶时收集了巨量粮食，通过增加粮仓数目为将来的饥荒做好准备。那时，埃及的希望和安全全系于他一人。大约与此同时，\Person{阿斯纳特}为他生了两个儿子：\Person{默纳协}和\Person{厄弗辣因}。\Person{若瑟}从国王手中接受大权时年三十岁，因为他被哥哥们卖掉时才十七岁。\par

% structure: p0024 source=h2 target=subsection reason=relative-heading-rank; base=h2

\subsection{第十二章。}

其间，埃及为应对饥荒已妥善安排好一切，一场严重的粮荒开始困扰天下。\Person{雅各伯}迫于这种需要，派他的儿子们下到埃及去，只留下\Person{本雅明}在家。此时\Person{若瑟}掌管一切事务，对粮食供应有全权，他的哥哥们来到他面前，向他行了如同对国王一样的敬意。\Person{若瑟}见到他们，却巧妙地隐藏了认出他们的迹象，并指责他们作为敌人前来，诡诈地窥探这地。但他因为没有看到弟弟\Person{本雅明}而烦恼。事情到了这一步：他们承诺会带他前来，特别是为了询问他们是否为了窥探这地而进入埃及。为保证这一承诺的实现，\Person{西默盎}被扣留作为人质，而他们则免费得到了粮食。于是他们回去，按照约定带回了\Person{本雅明}。然后\Person{若瑟}向他的哥哥们显明了自己的身份，使这些恶人蒙羞。他再次打发他们回家，满载粮食并赐予许多礼物，预先警告他们还有五年的饥荒，并劝他们带着父亲、孩子和所有亲属下到埃及去。于是\Person{雅各伯}下到了埃及，令埃及人和国王本人非常欢喜，并受到他儿子热情的欢迎。这事发生在\Person{雅各伯}一百三十岁时，洪水后一千三百六十年。但从\Person{亚巴郎}定居客纳罕地到\Person{雅各伯}进入埃及，可推算为二百一十五年。此后，\Person{雅各伯}在埃及居住的第十七年，因重病恳求\Person{若瑟}将他的遗体安葬在祖坟中。\Person{若瑟}将他的儿子们带来接受祝福；祝福完成时，\Person{若瑟}将幼子置于长子之前，在祝福的价值上也是如此，然后\Person{雅各伯}按顺序祝福了他的所有儿子。他享年一百四十七岁。他的葬礼极其隆重，\Person{若瑟}将他的遗体安放在他父亲的坟墓中。\Person{若瑟}继续善待他的哥哥们，尽管父亲死后，他们因意识到自己的过错而感到恐惧。\Person{若瑟}本人活到一百一十岁去世。\par

% structure: p0026 source=h2 target=subsection reason=relative-heading-rank; base=h2

\subsection{第十三章。}

几乎难以置信，那些下到\Place{埃及}的希伯来人竟然如此迅速地增多，以他们众多的后裔充满了\Place{埃及}。但那位因\Person{若瑟}的功绩而善待他们的国王去世后，他们便被后继的政府压制。因为不仅建造城市的沉重劳役加在他们身上，而且由于他们人数众多引人恐惧，唯恐有朝一日他们会以武力获得独立，于是国王下令，必须将他们新生的男婴溺死。没有人被允许逃避这道残酷的命令。恰在此时，\Person{法郎}的女儿在河中发现了一个婴儿，便收养他作自己的儿子，给他起名叫\Person{梅瑟}。这\Person{梅瑟}长大后，看见一个希伯来人被埃及人殴打；他见状满怀悲痛，便解救了他的弟兄，用石头打死了那埃及人。不久，因害怕自己所做的事受惩罚，他逃往\Place{米德扬}地，寄居在那地区的司祭\Person{耶特洛}家里，娶了他的女儿\Person{漆颇辣}为妻，她为他生了两个儿子，\Person{革尔熊}和\Person{厄里厄则尔}。在这个时代，\Person{约伯}生活着，他仅凭自然律就获得了对天主的认识和一切正义。他极其富有，因而更加显赫，因为他既未在财富完整时被其腐化，也未在失去时被其扭曲。因为，当藉魔鬼的作为，他被剥夺了财产，失去了子女，最后自身又长满了可怕的脓疮时，他并未被击垮，以至于因不堪痛苦而犯任何罪过。最终他获得了天主赞许的赏报，恢复健康后，他得以双倍收回他所失去的一切。\par

% structure: p0028 source=h2 target=subsection reason=relative-heading-rank; base=h2

\subsection{第十四章。}

但希伯来人因奴隶制的重重苦难而备受压迫，遂将哀怨诉诸上天，并怀抱从天主而来的援助希望。那时，当\Person{梅瑟}在牧羊时，突然有一丛荆棘向他显现，焚烧着，但令人惊奇的是，火焰并未伤害它。面对如此奇异的景象，他惊愕不已，便走近荆棘，天主立即向他说话，大意如下：祂是\Person{亚巴郎}、\Person{依撒格}和\Person{雅各伯}的天主，并愿意将受埃及人暴政压制的他们的后裔从苦难中解救出来；因此，他应当去埃及国王面前，并自荐为恢复他们自由的领袖。当他犹豫时，天主以能力坚固了他，并赐予他施行奇迹的恩赐。于是\Person{梅瑟}进入\Place{埃及}，先在同胞面前行了奇迹，并与他的兄弟\Person{亚郎}同去，晋见国王，宣布他是天主所派遣的，如今他以天主的话告诉国王，要释放希伯来人。但国王声称不认识主，拒绝服从给他的命令。当\Person{梅瑟}为证明他所传的命令来自天主，将他的棍杖变为蛇，不久又将所有的水变为血，并使全地布满青蛙时，由于加色丁人也在做类似的事，国王宣称\Person{梅瑟}所行的奇事不过是魔术，而非天主的能力，直到地上布满了飞来的昆虫，加色丁人承认这是神圣的威能所致。于是国王迫于苦难，将\Person{梅瑟}和\Person{亚郎}召来，同意让百姓离去，前提是国中的灾祸被移除。但灾祸停止后，他的内心不受控制，回到先前状态，并不按约定允许以色列人离开。最终，他被降在他自身及国上的十灾所压倒和征服。\par

% structure: p0030 source=h2 target=subsection reason=relative-heading-rank; base=h2

\subsection{第十五章。}

但是，在百姓出埃及的前一日，他们尚不熟悉日期，却因天主的命令而受教，要把那正在经过的那一月算作万月之首；并被告知，当日的祭献要在未来的岁月中庄重而按时地奉献，即在月之第十四日，要宰杀一只无残疾的一岁羔羊作为祭品，并用它的血涂在门框上；它的肉要全部吃掉，但一根骨头也不可折断；他们要在七日内戒绝发酵之物，只吃无酵饼；并且要将这规例传给后代。于是百姓带着财富出发，既有自己的财物，更有从埃及夺来的战利品。他们的人数从最初下到埃及的那七十五位希伯来人，增长到了六十万男子。从亚巴郎最初到达客纳罕地算起，已过去了四百三十年，而从洪水算起，则过去了五百七十五年。当他们匆忙前行时，白天有云柱，夜间有火柱，在他们前面行走。但是，由于红海海湾横亘其间，道路原本经过培肋舍特人的地方，为了避免日后希伯来人因畏惧旷野而有机会通过一条熟悉的陆路返回埃及，他们遵照天主的命令转向，朝红海行进，在那里停下并安营。当有人报告给法老说，希伯来百姓因迷路而来到海边，前有大海阻挡，无路可逃时，法老既恼怒又忿恨，如此众多的人竟能逃脱他的王国与权势，便急忙率领军队出发。当兵器、军旗和排列在广阔平原上的阵线已清晰可见时，希伯来人恐惧万分，仰望上天，梅瑟受天主指示，用杖击打海水，将海分开。这样，百姓便在干地上有了一条道路，海水在两边退去。埃及法老毫不迟疑地追赶前进的以色列人，进入了分开的海中；但海水迅速合拢，他和他所有的军队都被毁灭了。\par

% structure: p0032 source=h2 target=subsection reason=relative-heading-rank; base=h2

\subsection{第16章。}

于是，梅瑟因自己百姓获救、敌人遭此奇迹般毁灭而欢欣鼓舞，向天主唱了一首赞美之歌，全体民众，无论男女，都一同歌唱。但进入旷野，行了三天的路程之后，缺水的困苦临到他们；找到水时，却因水苦而不能用。那时，不耐烦之百姓的倔强首次显露出来，爆发对梅瑟的怨言；梅瑟受天主指示，将一些木头投入水中，木头的力量使水变得甘甜。由此前进，众民在厄林找到了十二股水泉和七十棵棕树，便在那里安营。百姓又因饥荒而抱怨，将辱骂加于梅瑟，并怀念埃及的奴役生活，因为那里有满足他们口腹之欲的丰盛；此时，有一群鹌鹑由天主降下，充满了营地。此外，次日，那些从营中出来的人看见地上覆盖着一种类似荚果的东西，其外观像雪白的胡荽籽，正如我们在冬季常见地上铺满白霜一般。于是，百姓通过梅瑟得知，这食粮是天主恩赐给他们的；各人要用预备好的器皿按人数收取足够一日之用；但第六日要收取双倍，因为安息日不可收取。然而，百姓向来不惯于服从，按人性不禁欲望，不仅为当日，也为次日储备。但所储藏之物生出了蛆虫，发出恶臭，然而第六日为了安息日所储藏之物却丝毫未腐坏。希伯来人食用这食物达四十年；其味道很像蜜；其名流传下来，称为\OriginalTerm{玛纳}{manna}。此外，作为天主恩赐的永久见证，据说梅瑟将满满一曷默尔放在金器中存留。\par

% structure: p0034 source=h2 target=subsection reason=relative-heading-rank; base=h2

\subsection{第17章。}

百姓从那里起行，再次遭遇缺水之苦，几乎难以自控，要杀死他们的领袖。于是梅瑟奉天主之命，用杖击打那称为曷勒布地方的磐石，便涌出丰沛的水源。但他们到了辣非丁时，阿玛肋克人发动攻击，杀害了许多百姓。梅瑟率众出战，派若苏厄率领军队；他自己则与亚郎和胡尔一同，只作战事的旁观者，同时为了向主祈祷，他上了一座山顶。两军交战，胜负未分时，因着梅瑟的祈祷，若苏厄击杀敌人直到日暮。同时，梅瑟的岳父耶特洛偕同他的女儿漆颇辣（她曾嫁给梅瑟，丈夫去埃及时却留在家中）和孩子们，得知梅瑟所行的事，便来到他那里。依他的建议，梅瑟将百姓分成各种等级，并设立千夫长、百夫长、五十夫长管理他们，从而为后世提供了一种纪律与秩序的规范。耶特洛随后返回本国，以色列子民则来到西乃山。在那里，主告诫梅瑟，百姓应当圣洁，因为他们将聆听天主的言语；这事被谨慎遵行。当天主降在山巅时，空中响彻嘹亮的号声，浓云翻滚，闪电频频。梅瑟和亚郎上到山顶，陪伴在上主旁边，而百姓则站在山脚下。于是颁赐了法律，内容丰富，充满天主的话语，且屡次重复；若有人想详细了解，须查阅原典，因为我们这里只是扼要提及。\Quote{你们中间不可有别的神，}天主说，\Quote{但你们要独尊崇我；你们不可为自己制造任何偶像；你们不可妄呼你们天主的名；安息日不可工作；孝敬你们的父亲和母亲；不可杀人；不可奸淫；不可偷盗；不可作假见证陷害你们的邻人；不可贪恋邻人的一切。}\par

% structure: p0036 source=h2 target=subsection reason=relative-heading-rank; base=h2

\subsection{第十八章}

天主说了这些话，号声响起，电光闪烁，浓烟笼罩了山岭，百姓因恐惧而战栗；他们恳求梅瑟，让天主只与他单独说话，而他则将他所听到的转告百姓。天主对梅瑟的诫命如下：用银子买来的希伯来奴仆要服事六年，第七年可自由；但他的耳朵要被穿孔，若他甘愿为奴。凡杀人的，应处死刑；误杀人的，须按规程流放。凡殴打父亲或母亲，或出言辱骂他们的，应处死刑。若有人出卖偷来的希伯来人，应处死刑。若有人打自己的男仆或女仆，致其当场死亡，应受审讯。若有人致使妇女流产，应处死刑。若有人打掉仆人的眼或牙，该仆人应依法获得自由。若公牛顶死人，该牛应被砸死；若主人知道牛的凶性而未加防范，他也应被砸死，或按控告者所要求的价金赎命。若公牛顶死仆人，应向其主人支付三十双德拉克玛银子。若有人挖坑不加盖，致使牲畜掉入，应赔偿牲畜的价值给其主人。若一牛顶死另一人的牛，应卖掉活牛，二人平分价钱；死牛也平分。但若主人知道牛的凶性而未加防范，应赔出公牛。若有人偷牛犊，应赔五头；若偷羊，应赔四倍；若赃物尚在他手中，应加倍偿还。夜间偷盗者可被杀死，白日则不可。若牲畜吃了别人的庄稼，牲畜的主人应赔偿损失。若寄存之物丢失，代管者应发誓未行欺诈。被抓住的贼应加倍赔偿。托管的牲畜若被野兽吞噬，无须赔偿。若有人玷污尚未许配的童贞女，应给女子聘礼，娶她为妻；但若女子父亲拒绝嫁女，则玷污者仍须给她聘礼。若有人与牲畜交合，应处死刑。献祭给偶像的，应灭亡。不可欺压寡妇和孤儿；不可苛刻对待贫困的债户，不可索取高利；不可拿穷人的衣服作抵押。不可毁谤百姓的首领。所有头胎的当献给天主。被野兽撕裂的肉不可吃。不可订立虚假见证或任何邪恶的协议。不可对你仇敌走失的牲畜视而不见，应牵回。若你看见仇敌的牲畜被重担压倒，应扶起它。不可杀害无辜和义人。不可因贿赂而称恶人为义。不可收受礼物。对外方人要友善。六日应工作，安息日应休息。第七年的收成不可收割，应为穷苦和需要的人留下。\par

% structure: p0038 source=h2 target=subsection reason=relative-heading-rank; base=h2

\subsection{第十九章}

\Person{梅瑟}将这些天主的话报告给百姓，并在山脚下用十二块石头立了一座祭台。然后他再次上山，主已在那里，他带了\Person{亚郎}、\Person{拿答}和七十位长老同去。但他们不能瞻仰主的面；然而，他们看见了天主所站之处，据说他的形像奇妙，他的荣光灿烂。\Person{梅瑟}受天主召唤，进入围绕天主的云层深处，据记载他在那里住了四十昼夜。在此期间，他受天主的话教导，关于建造会幕和约柜，以及祭献的礼仪——这些事情，因显然叙述冗长，我认为不宜插入这样简略的著作中。但因\Person{梅瑟}离开日久——他在主面前度过了四十天——百姓绝望于他的归来，便强迫\Person{亚郎}制造神像。于是，从熔化的金属中造出了一个牛头。百姓忘记了天主，向它献祭，纵情于吃喝。天主看到这些，义愤填膺，本要消灭这邪恶的百姓，若不是\Person{梅瑟}恳求他不要这样做。但\Person{梅瑟}下山时，手里拿着两块由天主亲手书写的石版，看见百姓沉溺于享乐和亵渎，便摔碎石版，认为这民族不配领受主的法律。然后他召集被多方辱骂的肋未人到自己身边，命令他们拔剑击杀百姓。在这次袭击中，据说有二万三千人被杀。随后\Person{梅瑟}在营外支搭会幕；每当他进入会幕，就见云柱停在门口；天主与\Person{梅瑟}面对面说话。但当\Person{梅瑟}恳求看见主独特的威严时，他被告知天主的形像非肉眼所能见；但允许看见他的后背；\Person{梅瑟}先前摔碎的石版被重新制作。据记载，\Person{梅瑟}在这次与天主的交谈中，与主同住了四十天。此外，当他从山上下来，带着石版时，他的脸面发光，如此明亮，以致百姓不能看他。因此，规定当他向百姓传达天主的命令时，他用帕子蒙上脸，就这样以天主的话向百姓说话。在此段历史中，记述了会幕及其内部的建造。完工后，有云从上降下，笼罩会幕，甚至阻止\Person{梅瑟}本人进入。以上是\WorkTitle{创世记}和\WorkTitle{出谷纪}两书所包含的主要内容。\par

% structure: p0040 source=h2 target=subsection reason=relative-heading-rank; base=h2

\subsection{第二十章}

接下来是\WorkTitle{肋未纪}，其中阐述了关于祭献的规条；又对先前所赐的法律增添了诫命；几乎全部充满与司铎有关的训示。若有人想了解这些，可以从那本书中获得更详细的信息。因为我们在所承担工作的范围内，只涉及历史。肋未支派被分别出来担任司铎后，其余支派被统计，人数共计六十万三千五百人。当百姓以玛纳为食物时，如上所述，尽管天主如此多、如此大的恩惠，他们却一如既往地忘恩负义，渴望在埃及所习惯的粗劣食物。于是主带来大量的鹌鹑进入营中；当他们急切地撕食这些鹌鹑时，嘴唇一碰到肉，便丧命了。那一日，营中果然大遭毁灭，据说有二万三千人死亡。百姓就这样因他们所贪求的食物而受惩罚。然后队伍前行，来到\Place{帕兰}；\Person{梅瑟}受主指示：主所应许赐给他们的土地已经近了。于是派遣探子进入那地，他们报告说那是一块物产丰饶的土地，但当地的民族强大，城镇有巨大的城墙环绕。这消息传到百姓耳中，众人心中充满恐惧；他们竟邪恶到如此地步，藐视\Person{梅瑟}的权威，准备为自己立一个首领，在他的带领下返回埃及。那时，\Person{若苏厄}和\Person{加肋布}——两人本是探子之一——撕裂衣服，流泪恳求百姓不要相信那些报告恐怖消息的探子；因为他们自己也同去，并未发现那地有什么可怕之处；他们应当相信天主的应许，这些敌人将成为他们的战利品，而非他们的毁灭。但这顽固的民族反抗一切善意的劝告，冲向他们要杀害他们。主因这些事发怒，使一部分百姓被敌人杀死，而探子们因在百姓中引起恐惧而被杀。\par

% structure: p0042 source=h2 target=subsection reason=relative-heading-rank; base=h2

\subsection{第二十一章}

随后发生了叛乱，那些以达堂和阿彼兰为首的叛党企图反抗梅瑟和亚郎；但大地裂开口，将他们活活吞下。不久之后，全体百姓又起来叛乱，反对梅瑟和亚郎，他们冲进了只有司铎才能进入的会幕。那时，死亡真的成堆地收割他们；若不是主因梅瑟的祈祷而息怒，转避了灾祸，所有人都会在顷刻间灭亡。尽管如此，被杀的人数仍达到了七十一万四千。又过了不久，如同先前屡次发生的那样，百姓因缺水而叛乱。于是，梅瑟奉天主之命用棍杖击打岩石，用他早已熟悉的那种试探（因为他先前已做过同样的事），两次击打岩石，水便从岩石中流了出来。然而，在这件事上，据说梅瑟受到了天主的责备：由于缺乏信德，他只得通过反复击打才引出水来；事实上，因这过犯，他未能进入应许之地，我将在后面说明。梅瑟离开那地方，准备带领队伍沿着厄东的边界前进，便派使者去见厄东王，请求允许借道；他认为出于血缘关系，应当避免战争；因为那民族是厄撒乌的后裔。但厄东王藐视了请求者，拒绝了借道的要求，并准备以武力相争。于是，梅瑟率军转向曷尔山，避开那禁行之路，以免在同族之间挑起战争，并在这条路上消灭了客纳罕人的王。他也击败了阿摩黎人的王息红，占据了他们所有的城镇；他还征服了巴商和巴拉克。他在约但河对岸扎营，离耶里哥不远。随后，与米德杨人发生了战斗，他们被击败并臣服了。梅瑟在旷野治理百姓四十年后去世。希伯来人之所以在旷野停留如此之久，据说是为了让所有不信天主言语的人都死去。因为除了若苏厄和加肋布，离开埃及时年满二十岁以上的人，没有一人越过约但河。梅瑟本人只是看见了应许之地，却未能到达，这归咎于他的罪，因为在那次奉命击石取水时，尽管有了如此多奇迹能力的证据，他仍心存疑虑。他享年一百二十岁去世。关于他的埋葬地点，无人知晓。\par

% structure: p0044 source=h2 target=subsection reason=relative-heading-rank; base=h2

\subsection{第二十二章}

梅瑟去世后，首权传到了农的儿子若苏厄手中，因为梅瑟立他为自己的继承人，他是一位在显示的种种美德上与梅瑟极为相似的人。在他统治之初，他派使者走遍军营，吩咐百姓准备粮谷，并宣布他们应在第三天起程。但约但河，一条水流湍急的大河，阻碍了他们的渡河，因为他们没有足够的船只可用于渡河，而且当时河水正泛滥，无法涉水而过。于是，他命令司铎们抬着约柜前行，并逆水而立。这事一办，据说约但河便分开了，这样军队就在干地上渡过了河。那地方有一座城，名叫耶里哥，城墙坚固，不易攻取，无论是强攻还是围困都难以奏效。但若苏厄信赖天主，没有用武器或武力攻城；他只是命令抬着天主的约柜绕城而行，司铎们走在约柜前，吹着号角。当约柜绕城七周后，城墙和城楼都倒塌了；于是城被洗劫并焚烧。然后，据说若苏厄向主祈祷，并诅咒任何试图重建这座被天主相助而摧毁之城的人。接着，军队被带去攻打革特城，并设伏兵于城后；若苏厄假装畏惧，在敌人面前逃跑。城里的人见了，便打开城门，开始追击溃退的敌人。于是，伏兵攻占了城，城中居民全被杀尽，无一逃脱；那王也被擒获，并被处以极刑。\par

% structure: p0046 source=h2 target=subsection reason=relative-heading-rank; base=h2

\subsection{第二十三章}

此消息为邻近诸王所知后，他们便结成战争同盟，企图以武力征服希伯来人。但基贝红人，一个强盛民族，拥有富庶的城市，却自愿降服于希伯来人，承诺遵行一切命令，因而受到保护，并被指派搬运柴水。然而，他们的归顺激起了附近城王的愤怒。于是，那些王调集军队，围攻他们名为\Place{加巴敖特}的城镇。基贝红人陷入困境，遂派使者向\Person{若苏厄}求救，请求他援助被围之城。\Person{若苏厄}便以强行军突然出现在敌军面前，歼灭数千之众。胜利者在白昼将尽时，似乎黑夜将为败军提供庇护；但希伯来统帅凭借信德之力，阻止了黑夜，白昼得以延续，敌人无从逃脱。被擒获的五王悉数处死。在同样的攻击下，邻近城市也落入\Person{若苏厄}手中，其诸王被剪灭。然而，由于我力求简洁，无意逐一详述这些事件，只留心注意到：有二十九个王国被置于希伯来人轭下，其土地分配给了十一个支派，各按人口分给。至于被选为司祭的肋未人，则未分给土地，以便他们能更自由地事奉天主。我不愿默默略过这一范例，反而要热切指出，它值得教会牧者诵读。在我看来，他们不仅忘记这条诫命，甚至全然无知——因为在这个时代，一种占有的贪欲像不治之症一样占据了他们的心灵。他们渴望产业，经营田产，倚靠金银，买进卖出，用尽一切手段牟利。即便他们中有人似乎有更高的生活目标，既不拥有也不经商，但仍然（更可耻的是）无所事事，期待馈赠，并因唯利是图的性情败坏了生活的全部荣耀，同时摆出圣善的外表，仿佛连圣善也可以成为获利之源。但我对当代品性表达了憎恶与厌烦，这已超出了原定范围；我急忙回到正题。被征服的土地，如我所说，分给各支派后，希伯来人享受了深切的和平；邻国因战事惊恐，不敢对那些功勋卓著者发动敌对行动。同一时期，\Person{若苏厄}在第一百一十岁上去世。关于他统治的时间长短，我不下明确判断；但普遍观点认为他领导希伯来事务二十七年。若果真如此，那么从世界之初直到他去世，共经过了三千八百八十四年。\par

% structure: p0048 source=h2 target=subsection reason=relative-heading-rank; base=h2

\subsection{第二十四章}

若苏厄死后，百姓没有领袖行事。但因与客纳罕人作战的需要，犹大被任命为作战统帅。在他的引导下，诸事顺利：内外均极其安宁；百姓统治着那些被征服或按条件投降的民族。随后，如同太平时期几乎总会发生的那样，他们渐渐忘记道德和纪律，开始与征服者通婚，并逐渐采纳外邦习俗，甚至以亵圣的方式向偶像献祭：与外邦人结盟如此有害。天主早已预见这些事，曾以一条有益的命令吩咐希伯来人将征服的民族完全消灭。然而，百姓因贪求权力，宁愿（自取灭亡地）统治被征服者。因此，当他们离弃天主、崇拜偶像时，便失去了天主的助佑，被美索不达米亚王击败制服，受了八年俘虏的惩罚，直到以敖特尼耳为领袖，才重获自由，并享有五十年独立。随后，久安带来的恶果再次腐蚀他们，他们又开始向偶像献祭。他们如此犯罪，报应迅速降临。被摩阿布王厄革隆征服后，他们服事了他十八年，直至因天主的推动，厄胡得以计谋杀死敌王，并迅速召集一支军队，以武力恢复他们自由。此人治理希伯来人平安四十年。沙默加尔继他之后，与培肋舍特人交战，取得决定性胜利。然而，再度拜偶像的希伯来人又被客纳罕王雅宾制服，受其暴虐统治二十年，直至一位妇女德波辣使他们恢复原状。他们已对领袖失去信心，以至现在竟靠一位妇女保护。但值得注意的是，这种拯救方式预先安排，作为教会的预像，藉其帮助，人得以逃脱魔鬼的奴役。希伯来人在此领袖或民长治下四十年。又因罪恶被交于米德杨人手中，受严酷统治；他们受奴隶之苦，恳求天主的帮助。果然，当顺遂时，他们总忘记天主的恩惠，祈求偶像；但在逆境中，他们向天主呼号。因此，每当我思量那蒙受天主如此多恩惠的民族，当他们犯罪时受如此多灾祸惩罚，既经历天主的仁慈也经历天主的严厉，却丝毫未见改善，并且虽然他们每次犯罪都获得宽恕，但在宽恕后又不断犯罪，那么基督来临时不被他们接纳，就不足为奇了，因为从一开始，他们就多次背叛主。事实上，更奇妙的是，天主的仁慈从未在他们犯罪时抛弃他们，只要他们呼号他的名。\par

% structure: p0050 source=h2 target=subsection reason=relative-heading-rank; base=h2

\subsection{第二十五章}

\Emph{因此}，当米德杨人如上所述统治他们时，他们转向主，恳求他惯常的柔慈，并得蒙垂允。当时希伯来人中有一人名基德红，是位正义之人，为天主所喜爱和悦纳。天使在他从田里回家时站在他旁边，对他说：\Quote{主与你同在，英勇的壮士。}但他以谦卑的声音抱怨主不同他同在，因为俘虏之苦重压他的百姓，他含泪记起主所行的奇事，就是领他们出埃及地。天使便说：\Quote{去吧，就照你所说的这精神，拯救百姓脱离被掳。}但他声明以他微弱的力量不能承担如此重任，因他乃极卑微之人。天使却坚持敦促他，不要怀疑主所说的事必能成就。于是，他献了祭，推翻了米德杨人为巴耳像所筑的祭坛，便往自己百姓那里去，在敌营附近安营。阿玛肋克人也加入米德杨人一边，而基德红只聚集了三万二千人的军队。战前，天主对他说，这支人数比他希望他带去交战的更多；如果他使用这么多人，希伯来人将照他们惯常的邪恶，将战果不归于天主，而归于自己的勇力；因此他应给那些愿意离开的人提供机会。这话告知百姓后，有两万二千人离营。但剩下的一万人中，基德红遵照天主的指示，只保留了三百人，其余都从战场遣散。于是，在夜间三更时分，他进入敌营，命令所有人吹号角，使敌人大起恐慌，无人敢抵抗，都四散狼狈逃窜。然而，希伯来人从各方截击，击杀逃敌。基德红追赶诸王到约但河对岸，擒获后处死他们。在那场战役中，据说敌人死了十二万，俘虏了一万五千。于是，全民一致提议立基德红为王。但他拒绝此议，宁可与同胞平等相处，也不愿做他们的统治者。因此，他们从压了百姓七年的被掳中解脱后，享受了四十年的和平。\par

% structure: p0052 source=h2 target=subsection reason=relative-heading-rank; base=h2

\subsection{第二十六章}

\Person{基德红}死后，他的儿子\Person{阿彼默肋客}（其母为妾侍）在众多恶人的协助下，尤其依靠\Place{舍根}首领们的帮助，杀死了他的兄弟们，夺取了王国。他因内乱而受困扰，同时以战争压迫他的人民，试图攻占某座塔楼，城中居民失城后逃入其中。然而，当他不够谨慎地靠近该地时，被一个女人投掷的石头击中而死，他统治了三年。继他之后是\Person{托拉}，执政二十二年。他之后是\Person{雅依尔}；\Person{雅依尔}也统治了二十二年后，百姓离弃了天主，献身于偶像。为此，以色列子民被\Place{培肋舍特}人和\Place{阿孟}人征服，在他们的统治下十八年。在此期间结束时，他们开始呼求天主；但天主的答复是，他们不如去祈求他们偶像的帮助，因为他不再向如此忘恩负义的人施恩。他们却流泪承认自己的过错，恳求宽恕；同时抛弃偶像，热切呼求天主，终于获得了天主的怜悯，尽管起初被拒绝。因此，在\Person{依弗大}率领下，他们大量聚集，意图以武力恢复自由，并先派遣使者到\Place{阿孟}王那里，恳求他满足于自己的领土，不要与他们交战。但他非但不拒绝战斗，反而立即整军备战。然后，据说\Person{依弗大}在开战信号发出前许下誓愿：如果他获得胜利，回家时首先迎接他的人，应作为祭献献给天主。于是，敌人被打败后，\Person{依弗大}回家时，他的女儿出来迎接他，欢欣鼓舞地打着鼓跳舞，迎接父亲凯旋。\Person{依弗大}满怀忧伤，痛苦地撕裂衣服，向女儿说明了他誓愿的严厉约束。但她以超乎女子的勇气，并未拒绝死亡；只恳求宽限两个月，以便在死前得以与同龄的朋友见面。此事完成后，她甘愿回到父亲那里，向天主履行了誓愿。\Person{依弗大}执政六年。继他之后是\Person{厄色本}，平安统治七年后去世。其后，\Place{则步隆}人\Person{厄隆}统治了十年，\Person{阿贝冬}统治了八年；但因他们的统治太平无事，没有做出什么值得历史记载的事迹。\par

% structure: p0054 source=h2 target=subsection reason=relative-heading-rank; base=h2

\subsection{第二十七章。}

以色列子民再次转向偶像；因此失去天主的保护，被\Place{培肋舍特}人征服，因背信而受罚被囚四十年。那时，\Person{三松}诞生了。他的母亲久不生育，曾看见天使显现，被告诫要戒绝清酒、烈酒及一切不洁之物；因为她要生一个儿子，将成为以色列子民恢复自由者，并向他们的敌人复仇。他蓄着长发，据说拥有惊人的力量，竟能手撕一头拦路的狮子。他娶了一位\Place{培肋舍特}女子为妻；当她在丈夫不在时嫁给另一个人时，\Person{三松}因妻子被夺而愤怒，遂向她的民族施行毁灭。他依靠天主和自己的勇力，公开给那些至今的胜利者带来灾难。他捉了三百只狐狸，将燃烧的火把绑在它们尾巴上，放入敌人的田地里。恰逢收割时节，庄稼成熟，火势容易蔓延，葡萄园和橄榄树都被烧成灰烬。这样，他通过给\Place{培肋舍特}人造成重大损失，报复了夺妻之辱。\Place{培肋舍特}人为这场灾难所激怒，用火烧死了那位导致如此大祸的女子，连同她的房屋和父亲。但\Person{三松}认为自己尚未充分复仇，仍以各种恶计不断骚扰异族。于是犹太人被迫将他作为俘虏交给\Place{培肋舍特}人；但被交出后，他挣断绳索，顺手拾起一块驴腮骨作为武器，击杀了一千敌人。当日间酷热难当，他开始口渴难耐时，他呼求天主，水便从他手中的骨头里流了出来。\par

% structure: p0056 source=h2 target=subsection reason=relative-heading-rank; base=h2

\subsection{第二十八章。}

那时，\Person{三松}统治希伯来人，培肋舍特人因一人的英勇而被征服。于是他们设谋取他的性命，不敢公开攻击他，为此贿赂他的妻子（他在所述之事后所娶的），让她出卖丈夫的力量所在。她用女性的谄媚攻击他；他多次欺骗她，拖延她的目的，最后她诱他说出他的力量在于头发。不久，趁他睡着时，她偷偷剪去他的头发，就这样将他交在培肋舍特人手中；尽管他以前常被交在他们手中，他们却无法拘禁他。他们剜了他的眼睛，用镣铐捆住他，把他关进监狱。但过了一段时间，他被剪去的头发开始重新生长，力量也随之恢复。\Person{三松}意识到自己恢复了力量，只等待时机进行正义的报复。培肋舍特人有个习俗，在节日时把\Person{三松}带出来当众戏弄，嘲弄他们这位显赫的俘虏。于是在某一天，当他们为自己的偶像举行盛宴时，他们命令把\Person{三松}展示出来。那庙宇——所有百姓和培肋舍特众首领都在其中宴饮——靠两根极大的柱子支撑；\Person{三松}被带出来时，被置于两柱之间。他先呼求了主，然后抓住时机，推倒了柱子。全体群众都被压在建筑物的废墟下，\Person{三松}自己也与敌人同死，在统治希伯来人二十年之后，总算向他们报了仇。\Person{西米加尔}接续他，关于此人圣经只记了这点事实。因为连他的统治何时结束我也没找到记载，而且我看到百姓一度没有领袖。后来对\Person{本雅明}支派爆发内战，\Person{犹大}被选为临时战争领袖。但多数记录这些时期的作者指出他的统治仅一年。因此，许多人完全略过他，把司祭\Person{厄里}直接排在\Person{三松}之后。我们将此点存疑，视为未确定之事。\par

% structure: p0058 source=h2 target=subsection reason=relative-heading-rank; base=h2

\subsection{第二十九章}

大约此时，正如我们所说，内战爆发了。骚乱的起因如下：一个肋未人带着他的妾旅行，因夜幕降临，便在\Place{加巴}镇投宿，该镇住着本雅明人。一位老人善意地接待了他们留宿；但镇上的年轻人包围了客人，想要对他行不轨之事。老人严厉斥责他们，好不容易劝阻了他们的意图，他们最后同意用他的妾代替他本人供他们淫乐；这样他们放过了陌生人，却凌辱了那妇人一整夜，直到次日才将她归还。但她（不知是因他们卑劣行为所致伤害，还是出于羞耻，我不敢断言）在再次见到丈夫时便死了。肋未人为了证明这可怕的罪行，将她的尸体切成十二块，分给十二支派，以便更容易激起所有人的义愤。众人得知此事后，其他十一支派便结成军事联盟攻打本雅明。在这次战争中，如前所述，\Person{犹大}担任统帅。但头两次战斗他们都不顺利。最终在第三次，本雅明人被击败，全部遭杀戮；如此，少数人的罪行由多人的毁灭而受惩罚。这些事也记载在\BibleBook{民长}{纪}中；接着是\BibleBook{列王}{纪}。但对我这个遵循年代顺序和日期次序的人来说，历史似乎并不严格符合年代学。因为在\Person{三松}作民长之后，有\Person{舍米加尔}，稍后历史证实百姓没有民长，而\BibleBook{列王}{纪}中记载司祭\Person{厄里}也曾是民长，但圣经没有说明\Person{厄里}与\Person{三松}之间有多少年。我看到这两者之间有一段时期被模糊了。但从\Person{若苏厄}去世之日直到\Person{三松}去世之时，共算四百一十八年，从创世之初算起则为四千三百零三年。然而，我知道别人与我们的计算有出入；但同时我意识到，我并非没有经过一些努力才把逐年事件顺序（这是迄今模糊不清的）陈述出来，直到我遇到这段时期，我承认对此存疑。现在我继续叙述余下之事。\par

% structure: p0060 source=h2 target=subsection reason=relative-heading-rank; base=h2

\subsection{第三十章}

当时，希伯来人正如我上文所述，各自按自己的意愿生活，没有任何民长或将领。\Person{厄里}作司铎；在他日子，\Person{撒慕尔}诞生了。他父亲名叫\Person{厄耳卡纳}，母亲名叫\Person{亚纳}。她长期不孕，据说当她向天主求子时，曾许愿说，若生男孩，就将他奉献给天主。于是，她生了一个男孩，便把他交给司铎\Person{厄里}。后来，孩子长大，天主对他说话。天主因\Person{厄里}儿子的行为宣布对司铎\Person{厄里}发怒，因为他的儿子们把他父亲的司铎职当作牟利的手段，向前来祭献的人索取礼物；虽然据说他们的父亲屡次责备他们，但他的责备过于温和，未能达到管束的目的。那时，培肋舍特人侵入\Place{犹大}，以色列人迎战。但希伯来人战败，准备再次交战；他们把主的约柜带到战场，司铎的儿子们也同约柜一起出战，因为\Person{厄里}本人年事已高、双目失明，不能担任那职务。但当约柜出现在敌人眼前时，他们好像被天主临在的威严吓倒，准备逃跑。但随后他们重振勇气，改变主意（并非没有天主的推动），全力冲入战斗。希伯来人被击败；约柜被夺走；司铎的儿子们阵亡。\Person{厄里}听到灾祸的消息，悲痛过度，断了气。他担任司铎共二十年。\par

% structure: p0062 source=h2 target=subsection reason=relative-heading-rank; base=h2

\subsection{第三十一章}

在这场胜仗中获胜的培肋舍特人，将落入他们手中的天主的约柜，带到了\Place{阿左托}城的\Person{达贡}庙。但那尊献给邪神的偶像，在约柜被抬进去时倒了下来；他们再把它立回原处，次日夜晚，偶像被摔得粉碎。随后老鼠在全地滋生，它们的有毒咬伤导致数千人死亡。\Place{阿左托}人受此痛苦所迫，为逃避灾祸，将约柜迁至\Place{加特}。但那里的人也遭受同样祸患，便将约柜转送到\Place{阿市刻隆}。然而，该地的居民召集了全国的首领，计划将约柜送回希伯来人那里。于是，依照首领、占卜者和司铎的意见，他们把约柜放在一辆车上，连同许多礼物一起送回。当时发生了一件奇异的事：他们让母牛套上车，却把小牛留在家中，这些牲畜竟无人引领，径直朝\Place{犹大}走去，丝毫没有因挂念留在身后的小牛而想返回。跟随约柜进入希伯来人境内的培肋舍特首领，被这事的奇妙所震撼，便举行了一项宗教仪式。犹太人一看见约柜被送回，都争相从\Place{贝特舍默士}城欢天喜地地冲出去迎接，匆忙、欢跃、感谢天主。随后，肋未人（他们的职务就是如此）向天主献祭，并献上那两头拉了约柜的母牛。但约柜不能保存在我上文提到的城，于是天主安排全城遭受重病。约柜就被运到\Place{克黎雅特耶阿陵}城，并在那里停留了二十年。\par

% structure: p0064 source=h2 target=subsection reason=relative-heading-rank; base=h2

\subsection{第三十二章}

此时，司铎\Person{撒慕尔}治理希伯来人。由于战争止息，民众安居乐业。但这份安宁却被培肋舍特人的入侵打破，各阶层的人都因意识到自己的罪孽而惊恐万状。\Person{撒慕尔}先行了祭献，并信赖天主，率领他的人出战，敌人一交锋就溃败，希伯来人取得了胜利。然而，当敌人的恐惧消除，国泰民安之后，民众转变观念，每况愈下——暴民总是厌倦已有的，向往未曾经历的——他们竟渴望君王的称号，这个称号几乎被所有自由民族深恶痛绝。是的，这真是一个极其疯狂的愚行例证，他们现在宁愿用自由换取奴役。于是，他们成群结队来到\Person{撒慕尔}面前，因为他自己年事已高，要求他为他们立一位君王。\Person{撒慕尔}则用一番有益的讲话，试图安抚民众的非分之想；他陈明了君王的暴政和傲慢统治，却赞扬自由、谴责奴役；最后，他以天主的义怒警告他们，如果他们竟愚昧到这样的地步，既然有天主作他们的君王，却还要从人类中要求一个君王。他说了这些以及类似的话，但都白费口舌，见民众固执己见，他便请示天主。天主被那疯狂民族的愚行所触动，回答说：他们要求不利于自己的事，什么也不可拒绝他们。\par

% structure: p0066 source=h2 target=subsection reason=relative-heading-rank; base=h2

\subsection{第三十三章}

因此，撒乌耳先由撒慕尔以司铎之油傅油后，被立为王。他属本雅明支派，父亲名叫克士。他心性谦逊，容貌异常俊美，其人之尊严与王权相称。但在其统治初期，部分百姓背叛他，拒绝承认其权威，并投靠了阿孟人。然而，撒乌耳有力地报复了这些人；敌人被征服，希伯来人获得宽赦。随后，撒乌耳据说由撒慕尔第二次傅油。接着，因培肋舍特人入侵，爆发了一场血战；撒乌耳指定基耳加耳为军队集结之地。他们在那里等候撒慕尔七日，以便他向天主献祭，但因他的迟延，百姓逐渐散去，而王竟以非法僭越，献上了全燔祭，如此僭取了司铎的职责。为此他受到撒慕尔的严厉斥责，并承认自己的罪，但忏悔已为时过晚。因为王的罪，恐惧弥漫全军。敌人营地相距不远，显明了实际的危险，无人有勇气出战；多数人逃往沼泽。因为除了那些因王的罪感到天主疏远的人缺乏勇气外，军队严重缺少铁制武器，以至于除了撒乌耳和他的儿子约纳堂外，无人拥有刀或枪。原来培肋舍特人在之前的战争中得胜，已剥夺了希伯来人使用武器的权利，无人能锻造任何兵器，甚至制作农具也不能。在此情形下，约纳堂以大胆的计谋，仅带执械者一人，进入敌营，杀了约二十人，使全军惊恐。然后，借着天主的安排，敌人转而逃跑，他们既不服从命令也不保持队列，只将安全希望寄托于逃跑。撒乌耳察觉后，迅速率兵追击逃敌，获得胜利。据说王在那天宣布命令，在敌人被消灭前，任何人不得进食。但约纳堂不知此禁令，找到蜂巢，用杖头蘸蜜吃了。此事因随后天主的愤怒而被王知晓，他下令处死儿子。但百姓相助，使他免于死亡。那时，撒慕尔受天主指示，去见王，以天主的话告知他要攻打阿玛肋克人，他们曾在希伯来人出埃及之时阻碍他们；并附加禁令，不可贪图被征服者的任何战利品。于是，军队开入敌境，王被俘，该民族被征服。但撒乌耳无法抗拒大量战利品的诱惑，忘记了天主的命令，下令保存并聚集掠物。\par

% structure: p0068 source=h2 target=subsection reason=relative-heading-rank; base=h2

\subsection{第34章}

天主对所行之事不悦，对撒慕尔说，祂后悔立撒乌耳为王。司铎将所闻报告给王。不久，受天主指示，他给达味傅了王油，那时达味尚是幼童，在父亲照管下牧羊，且常弹竖琴。因此，后来他被撒乌耳带走，列入王之仆役中。此时培肋舍特人与希伯来人激战正酣，两军对垒，培肋舍特人中有一人名叫哥肋雅，身材异常高大有力，沿己方阵营用最激烈的言词辱骂敌人，并向任何人挑战单挑。于是王应许重赏并将女儿嫁给能夺回那夸口者战利品的人；但如此众多的人中无人敢尝试。在此情形下，达味虽尚为少年，自请出战，并拒绝那压垮他幼年之身的武器，仅以一根杖和五块石头迎战。第一击，他用投石器射出一石，击倒了培肋舍特人；随后割下被征服敌人的头，夺其战利品，后将他的刀存放在圣殿中。与此同时，所有培肋舍特人逃窜，将胜利交给希伯来人。但达味从战场归来时所受的极大宠爱引起了王的嫉妒。然而，恐怕杀害如此受众人爱戴的人会招致对自己的憎恨并带来灾祸，他便决定在尊荣的外表下使达味陷入危险。首先，他任命达味为军官，负责战事；其次，虽然他曾应许将女儿嫁给他，却失信，将她给了别人。不久，王的一个小女儿名叫米加耳，狂热地爱上了达味。于是，撒乌耳向达味提出娶她的条件：若他带回一百个敌人的包皮，王的女儿就嫁给他；他指望那少年冒险如此大的危险时可能丧命。但结果完全出乎他的想象，因为达味按所提的条件，迅速带回了一百个培肋舍特人的包皮；这样他便娶了王之女。\par

% structure: p0070 source=h2 target=subsection reason=relative-heading-rank; base=h2

\subsection{第35章}

君王对他的仇恨与日俱增，这是由于嫉妒的影响，因为恶人总是迫害善人。因此，他命令他的臣仆和他的儿子\Person{约纳堂}设下圈套谋害他的性命。但\Person{约纳堂}从一开始就对\Person{达味}怀有极大的敬重和感情；因此，君王受到儿子的责问，便撤销了他所下的残酷命令。但恶人不会长久为善。因为当\Person{撒乌耳}被错谬之灵折磨时，\Person{达味}站在他旁边，用竖琴安抚他痛苦中的心神，\Person{撒乌耳}试图用矛刺穿他，若不是他迅速躲开了那致命的一击，他早已得逞。从那时起，君王不再秘密，而是公开地设法杀害他；\Person{达味}也不再信任他的权力。他逃走了，起初投奔\Person{撒慕尔}，然后投奔\Person{阿希默肋客}，最后逃到摩阿布王那里。不久，在先知\Person{加得}的指示下，他回到犹大境内，在那里生命处于危险之中。那时，\Person{撒乌耳}杀了司铎\Person{阿希默肋客}，因为他接待了\Person{达味}；当君王的臣仆中没有人敢下手加害司铎时，叙利亚人\Person{多厄格}执行了这项残酷的任务。此后，\Person{达味}逃往旷野。\Person{撒乌耳}也跟踪他到那里，但他想杀害他的一切努力都是枉然，因为天主保护了他。旷野中有一个洞穴，入口之后有极大的空间。\Person{达味}藏身于洞穴深处。\Person{撒乌耳}不知道他在那里，进去为了休息身体，在那里因困倦而睡着了。当\Person{达味}看到这情景时，虽然众人都劝他利用这个机会，但他没有杀害君王，只取走了他的外套。随后他出去，从后面一个安全的位置对君王说话，历数他为君王所做的事，他如何多次为王国冒生命危险，以及最后这次，当天主将君王交在他手中时，他没有设法杀他。\Person{撒乌耳}听了这些话，承认了自己的过错，恳求宽恕，流下眼泪，赞扬\Person{达味}的虔诚，责备自己的邪恶，并称\Person{达味}为王和儿子。他变得与先前凶猛的性情判若两人，以至于现在没有人会想到他会再对他的女婿下手。但\Person{达味}，他彻底考验并了解了他的邪恶性情，认为把自己置于君王的权力之下是不安全的，便留在旷野中。\Person{撒乌耳}几乎因愤怒而疯狂，因为他无法抓获他的女婿，便将他的女儿\Person{米加耳}（如上所述，她曾嫁给\Person{达味}）嫁给了\Person{法耳提}。\Person{达味}逃到培肋舍特人那里。\par

% structure: p0072 source=h2 target=subsection reason=relative-heading-rank; base=h2

\subsection{第三十六章}

那时\Person{撒慕尔}死了。\Person{撒乌耳}当培肋舍特人向他开战时，求问天主，但没有任何回答给他。于是，他借助一个被错谬之灵充满内脏的妇人，招上\Person{撒慕尔}并求问他。\Person{撒乌耳}被他告知，第二天他与他的儿子们必被培肋舍特人击败，阵亡在沙场。培肋舍特人于是在敌军境内扎营，第二天列阵，达味却被从营中遣走，因为培肋舍特人不相信他会忠于他们而攻击自己的同胞。但战斗发生后，希伯来人溃败，君王的儿子们阵亡；\Person{撒乌耳}从马上跌下，以免被敌人活捉，就伏剑自刎。关于他统治的年限，我们找不到任何确定的记载，除非在\BibleBook{宗徒}{大事录}中说他统治了四十年。关于这点，我倾向于认为，在宣讲中作此声明的\Person{保禄}，当时意欲也将\Person{撒慕尔}的年数包括在那位君王的统治年限内。然而，大多数记述这些时代的人都说他统治了三十年。我绝不能同意这种意见，因为当天主的约柜被移到\Place{克黎雅特耶阿陵}城时，\Person{撒乌耳}尚未开始统治，据说\Person{达味}王在约柜在那里二十年后，才将它从那城迁走。因此，既然\Person{撒乌耳}在那期间统治并死去，他掌权的时间必定非常短暂。关于\Person{撒慕尔}的时代，我们也发现同样的不清楚，他生于\Person{厄里}作司铎时，据说在年老时履行了司铎的职责。然而，有些记载这些时代的人（因为圣史几乎没有记录他的年数），但大多数说他统治了百姓七十年。然而，我未能发现这种假设有何权威根据。在如此多的错误说法中，我们依据\BibleBook{}{编年纪}的记载，因为我们相信它（如上所述）取自\BibleBook{宗徒}{大事录}，我们重申\Person{撒慕尔}和\Person{撒乌耳}共同掌权四十年。\par

% structure: p0074 source=h2 target=subsection reason=relative-heading-rank; base=h2

\subsection{第三十七章}

撒乌耳既如此被剪除，当达味在培肋舍特地方得知他死亡的消息时，据记载他哭泣了，并显出了他情感的奇妙明证。随后他前往犹太的一个城镇赫贝龙，在那里再次受傅王油，接受了君王的称号。但撒乌耳王的军长阿贝乃尔轻视达味，立撒乌耳王的儿子依市巴耳为王。随后两王的将领之间发生了多次战斗。阿贝乃尔常被击败；但他在败逃中斩杀了达味方面军长约阿布的兄弟。约阿布因对此事感到悲伤，后来当阿贝乃尔向达味王投降时，便下令将他杀害，而达味王对此深感惋惜，因他的名誉因此受损。与此同时，希伯来人的几乎所有长老都公意将全国的主权授予了他；因为七年间他仅在赫贝龙为王。这样，他第三次受傅为王，年约三十岁。他成功击退了侵入其王国的培肋舍特人。就在那时，他将天主的约柜迁至熙雍，如上所述，约柜原在克黎雅特耶阿陵镇。当他有意为天主建造圣殿时，天主的答复是，这事要留给他的儿子。随后他在战争中征服了培肋舍特人，制服了摩阿布人，并降服了叙利亚，使其纳贡。他带回了大量的金铜战利品。接着，由于阿孟人的君王哈农所行的侮辱，一场战争爆发了。当叙利亚人再次背叛，与阿孟人结盟作战时，达味将军队的指挥权交给军长约阿布，而他自己则留在远离战场的耶路撒冷。\par

% structure: p0076 source=h2 target=subsection reason=relative-heading-rank; base=h2

\subsection{第三十八章}

此时，他以有罪的方式认识了贝特舍巴，一位容貌极美的妇人。据说她是名叫乌黎雅之人的妻子，此人当时正在营中。达味将他置于战事危险之处，使他被杀。这样，他将那已脱离婚姻束缚、却因通奸而怀孕的妇人添为妻室。随后，达味虽受纳堂先知严厉斥责后承认了自己的罪，却未能逃脱天主的惩罚。因他在数日内便失去了那从隐秘结合所生的儿子，且他的家室发生了许多可怕之事。最后，他的儿子阿贝沙隆对父亲举起不敬的武器，企图将他从王位上赶走。约阿布在战场上与他对阵，王曾恳求他手下留情，饶恕战败的青年；但他违命，以刀剑惩罚了弑父的企图。据说那次胜利对王是悲哀的：他天性仁爱，甚至希望他那弑父的儿子得到宽恕。这场战争几乎尚未结束，又一场战事兴起，由一名叫舍巴的将领发动，他煽动所有恶人武装起来。但整个骚乱因首领被杀而迅速平息。此后达味与培肋舍特人进行数次战斗，结果有利；所有外敌内乱皆被平定后，他平安地拥有一个极其兴盛的王国。那时他突然生出统计百姓的欲望，以了解其帝国的兵力；于是约阿布军长数点了百姓，共有一百三十万公民。达味随即后悔并痛悔此事，恳求天主宽恕他竟起了这样的念头，以为国力的强盛在于臣民众多，而非天主的恩宠。因此，一位天使被派遣到他面前，向他启示三种惩罚，并让他选择其一。当三年饥荒、三月逃避敌人和三日瘟疫摆在他面前时，他既逃避了逃亡与饥荒，便选择了瘟疫，几乎瞬间七万人丧命。于是达味看见那位右手击倒百姓的天使，便恳求宽恕，情愿独自代众人受罚，说那该被毁灭的是他，因为是他犯了罪。这样，百姓的刑罚被免除了；达味在他看见天使的地方建造了一座祭台献给天主。此后，因年迈病弱，他立了由乌黎雅的妻子贝特舍巴所生的儿子撒罗满为王国继任者。撒罗满由司祭匝多克傅以王油，在父亲尚在世时接受了君王称号。达味在位四十年后去世。\par

% structure: p0078 source=h2 target=subsection reason=relative-heading-rank; base=h2

\subsection{第三十九章}

\Person{撒罗满}在开始统治时，用城墙环绕了城市。当他睡着时，\Person{天主}显现站在他旁边，赐给他所渴望的任何事物的选择。但他只求赐予他智慧，认为其他一切都没什么价值。因此，当他从睡眠中醒来，站在\Person{天主}的圣所前，他证明了\Person{天主}所赐予他的智慧。因为有两个妇人同住一所房屋，同时生了男婴；三天后夜里其中一个孩子死了，死婴的母亲趁另一妇人熟睡时，偷偷地调换了孩子，取走了活婴。于是她们之间起了争执，事情最终被带到王面前。因为没有证人，在双方都否认罪责的情况下，裁决是件难事。于是\Person{撒罗满}运用他天赋的神圣智慧，下令将孩子杀死，将尸体分给两个有争议的声称者。果然，当其中一人默认了这个判决，而另一人宁愿放弃孩子也不愿他被肢解时，\Person{撒罗满}从这妇人所表现的情感推断她是真正的母亲，便将孩子判给了她。旁观者不禁对这一裁决赞叹不已，因为他以如此巧妙的方式揭示隐藏的真理。因此，邻国的君王出于对他才能和智慧的钦佩，纷纷寻求他的友谊和结盟，准备执行他的命令。\par

% structure: p0080 source=h2 target=subsection reason=relative-heading-rank; base=h2

\subsection{第四十章}

仰仗这些资源，\Person{撒罗满}开始建造一座极其宏伟的圣殿献给\Person{天主}，为此筹备了三年资金，并在其统治第四年左右奠基。这是在希伯来人离开埃及后大约第五百八十八年，尽管在《\WorkTitle{列王纪第三卷}》中记载为四百四十年。这绝不准确；因为按照我上面给出的年代顺序，也许我应该算更少的年份而不是更多的。但我毫不怀疑，真相因抄写者的疏忽而被篡改，尤其是经历了如此多的世代，而不是神圣的作者犯错。同样，对于这本小作品，我们相信也会发生这种情况：由于誊录者的疏忽，那些我们并非不谨慎编纂的内容会被损坏。那么，\Person{撒罗满}从动工起在第二十年完成了圣殿的建造。然后，他就在那地方献了祭，并做了祈祷，以此祝福百姓和圣殿，\Person{天主}对他说话，宣告说，如果他们任何时候犯罪离弃\Person{天主}，他们的圣殿将被夷为平地。我们看到这在很久以前已经应验，我们将在适当的时候叙述事件的连贯顺序。在此期间，\Person{撒罗满}财富丰裕，事实上是有史以来所有君王中最富有的。但是，正如在这种情况下常发生的，他从财富沉沦到奢侈和邪恶，违背\Person{天主}的禁令，与外邦女子通婚，直到有七百个妻子和三百个妃嫔。结果，他为他们按照她们民族的样式设立了偶像，供她们献祭。\Person{天主}因这些行为远离他，严厉责备他，并告诉他作为惩罚，他的王国大部分将从其子手中夺走，赐给一个仆人。这事果然发生了。\par

% structure: p0082 source=h2 target=subsection reason=relative-heading-rank; base=h2

\subsection{第四十一章}

因为，\Person{撒罗满}在位第四十年去世后，其子\Person{勒哈贝罕}在十六岁继承父位，一部分百姓因受冒犯而背叛了他。因为，他们请求减轻\Person{撒罗满}加在他们身上的沉重赋税，他拒绝了这些恳求者的请求，从而失去了全体百姓的欢心。于是，经一致同意，政权交给了\Person{雅洛贝罕}。他出身中等家庭，曾在\Person{撒罗满}手下服役一段时间。但当王得知，\Person{天主}已通过先知\Person{阿希雅}的答复将希伯来人的主权应许给他时，便密谋要除灭他。\Person{雅洛贝罕}因恐惧而逃往埃及，并在那里娶了一位王族女子为妻。但当他终于听到\Person{撒罗满}的死讯时，他返回故土，并按百姓的意愿，如我们上面所说，执掌了政权。然而，有两个支派，即\Place{犹大}和\Place{本雅明}，仍然受\Person{勒哈贝罕}管辖；他从这些支派中准备了三万人的军队。但当两军前进时，百姓因\Person{天主}的话语被指示不要作战，因为\Person{雅洛贝罕}是受天命得到王国的。于是军队藐视王的命令，解散了，而\Person{雅洛贝罕}的势力增强了。但由于\Person{勒哈贝罕}占据\Place{耶路撒冷}，百姓习惯于在那里\Person{撒罗满}所建的圣殿中向\Person{天主}献祭，\Person{雅洛贝罕}担心他们的宗教情感会使百姓疏远他，便决意用迷信充满他们的思想。于是，他在\Place{贝特耳}和\Place{丹}各设立了一只金牛犊，让百姓在那里献祭；并且越过\Place{肋未}支派，他从百姓中立了司祭。但这种令\Person{天主}憎恶的罪孽招来了谴责。此后两王之间频繁交战，他们在不确定的条件下各自保有王国。\Person{勒哈贝罕}在位第十七年末去世。\par

% structure: p0084 source=h2 target=subsection reason=relative-heading-rank; base=h2

\subsection{第四十二章}

在\Person{阿彼雅}的儿子继位，在\Place{耶路撒冷}统治了六年，虽然\WorkTitle{编年纪}说他统治了三年。他的儿子\Person{阿撒}继位，是\Person{达味}的第五代，即他的曾曾孙。他是一位虔诚敬拜天主的人；因为摧毁了偶像的祭坛和木偶，他清除了他父亲不忠的痕迹。他与\Place{叙利亚}王结盟，并借助他的帮助严重地打击了\Person{雅洛贝罕}的王国，当时\Person{雅洛贝罕}的王国由他的儿子统治，且常常在战胜敌人后携获战利品。四十一年后他去世，患了脚疾。他被归咎于三种罪：第一，他过于信赖与\Place{叙利亚}王的结盟；第二，他把斥责他的天主先知囚禁；第三，当他患脚疾时，他寻求医治，不是求天主，而是求医师。在他统治初期，十个支派的王\Person{雅洛贝罕}去世，将王位留给他的儿子\Person{纳达布}。他因邪恶的行为，且因自己和父亲的行为而成为天主所憎恶的人，统治不到两年，他的儿女因不配而被剥夺了统治权。他的继位者是\Person{巴厄沙}，\Person{阿希雅}的儿子，他也证明了自己与天主同样疏远。他在位第二十六年去世，权力传给他的儿子\Person{厄拉}，但统治不到两年。因为他的骑兵队长\Person{齐默黎}在宴席中杀了他，篡夺了王位——一个同样为天主和人所憎恶的人。一部分民众反叛他，王权被授予一个叫\Person{塔木尼}的人。但\Person{齐默黎}在他之前统治了七年，同时与他一起统治了十二年。在\Person{阿撒}去世后，他的儿子\Person{约沙法特}开始统治犹大支派的一部分，他是一个因虔诚美德而非常著名的人。他与\Person{齐默黎}和平相处；他统治二十五年后去世。\par

% structure: p0086 source=h2 target=subsection reason=relative-heading-rank; base=h2

\subsection{第四十三章}

在他在位期间，\Person{阿哈布}，\Person{敖默黎}的儿子，是十个支派的王，其不虔敬超过任何人。因为他娶了\Place{漆冬}王\Person{巴舍}的女儿\Person{依则贝耳}，并为巴耳偶像设立了祭坛和木偶，并杀害了天主的先知。那时，先知\Person{厄里亚}藉祈祷封闭了天，使天不降雨在地上，并告诉国王，好让他在不虔敬中知道自己就是灾祸的原因。因此，天上的雨水被阻，整个大地被太阳烤焦，不供给人和牲畜食物，先知甚至一度濒临饿死。那时，他退到荒野，靠乌鸦给他送食物维持生命，附近的小溪为他供水，直到溪水干涸。随后，他蒙天主指示，前往\Place{匝尔法特}镇，投宿在一个寡妇家。当他饥饿地向她乞食时，她抱怨说只有一把面和一点油，用完这些后她和孩子们就要等死。但当\Person{厄里亚}以天主的话应许说坛里的面不会减少，瓶里的油不会缺少时，那女人不再犹豫，相信了先知要求信德的话，并得到了所应许的应验，因为每日的增加补充了每日的消耗。同时，\Person{厄里亚}使这位寡妇的死去的儿子复活。然后，他遵天主的命令去见国王，斥责了他的不虔敬，并命令所有民众聚集到他那里。民众迅速聚集，偶像和木偶的祭司大约四百五十人也应召而来。于是他们之间发生了争论，\Person{厄里亚}陈述天主的荣耀，而他们坚持自己的迷信。最后他们同意做这样的试验：如果从天上降下的火烧尽了任何一方所杀的牺牲，那行奇迹的宗教就是真宗教。于是，祭司们杀了一头牛犊，开始呼求巴耳偶像；在枉费了呼求之后，他们无言地承认了他们的神无能为力。然后\Person{厄里亚}嘲笑他们说：\Quote{你们要高声呼喊，免得他睡着了，这样你们才把他从沉睡中唤醒。}那些可怜的人只能战栗，彼此嘟囔，但仍然等着看\Person{厄里亚}会做什么。于是，他杀了一头牛犊，放在祭坛上，首先用水浸满了圣地；然后呼求主的名，火从天上降下，在众人眼前，同时烧尽了水和祭品。这时，民众真的俯伏在地，承认天主，咒骂偶像；最后，在\Person{厄里亚}的命令下，不虔敬的祭司被逮捕，被带下到溪边，在那里被处死。先知跟随国王从那里返回；但王后\Person{依则贝耳}设法取他的性命，他便退到更偏远的地方。在那里，天主对他说话，告诉他还有七千人没有向偶像屈膝。这对\Person{厄里亚}是惊人的宣告，因为他原以为只有自己保持了不受不虔敬的玷污。\par

% structure: p0088 source=h2 target=subsection reason=relative-heading-rank; base=h2

\subsection{第四十四章}

那时，\Place{撒玛黎雅}王\Person{阿哈布}觊觎\Person{纳波特}邻近自己田地的葡萄园。\Person{纳波特}不愿卖给他，因\Person{依则贝尔}的诡计，\Person{纳波特}被除灭。这样\Person{阿哈布}便占有了那葡萄园，但同时据说他对\Person{纳波特}的死感到后悔。他承认了自己的罪，身穿苦衣行补赎，以此避开了威胁的惩罚。\Place{叙利亚}王率领大军，与三十二个王结盟，进入\Place{撒玛黎雅}境内，开始围攻城和王。当时被围者处境极度困窘，\Place{叙利亚}王提出战争条件——若他们交出金银和妇女，他就饶他们的命。但面对如此不义的条件，忍受最极端的苦难似乎更好。当所有人都绝望于安全时，一位奉天主派遣的先知来到王那里，鼓励他出战，当他犹豫时，从多方面坚定他的信心。于是他们出击，敌人溃败，获得了大量战利品。但一年后，\Place{叙利亚}王重整兵力回到\Place{撒玛黎雅}，渴望复仇，但再次被击败。那一战\Place{叙利亚}人死了十二万；王被赦免，他的王国和先前地位得以归还。然后\Person{阿哈布}被先知以天主的话责备，因为他滥用了天主的仁慈，饶恕了交在他手中的敌人。因此，三年后，\Place{叙利亚}王向希伯来人发动战争。\Person{阿哈布}听信了某个假先知的建议，出战，他藐视先知\Person{米该亚}的话，并将他下狱，因为先知曾警告他这场战斗于他不利。于是，\Person{阿哈布}在那场战斗中被杀，将王国留给了他的儿子\Person{阿哈齐雅}。\par

% structure: p0090 source=h2 target=subsection reason=relative-heading-rank; base=h2

\subsection{第四十五章}

他身患疾病，派了一些仆人去求问偶像关于他的康复，\Person{厄里亚}奉天主指示在路上遇见他们，责备了他们，并叫他们告诉王，他必死于那病。于是王下令捉拿\Person{厄里亚}，带到他面前，但派去的人都被天火烧死。王如先知预言般死了。他的兄弟\Person{约兰}继位，治理了十二年。而在犹大一方，\Person{约沙法特}王驾崩，他的儿子\Person{约兰}获得王位十八年。他娶了\Person{阿哈布}的女儿为妻，证明自己更像他的岳父而非他的父亲。之后，他的儿子\Person{阿哈齐雅}继位。他在位期间，\Person{厄里亚}被提升天。同时，他的门徒\Person{厄里叟}行了许多大能奇事，众所周知，无需我笔描述。他使寡妇的儿子复活，洁净了\Place{叙利亚}的麻风病人，在饥荒时使敌人败退，城内物资丰富，为三支军队提供了水，用一点油因油大量倍增还清了一个妇人的债，她获得了足够的生活资料。在他之时，如我们所说，\Person{阿哈齐雅}治理犹大，而\Person{约兰}则治理以色列；他们之间结盟。因为他们联合对\Place{叙利亚}人和\Person{耶户}作战，\Person{耶户}曾被先知膏立为十支派的王；他们一同出战，双双战死于同一场战斗。\par

% structure: p0092 source=h2 target=subsection reason=relative-heading-rank; base=h2

\subsection{第四十六章}

\Person{耶户}取得了\Person{约兰}的王国。\Place{犹大}\Person{阿哈齐雅}死后（他在位一年），他的母亲\Person{阿塔里雅}夺取了最高权力，剥夺了她孙子\Person{约阿士}（当时还是小孩）的统治权。但八年后，通过司铎和人民，权力得以归还给他，他的祖母被放逐。他统治初期，最热心于敬拜天主，并花费巨资装饰圣殿；但后来，被首领们的谄媚所败坏，他们过分尊崇他，招致了愤怒。因为\Place{叙利亚}王\Person{哈匝耳}向他开战；战事不利，他便用圣殿的金子购买和平。然而并未如愿；因他所做之事，他在位第四十年被自己的人民杀害。他的儿子\Person{阿玛责雅}继位。在\Place{以色列}一方，\Person{耶户}死后，他的儿子\Person{约阿哈次}开始为王，因恶行不悦于天主，作为惩罚，他的王国被\Place{叙利亚}人蹂躏，直到因天主的仁慈，敌人被击退，当地居民恢复原位。\Person{约阿哈次}去世后，将王国留给了他的儿子\Person{约阿士}。他向犹大王\Person{阿玛责雅}发动内战；得胜后，掠获大量战利品带回国内。据记载，这事发生在\Person{阿玛责雅}因犯罪而受罚时，因为他战胜了\Place{厄东}人的领土后，接受了那民族的偶像。据\WorkTitle{列王纪}记载，他统治了九年。但在圣经\WorkTitle{编年纪}以及\WorkTitle{尤西比乌斯编年史}中，他被认为在位二十九年；在\WorkTitle{列王纪}中容易察觉的计算方式无疑导致这一结论。因为\Person{雅洛贝罕}开始作\Place{以色列}王是在\Person{阿玛责雅}在位第八年，他统治四十一年，最终在\Person{阿玛责雅}的儿子\Person{乌齐雅}在位第四年去世。按此计算，\Person{阿玛责雅}的在位年数是二十八年。因此，我们遵循这一点，因为本书的目的是按照时间顺序保持日期，我们采纳了\WorkTitle{编年纪}的权威。\par

% structure: p0094 source=h2 target=subsection reason=relative-heading-rank; base=h2

\subsection{第四十七章}

于是，\Person{阿玛息雅}的儿子\Person{乌齐雅}继承了他。因为，在十个支派方面，\Person{约阿士}到了他日子的尽头，将王位让给了他的儿子\Person{雅洛贝罕}；之后，他的儿子\Person{则加黎雅}又开始作王。关于这些君王，以及所有在十个支派方面统治\Place{撒玛黎雅}的，我们认为没有必要记录年代，因为为了简洁，我们省略了一切多余的事；我们认为应该仔细追溯年代，特别是为了认识那批犹太人的时间，他们比另一批后来被掳，因此作为王国经历了更长的时间。\Person{乌齐雅}获得\Place{犹大}王位后，将首要的注意力放在认识主上，大量使用先知\Person{则加黎雅}（据说\Person{依撒意亚}也是在这位王下开始说预言的）；因此，他同邻邦作战，得到理所当然的胜利，也征服了\Place{阿剌伯}人。他已经用他名的恐惧震撼了\Place{埃及}；但因得意而骄横，他胆敢做禁止之事，向天主献香，这件事已成惯例只可由司铎做。于是，他被司铎\Person{阿匝黎雅}斥责，被迫离开圣地，他勃然大怒，但最终退出时，全身长了大癞。在这种疾病下，他结束了日子，在位五十二年后。然后王位给了他儿子\Person{约堂}；据说他非常虔诚，治理成功：他征服了\Place{阿孟}人的国家，强迫他们进贡。他在位十六年，他的儿子\Person{阿哈次}接替他。\par

% structure: p0096 source=h2 target=subsection reason=relative-heading-rank; base=h2

\subsection{第四十八章。}

据说大约在这时候，\Place{尼尼微}人非凡的信德得以彰显。那座城市，由\Person{闪}的儿子\Person{阿苏尔}在古时建立，是\Place{亚述}王国的首都。当时它满有众多居民，养育着十二万人，并充满邪恶，正如在庞大的人群中常见的那样。天主因他们的罪恶而感动，命令先知\Person{约纳}从\Place{犹大}去，宣告那城将要毁灭，如同古时\Place{索多玛}和\Place{哈摩辣}被天火焚烧一样。但先知拒绝了那宣讲的职务，并非出于悖逆，而是出于预见，使他能看见天主因百姓的悔改而和好；他上了一艘开往\Place{塔尔史士}的船，方向完全相反。但船行至深海时，水手因海浪猛烈，用签询问谁是那苦难的原因。签落在\Person{约纳}身上，他被投入海中，作为平息风暴的祭品，并被一头海怪——鲸鱼捕获吞下。三天后，他被吐在\Place{尼尼微}人的海岸上，他按所受的命令宣讲，即那城将在三天内毁灭，作为百姓罪恶的惩罚。先知的声音被听了，不是像古时\Place{索多玛}那样虚伪；立刻，按国王的命令和榜样，全体百姓，甚至新生的婴儿，都被命令禁食禁饮：那地方的驮兽和各种动物，因饥渴所迫，呈现出与人同哀哭的样子。如此，那威胁的灾祸被避免了。\Person{约纳}向天主抱怨他的话没有应验，得到的回答是，悔改者永不蒙拒绝宽恕。\par

% structure: p0098 source=h2 target=subsection reason=relative-heading-rank; base=h2

\subsection{第四十九章。}

但在\Place{撒玛黎雅}，王\Person{则加黎雅}，他非常邪恶，我们上面提到他占了王位，被一个叫\Person{塞拉}的人杀死，\Person{塞拉}夺取了王国。他反过来又死于\Person{玛纳}的背叛，\Person{玛纳}只是重复了前任的行为。\Person{玛纳}掌握了他从\Person{塞拉}夺得的政府，并传给他儿子\Person{帕刻}。但一个同名的人杀了\Person{帕刻}，夺取了王国。不久，被\Person{奥西}除掉，他因夺取王位的同一罪行丧失了主权。这人比他以前所有的王更不敬畏天主，为自己招来天主的惩罚，并为他的民族带来永久的掳。因为\Place{亚述}王\Person{撒耳玛纳撒}与他作战，他战败后成为属国。但当他秘密谋划造反，并向当时占据\Place{埃及}的\Place{厄提约丕雅}王求援时，\Person{撒耳玛纳撒}发现后，将他戴上手铐脚镣投入监狱，永不释放，同时摧毁了城市，并将全体百姓掳到自己的王国，\Place{亚述}人被安置在敌国看守。因此那地区被称为\Place{撒玛黎雅}，因为在\Place{亚述}语中，守卫被称为撒玛黎雅人。这些移民中有许多人接受了犹太宗教的神圣礼仪，而其他人则留在异教的错误中。在这场战争中，\Person{多俾亚}被掳。但在两个支派方面，\Person{阿哈次}因不虔敬而不悦于天主，他发现自己在与邻邦的战争中多次失败，便决定敬拜异教的神，无疑因为靠他们的帮助他的敌人在多次战斗中获胜。他带着这罪恶的恶意结束了他的日子，在位十六年。\par

% structure: p0100 source=h2 target=subsection reason=relative-heading-rank; base=h2

\subsection{第五十章。}

他的儿子\Person{希则克雅}继位，品性与父亲大不相同。因为在位之初，他劝勉百姓和司铎敬拜天主，用许多话向他们讲论，表明他们屡次受主惩戒后如何获得了怜悯，以及十个支派最终被掳——正如最近所发生的那样——如今正在为他们的不虔敬受罚。他又补充说，他们应当谨慎小心，免得遭受同样的厄运。于是，众人的心都转向宗教，他任命肋未人和众司铎依照法律献祭，并安排庆祝长久被忽略的逾越节。当圣日临近时，他派遣使者走遍全国，宣布特别集会的日子，以便任何在十个支派被掳后仍留在撒玛黎雅的人，都能聚集参与这神圣的礼仪。因此，在一次极为盛大的集会中，众人以公共的喜乐度过了这圣日，经过漫长的间隔，正当的宗教礼仪由\Person{希则克雅}恢复了。随后，他以处理神事同样的勤勉处理军事事务，在多次战役中击败了培肋舍特人。直到亚述王\Person{散乃黑黎布}率领大军进入他的领土向他开战；那时，国土未经抵抗便被蹂躏，随后王城被围困。因为\Person{希则克雅}兵力不足，不敢与他交战，只守在城墙内自保。亚述王在城门前大声恐吓，威胁要毁灭，并劝降，喊叫说\Quote{希则克雅信靠天主是枉然，因为他反而是奉天主的命拿起武器；他既然征服了万民，又推翻了撒玛黎雅，无人能逃脱，除非王速速投降以确保安全。}在此情况下，\Person{希则克雅}信赖天主，求问先知\Person{依撒意亚}，从他的答复中得知敌人不足为惧，天主的助佑必不缺少。果然不久之后，\Person{塔辣卡}，\Place{埃塞俄比亚}王，入侵了亚述王国。\par

% structure: p0102 source=h2 target=subsection reason=relative-heading-rank; base=h2

\subsection{第五十一章}

这消息促使\Person{散乃黑黎布}返回保卫自己的领土，放弃了这场战争，同时抱怨并嚷着说胜利者被人夺去了胜利。他还写信给\Person{希则克雅}，用许多侮辱性的话语宣称，他处理完自己的事务后，将迅速回来毁灭\Place{犹大}。但\Person{希则克雅}毫不动摇，据说他向天主祈祷，求天主不容这样的人如此傲慢而不受惩罚。于是当夜，一位天使袭击了亚述人的营寨，导致数万人死亡。王惊恐地逃到\Place{尼尼微}城，在那里被自己的儿子杀死，结局恰如其人。与此同时，\Person{希则克雅}身体患病，卧病在床。当\Person{依撒意亚}以主的话告诉他生命将尽时，据说王哭泣了；于是他的寿命延长了十五年。这些岁月终结后，他在位第二十九年去世，将王国留给他的儿子\Person{默纳舍}。他远不如父亲，离弃了天主，转而行不敬神的崇拜；为此受罚被交在亚述人手中，因受苦才不得不承认自己的错误，并劝勉百姓离弃偶像，敬拜天主。他没有做出什么值得一提的事，在位五十五年。然后他的儿子\Person{阿孟}获得王位，但只统治了两年。他继承了父亲的不虔敬，表现出对天主的漠不关心：被朋友的一些诡计所诱，丧命了。\par

% structure: p0104 source=h2 target=subsection reason=relative-heading-rank; base=h2

\subsection{第五十二章}

政权随后传给他的儿子\Person{约史雅}。据说他非常虔诚，以极大的细心处理神事，大大得益于司铎\Person{希耳克雅}的帮助。他阅读了一本写有天主话语的书——是司铎在圣殿中找到的——书中说希伯来民族将因他们屡次的不虔敬和亵渎而灭亡；他通过虔诚的祈祷和不断的眼泪，向天主恳求，避免了即将来临的倾覆。当通过女先知\Person{胡耳达}得知这恩宠已赐予他时，他便更加用心地实践对天主的敬拜，因为他如今欠了天主的恩情。于是，他焚烧了所有先前诸王因迷信而奉献给偶像的器皿。因为邪妄的习俗已盛行到如此地步，他们竟向太阳和月亮敬拜，甚至为这些假神建造金属神龛。\Person{约史雅}将这些碾成粉末，并杀死了邪神庙的司祭。他甚至没有放过不虔敬者的坟墓；值得注意的是，这正应验了先知古时所预言的。他在位第十八年庆祝了逾越节。大约三年后，他出兵攻打向亚述人开战的埃及王\Person{乃苛}，在军队正式交战前，被一支箭射伤。他被送回城中，因伤而死，在位共二十一年。\par

% structure: p0106 source=h2 target=subsection reason=relative-heading-rank; base=h2

\subsection{第五十三章}

他的儿子\Person{约阿哈次}随后获得王国，在位三个月，因不虔而被注定被掳。\Place{埃及}王\Person{尼苛}将他捆绑并掳去，不久后，他仍在囚禁中结束了一生。犹太人被要求每年进贡，而胜利者按己意立了一位王。他名叫\Person{厄里雅金}，但后来改名为\Person{约雅金}。他是\Person{约阿哈次}的兄弟，\Person{约史雅}的儿子，但更像其兄而不像其父，因不虔而得罪了天主。因此，当他臣服于\Place{埃及}王并为此纳税时，\Place{巴比伦}王\Person{拿步高}夺取了\Place{犹太地}，作为胜者凭战争权占领了三年。因\Place{埃及}王当时退让，双方帝国边界划定，约定犹太人应归属\Place{巴比伦}。这样，\Person{约雅金}在位十一年后结束统治，传位给与他同名的儿子，而此子激起了\Place{巴比伦}王的愤怒（天主无疑统御万有，决意将犹太民族交付于被掳和毁灭），\Person{拿步高}率军进入\Place{耶路撒冷}，将城墙和圣殿夷为平地。他还掠走了大量金银，包括公家和私人的圣器，以及所有成年男女，只留下那些因虚弱或年老而给征服者带来麻烦的人。这群无用的残众被指派劳作和耕种田地，如同奴隶，以免土地荒芜。他们之上立了一位王，名叫\Person{漆德克雅}，但只给他王名的虚影，却剥夺了一切实权。至于\Person{约雅金}，他只掌权三个月。他与百姓一同被掳至\Place{巴比伦}，投入监牢；但三十年后获释，被王视为友人，与他同桌共饮、参与谋划，最终去世，也算因不幸得以解除而有些许安慰。\par

% structure: p0108 source=h2 target=subsection reason=relative-heading-rank; base=h2

\subsection{第54章。}

\Emph{同时}，\Person{漆德克雅}，这无用残众之王，虽无权力，却性本不忠，轻慢天主，不明白被掳是因民族的罪恶加于他们身上，最终导致他遭受所能承受的最后磨难，触怒了王。因此，九年后，\Person{拿步高}向他开战，逼他撤入城墙之内，围攻了三年。此时，他咨询了先知\Person{耶肋米亚}，先知早已多次宣告被掳将临于城，问是否还有一线希望。但先知并非不知上主的愤怒，因常有人向他提出同样的问题，终于回答，宣告君王将受特别的惩罚。于是\Person{漆德克雅}勃然大怒，命人将先知投入监牢。然而不久，他就为这残忍之举感到后悔，但因犹太人的首领（他们自古就有迫害义人的行径）反对他，他不敢释放这无罪之人。在同样的胁迫下，先知被投入一个极深的坑中，因其污秽和肮脏而令人作呕，且发出致命的恶臭。此举是为了不让他死于寻常的死法。但王虽然不虔，却比司铎们显得更仁慈一些，命人将先知从坑中拉出，重新关入监牢。与此同时，敌军的压力和匮乏开始紧逼被围者，一切可食之物都被吃尽，饥荒牢牢抓住了他们。守军因缺粮而精疲力竭，城被攻取并焚烧。正如先知所预言的，王被剜了眼睛，并被掳至\Place{巴比伦}。\Person{耶肋米亚}因敌人的仁慈而被从监牢中提出。当\Person{拿步匝辣当}，王室王子之一，正与其他俘虏一同将他带走时，先知获准选择：要么留在他那荒凉废弃的故土，要么随他同去，享有最高荣誉。\Person{耶肋米亚}宁愿留在故土。\Person{拿步高}掳走百姓后，立\Person{革达里雅}为征服者留下之人的总督；\Person{革达里雅}属同一民族。此举或出于战争状况，或出于对战利品的极度厌倦。然而，他没有给他王权的标志，甚至连总督的名号也没有，因为统治这少数可怜之人实在毫无荣耀可言。\par

\SourceRecord{Translated by Alexander Roberts.；Nicene and Post-Nicene Fathers, Second Series；Vol. 11.；Edited by Philip Schaff and Henry Wace.；Buffalo, NY: Christian Literature Publishing Co.,；1894.；<http://www.newadvent.org/fathers/35051.htm>.}

% structure: p0001 source=h1 target=section reason=page-title

\section{\WorkTitle{圣史（第二卷）}}

% structure: p0002 source=h2 target=subsection reason=relative-heading-rank; base=h2

\subsection{第一章}

被掳时期因先知们的预言和事迹而熠熠生辉，尤其是\Person{达尼尔}坚守法律的不屈精神，以及\Person{苏撒纳}藉天主的智慧而获救，并因它成就的其他事，这些我们将按顺序叙述。\Person{达尼尔}在\Person{约雅金}王在位时被掳，被带到\Place{巴比伦}，当时他还是个很小的孩子。后来，因他容貌俊美，被准许在王的仆从中占有一席之地，同他一起的还有\Person{阿纳尼雅}、\Person{米沙耳}和\Person{阿匝黎雅}。但当王命令供给他们更精美的食物，并责成太监\Person{阿斯凡纳}督办此事时，\Person{达尼尔}记念祖先的传统，禁止他吃外邦王桌上的食物，便恳求太监只准他食用蔬菜。\Person{阿斯凡纳}反对说，随后出现的消瘦会暴露王命被违背的事实。但\Person{达尼尔}信赖天主，应许说，靠吃蔬菜生活，会比享用王的美食容貌更俊美。他的话应验了，以致那些受公家供养的人的面容，被认为远不及\Person{达尼尔}及其同伴。因此，他们被王提升，得享尊荣和恩宠，不久，凭其谨慎和智慧的行为，被置于所有最接近王的人之上。大约同时，\Person{苏撒纳}，一个名叫\Person{约雅金}的人的妻子，容貌出众，被两位长老垂涎；当她不肯听从他们不洁的提议时，便遭诬告。这两位长老报告说，在一僻静之处发现一个青年与她在一起，但因他们年老体弱，那青年凭其年轻敏捷逃脱了他们的手。于是，信了这两位长老的话，\Person{苏撒纳}被民众的判决定罪。当她依法被带去受罚时，当时十二岁的\Person{达尼尔}，先责备犹太人是将无辜者交于死亡，然后要求将她带回重审，重新听取她的案件。因为在场的众多犹太人认为，一个年纪如此幼小、不足以令人尊敬的孩子，若非受天主感动，决不敢如此大胆，便允准了他的请求，重新开庭。于是再次进行审讯；\Person{达尼尔}被允许在长老中就座。随即，他命令将两个控告者彼此隔离，分别询问他们，是在哪种树下发现那奸妇的。因他们所答的不同，其谎言被揭露：\Person{苏撒纳}获释；而那两位使无辜者陷于危险的长老，被判处死刑。\par

% structure: p0004 source=h2 target=subsection reason=relative-heading-rank; base=h2

\subsection{第二章}

那时，\Person{拿步高}做了一个梦，因其预示未来的洞察力而令人惊奇。因他自己无法解释梦境，便派人召来\OriginalTerm{加色丁人}{Chaldæans}，他们被认为能藉法术和祭牲的内脏知晓隐秘之事并预言未来，为求解释。他随即担心，以免他们如常人一般，从梦境中得出的不是实情，而是取悦王的话，便隐瞒了所见之事，要求他们，若真有占卜之力，就当将梦本身告诉他；说，如果他们先借述梦证明自己的技能，他便相信他们的解释。但他们婉拒尝试如此重大之难事，承认此类事非人力所及。王大怒，因他们以虚假的占卜自居，以谬误愚弄众人，而今迫于此事不得不承认自己并无假装的那种知识，便以王谕揭露他们；所有操此技艺者均被公开处死。\Person{达尼尔}闻知此事，便对一位最接近王的人说话，应许说明那梦，并提供解释。此事奏报王，\Person{达尼尔}被召来。这奥秘已由天主启示给他；于是他叙述王的异象，并解释之。但此事要求我们陈述王的梦及解释，并说明他话语借随后之事应验。王在梦中看见一个像：头是金的，胸膛和手臂是银的，腹部和大腿是铜的，腿是铁的，脚部一部分是铁，一部分是泥。但铁和泥混合时，彼此不能粘合。最后，一块未经人手凿成的石头将像砸碎，整个像化为尘土，被风吹走。\par

% structure: p0006 source=h2 target=subsection reason=relative-heading-rank; base=h2

\subsection{第三章}

\Emph{因此}，正如先知对这件事的解释，所见的像代表了世界。金头是加色丁人的帝国；因为我们知道它是第一个也是最富有的。银的胸膛和双臂代表第二个王国；因为居鲁士在征服加色丁人和玛待人之后，将帝国交给了波斯人。铜的肚腹表明第三主权；我们看到这已应验，因为亚历山大从波斯人手中夺取了帝国，为马其顿人赢得了主权。铁腿指向第四强权，这被理解为罗马帝国，它比之前所有的王国都更强大。但脚一部分是铁、一部分是泥的事实，表明罗马帝国将要分裂，以致永远不能合一。这也已应验，因为罗马国不是由一位皇帝统治，而是由几位皇帝，并且他们总是在彼此争吵，无论是通过实际战争还是派系斗争。最后，泥和铁混合在一起，但在本质上从未完全结合，预示着人类未来那些混合体，他们彼此分歧，尽管表面上联合。因为很明显，罗马领土被外族或叛军占据，或者已被交给那些在和平假象下投降的人。同样明显的是，野蛮民族，尤其是犹太人，已混杂在我们的军队、城市和行省中；我们看到他们生活在我们的中间，却绝不接受我们的习俗。先知们宣称这是末期。但那块非人手凿出来的石头，打碎了金、银、铜、铁和泥，是基督的预像。因为他，不是在人的条件下出生的（因他不是由人的意愿所生，而是由天主的意愿所生），将消灭那个存在世上王国的世界，并建立另一个不朽且永恒的王国，即为圣人预备的未来世界。有些人的信德仍对此点犹豫，他们不相信将来的事，虽然对过去的事深信不疑。于是，君王赏赐了\Person{达尼尔}许多礼物，派他管理\Place{巴比伦}和整个帝国，并给予他最高的荣誉。通过他的影响，\Person{阿纳尼雅}、\Person{阿匝黎雅}和\Person{米沙耳}也被提升到最高的尊严和权力。大约在同一时间，\Person{厄则克耳}著名的预言出现了，未来的奥秘和复活向他启示。他的书是一部非常重要的著作，值得细读。\par

% structure: p0008 source=h2 target=subsection reason=relative-heading-rank; base=h2

\subsection{第四章}

但在\Place{犹大}，正如我们上文所述，\Place{耶路撒冷}毁灭后，\Person{哥多利亚}被立为总督，犹太人极为不满，因为他们认为征服者凭己意指定了一个非王族的人作统治者，于是以某个\Person{依市玛耳}为首，煽动可憎的阴谋，在宴席上以诡计杀害了\Person{哥多利亚}。然而，那些没有参与阴谋的人，想要为这一罪行报仇，急忙拿起武器对抗\Person{依市玛耳}。但当\Person{依市玛耳}得知毁灭临近时，他撇下他聚集的军队，仅带不超过八个同伴逃往\Place{阿孟人}那里。因此，恐惧降临到全体人民，生怕巴比伦王为了少数人的罪而毁灭所有人；因为除了\Person{哥多利亚}，他们还与他一同杀了许多加色丁人。于是，他们计划逃往\Place{埃及}，但先集体去见\Person{耶肋米亚}，请求他求问天主的旨意。\Person{耶肋米亚}便以天主的话劝告他们全部留在本国，告诉他们如果这样做，必蒙天主的能力保护，不会因巴比伦人而遭遇危险；但如果他们去\Place{埃及}，必在那里死于刀剑、饥荒和各种死亡。然而，群众照常显示出邪恶的倾向，不习惯于顺从有益的劝告和天主的能力，竟去了\Place{埃及}。\WorkTitle{圣经}对他们后来的命运保持沉默；而我未能发现关于此事的任何信息。\par

% structure: p0010 source=h2 target=subsection reason=relative-heading-rank; base=h2

\subsection{第五章}

在这个时期，\Person{拿步高}因兴旺而得意，为自己立了一座巨大的金像，命令人将它作为圣像来朝拜。当所有人都因普遍盛行的谄媚而心志败坏、热心地这样做时，\Person{阿纳尼雅}、\Person{阿匝黎雅}和\Person{米沙耳}避开了这种亵渎的仪式，因为他们深知这种崇敬只应归于天主。因此，按照国王的谕令，他们被视为罪犯，并为他们预备了火窑作为刑罚的手段，以便通过眼前的恐怖迫使他们朝拜金像。但他们宁愿被火焰吞没，也不愿犯这样的罪。于是他们被捆绑起来，扔进火中。但火焰却抓住了那些从事这可恶工作的人——他们正急切地把受害者推进火里；而奇妙的是——对所有人来说，若非亲眼目睹，简直难以置信——火焰根本没有触及那些希伯来人。围观者看见他们在火窑中行走，并唱着赞美天主的歌，同时还看见与他们一起的第四位，形貌似天使，\Person{拿步高}走近观看时，承认他是天主子。于是国王确信所发生的事中有天主的能力，便向全国发出通告，宣告所发生的奇迹，并承认只有天主才配受崇敬。不久之后，因出现在他面前的神视，以及随后从天上传来的声音指示，据说他做了补赎，放下王权，远离人群，仅靠野草维持生命。然而，他的帝国因天主的旨意而为他保留，直到时期满足，最终他正式承认天主，七年后恢复了王国和先前的地位。据记载，他在征服\Person{漆德克雅}（如我们上文所说，将他掳到巴比伦）之后，统治了二十六年，尽管我在神圣历史中并未找到此记载。但或许是因为我在探究许多问题时，在某位无名作者的作品中发现了这一评注，而该作品随时间推移被掺杂了内容，其中包含巴比伦诸王的年代。我认为不应忽略这一评注，因为它确实与编年史相符，因而其记载与我们一致，即通过诸王的继承（其年代由记录所载），直到居鲁士王元年，正好满了七十年，而这正是神圣历史中所说的从被掳到居鲁士时期的年数。\par

% structure: p0012 source=h2 target=subsection reason=relative-heading-rank; base=h2

\subsection{第六章}

\Person{拿步高}之后，王位传给他的儿子，我在编年史中看到其名为\Person{厄威耳默洛达客}。他在位第十二年去世，由其弟\Person{贝尔沙匝}继位。\Person{贝尔沙匝}在第十四年为他的首长和官员举行公宴时，命人将圣器（即\Person{拿步高}从耶路撒冷圣殿中取走、但从未用于王用、而是保存在库房中的器皿）取来。当所有男女，连同他的妻妾，在皇家宴会的奢华与放纵中使用这些器皿时，突然国王看见有手指在墙上写字，并注意到字母组成了文字。但没有人能读出这文字。国王因此惊慌，召来术士和\Person{加色丁人}。他们只是彼此嘀咕，一无所答。王后提醒国王，有一位名叫\Person{达尼尔}的希伯来人，曾为\Person{拿步高}揭示包含隐秘奥迹的梦，并因他的卓越智慧被擢升到最高荣誉。于是他被召来，读出了文字并解释其意：因国王亵渎了神圣的器皿而犯罪，毁灭将临于他，他的王国将归于\Person{玛待人}和\Person{波斯人}。这事随即发生。因为当夜\Person{贝尔沙匝}丧命，\Person{达理阿}——一位玛待人——夺取了他的王国。\Person{达理阿}见\Person{达尼尔}享有极高声誉，便立他为全帝国的首脑，这遵循了他前任诸王的判断。因为\Person{拿步高}也曾立他管理王国，\Person{贝尔沙匝}曾赐他紫袍和金链，并立他为王国中的第三位统治者。\par

% structure: p0014 source=h2 target=subsection reason=relative-heading-rank; base=h2

\subsection{第七章}

因此，那些与他一同掌权的人因嫉妒而激动，因为一个属于被掳民族的外国人竟被置于与他们平等的地位，他们迫使那已被谄媚腐蚀的国王颁布法令，在接下来的三十天内应向他（国王）献上神性的尊荣，并且任何人不得向神祈祷，除非向国王。达理阿很容易就被说服了，因为所有国王都愚蠢地为自己要求神性的尊荣。在这种情况下，达尼尔并非不了解所发生的事，也非不知道祈祷应当向天主而非人呈上，他被指控没有遵守国王的命令。尽管达理阿很不情愿——达尼尔一向为他所疼爱和悦纳——但掌权者们仍然占了上风，结果达尼尔被投入狮穴。然而，当他暴露在野兽面前时，并未遭受任何伤害。国王发现这事后，下令将他的控告者交给狮子。但他们却没有经历同样的遭遇，因为他们立刻被吞食，满足了猛兽的饥饿。达尼尔先前已很有名，如今更被视为杰出；国王撤销了先前的谕令，颁布了一道新谕令，规定要抛弃一切错误和迷信，敬拜达尼尔的天主。此外，还有达尼尔神视的记录，他在其中揭示了未来世代事件的次序，也包含了年数，在这期间他宣布基督将降临人间（正如已发生的那样），并清楚说明了假基督的将来来临。如果有人渴望探究这些要点，他将在\WorkTitle{达尼尔书}中找到更详尽的论述；我们的设计只是简要地连贯陈述事件。据记载，达理阿在位十八年；此后，阿斯提亚格开始统治玛代人。\par

% structure: p0016 source=h2 target=subsection reason=relative-heading-rank; base=h2

\subsection{第八章}

他的外孙居鲁士利用波斯人的军队将他驱逐出王国，因此大权转移到了波斯人手中。巴比伦人也落入了他的权力和统治之下。在他统治之初，他颁布公开谕令，允许犹太人返回本国，并且归还了拿步高从耶路撒冷圣殿掳去的圣器。于是，少数人返回了犹太；至于其余的人，我们无法查明他们是缺乏归回的愿望还是能力。那时，在巴比伦人中有一座非常古老的国王贝鲁斯的铜像，维吉尔也曾提到过它。这座像因民众的迷信而被视为神圣，居鲁士也受了司铎们诡计的欺骗，惯常敬拜它。他们声称这像会饮食，而他们自己却暗中偷走每日供奉给偶像的供物。居鲁士与达尼尔交好，问达尼尔为何不敬拜这座像，因为它似乎是生活天主的明显象征，能消耗供奉给它之物。达尼尔嘲笑这个人的错误，回答说这是不可能的，那铜制品——纯粹无感觉的物质——不可能享用饮食。于是国王下令召来司铎（他们约有七十人），并以威吓使他们恐惧，责问他们是谁消耗了所供奉之物，因为达尼尔，一位以智慧著称的人，坚持认为无感觉的偶像不可能做到。于是他们仗着自己准备好的诡计，吩咐照常献祭，并由国王封印殿宇，约定除非次日所有供物都被发现已被消耗，他们就会被处死；反之，如果发现相反的情况，达尼尔也要遭受同样的命运。于是殿宇被国王的印玺封住；但达尼尔事先在未通知司铎的情况下用灰覆盖了地板，以便他们的脚印能暴露那些进入者的秘密行踪。次日，国王进入殿宇，发现他下令供奉给偶像的那些东西已被取走。于是达尼尔借着出卖行踪的脚印揭露了秘密的欺诈，表明司铎们带着妻子儿女从下面开启的洞口进入殿宇，吃掉了供奉给偶像的东西。因此，他们全体都在国王的命令下被处死，而殿宇和偶像则被交付给达尼尔处置，并依照他的命令被销毁。\par

% structure: p0018 source=h2 target=subsection reason=relative-heading-rank; base=h2

\subsection{第九章}

在此期间，那些我们上文谈及的、蒙\Person{居鲁士}许可回归故土的犹太人，曾试图重建他们的城市和圣殿。但因他们人数稀少且贫穷，进展甚微；直到大约百年之后，当\Person{阿尔塔薛西斯}王统治波斯人时，他们最终被当地掌权者彻底阻止了建筑。因为在那时，\Place{叙利亚}和整个\Place{犹太地}，在波斯帝国治下由长官和总督统治。于是，这些人商议致信\Person{阿尔塔薛西斯}王，说不应当给予犹太人重建其城市的机会，以免他们固守本性，惯于统治他族，一旦恢复力量，便不肯顺从于外邦势力的治权。此计划得到王的认可后，城市的建筑工程便被叫停，推迟到达理阿王第二年。然而，我们要将这一时期中所有波斯王列出，以便年代顺序有条不紊地呈现。那么，在\OriginalTerm{玛待}{Mede}人达理阿（如上所述他统治了十八年）之后，\Person{居鲁士}执掌最高权力三十一年。他在征讨\OriginalTerm{西徐亚人}{Scythians}的战役中阵亡，时值\OriginalTerm{塔奎尼乌斯·苏佩布斯}{Tarquinius Superbus}在罗马开始统治后的第二年。\Person{居鲁士}的儿子\Person{冈比西斯}继位，在位八年。他以战争骚扰\Place{埃及}和\Place{埃塞俄比亚}，征服这些国家后，凯旋返回波斯，但意外受伤，因伤而死。他死后，两位兄弟（是\OriginalTerm{贤士}{magi}，\OriginalTerm{玛待}{Mede}族人）统治波斯七个月。为了除掉他们，七位最高贵的波斯人结成了阴谋，为首的是\Person{希斯塔斯普}的儿子\Person{达理阿}，他是\Person{居鲁士}的堂兄弟，众人一致同意将王国授予他：他统治了三十六年。他在死前四年，参加了\Place{马拉松}战役，这场战役在希腊和罗马历史中都极为著名。那件事大约发生在罗马建城后第二百六十年，\Person{马克里努斯}和\Person{奥古里努斯}担任执政官时，即八百八十八年前，只要我关于罗马执政官继承顺序的研究没有欺骗我；因为我已自始至终计算到\Person{斯提利科}时代。\Person{达理阿}之后是\Person{薛西斯}，据说他统治了二十一年，尽管我发现大多数版本将其统治期记载为二十五年。继承他的是我们上文提及的\Person{阿尔塔薛西斯}。由于他下令停止犹太城市和圣殿的建造，工程暂停到达理阿王第二年。但为了使年代的顺序完整地延伸到达理阿，我有必要说明：\Person{阿尔塔薛西斯}在位四十一年，\Person{薛西斯}两个月，在他之后\Person{苏克狄安努斯}统治了七个月。\par

% structure: p0020 source=h2 target=subsection reason=relative-heading-rank; base=h2

\subsection{第十章}

接下来，\Person{达理阿}（当时名为\Person{奥库斯}）获得了王国，在他的统治下圣殿得以重建。他有三名忠诚可靠的希伯来人充当他的护卫，其中一人因其所显示的智慧而赢得了王的钦佩。于是，这人被给予机会，可以要求任何他心中向往的东西。他因其故土的废墟而叹息，请求准许重建城市，并从王那里获得命令，催促总督和统治者加紧建造圣殿，并提供所需费用。因此，圣殿在四年内完工；即\Person{达理阿}开始统治后的第六年，那时对犹太人民来说似乎已经足够了。因为重建城市是一项极其艰巨的工作，他们对自己资源缺乏信心，当时不敢着手如此困难的事业，而满足于重建了圣殿。同时，精通法律的经师\Person{厄斯德拉}，在圣殿完工约二十年后（\Person{达理阿}此时已去世，他在位十九年），蒙\Person{阿尔塔薛西斯二世}（不是那位介于两位薛西斯之间的，而是继承\Person{达理阿·奥库斯}的那位）许可，从\Place{巴比伦}出发，有许多人跟随他，他们携带了各种工艺的器皿以及王为天主的圣殿所送的礼物到\Place{耶路撒冷}。与他们同行的只有十二位\OriginalTerm{肋未人}{Levites}；因为据说那时该支派的人难得找到这个数目。他发现犹太人与\OriginalTerm{外邦人}{Gentiles}结婚，便严厉责备他们，命令他们断绝一切此类关系，并赶走这些婚姻所生的子女；所有的人都服从了他的话。于是，人民被圣化后，举行了古代法律所规定的礼仪。但我没有发现\Person{厄斯德拉}为重建城市做了任何事；因为依我看，他认为更紧迫的责任是改革人民，使他们摆脱所沾染的腐化习俗。\par

% structure: p0022 source=h2 target=subsection reason=relative-heading-rank; base=h2

\subsection{第十一章}

当时在巴比伦有一个名叫乃赫米雅的人，是国王的仆人，生为犹太人，因他所立下的功绩而深受阿尔塔薛西斯的宠爱。他询问同族犹太人关于他们祖先城市的情况，得知故土仍如先前一样处于沦陷状态，据说他忧心如焚，向天主祈祷，叹息流泪不已。他也想起本民族的罪过，恳切祈求天主的怜悯。因此，国王注意到他在侍宴时比往常更加忧愁，便询问他忧伤的原因。于是他开始哀叹民族的灾祸和祖先城市的废墟，该城如今已近两百五十年夷为平地，成为所受苦难的见证和敌人的笑柄。因此他请求国王准许他前去重建。国王应允了这一孝心的恳求，立即派他带着骑兵卫队出发，以便更安全地完成旅程，同时赐予他致各地首领的书信，要求他们提供一切所需。他抵达耶路撒冷后，将城建工作按人分配给了民众；众人竞相执行所接受的命令。重建工程已进行到一半时，周围外邦人的嫉妒爆发了，邻近的城市合谋阻挠工程，吓阻犹太人建造。但乃赫米雅部署了守卫来抵御那些袭击民众的人，毫不惊慌，继续推进已开始的工作。这样，城墙建成，城门入口完工后，他便为城内各家族划分地段建造房屋。他也估算了人口与城市规模不相称：男女老少总数不超过五万人——从前庞大的人数因频繁的战争和大批被掳而锐减至此。因为古时那两个支派（余民全属他们）在十支派分离后，尚能提供三十二万武装人员。但因罪过，天主将他们交付死亡和掳掠，他们便衰落到如今这般可怜的小数目。然而，这群体，如我所说，仅由两个支派组成：先前被掳的十支派散居在帕提亚人、玛待人、印度人和厄提约丕雅人中间，从未返回故土，至今仍受野蛮民族统治。被重建城市的完工记载于阿尔塔薛西斯在位第三十二年。从那时起到基督被钉十字架，即傅菲乌斯·革米努斯和鲁贝利乌斯任执政官之时，其间过去了三百九十八年。但从圣殿重建到它在韦斯帕西安治下被提托摧毁（时值奥古斯都任执政官），其间有四百八十三年。这昔日由达尼尔预言，他宣告从圣殿重建到其倾覆将有七十九周。从犹太人被掳到城市重建，其间有二百六十年。\par

% structure: p0024 source=h2 target=subsection reason=relative-heading-rank; base=h2

\subsection{第十二章}

在这一时期，我们认为艾斯德尔和友弟德生活过，但我承认我难以确定她们的事迹应特别与哪些国王相联系。因为艾斯德尔据说生活在阿尔塔薛西斯王时代，但我发现波斯有两位同名的国王，因此很难断定她属于哪一位。然而，我认为最好将艾斯德尔的事迹与耶路撒冷重建时期的阿尔塔薛西斯联系起来，因为如果她生活在先前的阿尔塔薛西斯时代（厄斯德拉记述过他的时代），厄斯德拉不可能对如此杰出的女子只字不提。这一点尤其令人信服，因为我们知道圣殿的建造曾被那位阿尔塔薛西斯禁止，而如果艾斯德尔当时已与他成婚，她必不会允许此事。现在我将叙述她所成就的事。当时有一位名叫瓦市提的女子与国王成婚，容貌惊人。国王常向众人夸耀她的美丽，一日他设宴款待，命王后前来展示她的美貌。但她比愚昧的国王更为明智，且过于谦逊，不愿在男人面前裸露身体，便拒绝服从命令。国王因这侮辱而勃然大怒，将她逐出婚姻和王宫。于是，当寻求一位年轻女子取代她成为国王的妻子时，艾斯德尔被发现美貌超群。她是一位犹太女子，属本雅明支派，是孤儿，无父无母，由堂兄摩尔德开抚养长大。与她成婚后，她依照抚育她的人的指示，隐瞒了自己的民族和祖国，并受他告诫不要忘记祖先的传统，也不可，尽管身为俘虏与外邦人结婚，却参与外邦人的食物。因此，她与国王结合后，不久，如所预料，轻易藉其美貌的力量完全俘获了国王的心，以致国王以王权的象征使她与自己平起平坐，赐予她紫红袍。\par

% structure: p0026 source=h2 target=subsection reason=relative-heading-rank; base=h2

\subsection{第十三章}

此时，\Person{摩爾德開}已位居王最亲近者之列，全权管理王室内务。他曾向王揭发两名太监所设的阴谋，因而更加受宠，被授予最高荣誉。当时有一位\Person{哈曼}，是王极为亲信的朋友，王使其与自己平等，并依照君主的惯例，命人向他敬拜。惟独\Person{摩爾德開}一人拒绝如此行，这大大激起了\Person{哈曼}的怒气，使之指向自己。于是，\Person{哈曼}立意要毁灭那希伯来人，便到王面前声称，在他的王国里有一族人，信仰邪恶的迷信，为天主与人所共憎。他说，这些人按外邦法律生活，理当被消灭；将整个民族交付死亡，乃正义之举。同时，他答应王从他们的财物中获取巨额财富。这位野蛮的君主轻易被说服，颁令灭绝犹太人，并立即派人将谕令传达全国，从\Place{印度}直到\Place{埃塞俄比亚}。\Person{摩爾德開}闻此，撕裂衣服，披上麻衣，将灰撒在头上，来到王宫，在那里使整个地方回荡着他的哀哭与怨诉，大声呼喊说，\Quote{一个无辜的民族竟要灭亡，且无任何毁灭的理由，这实属不当。}\Person{艾斯德尔}被哀哭之声吸引，得悉了实情。但当时她不知所措（因按波斯风俗，王后不得随意觐见君王，除非被召，而且不可在王愿见的任何时间进入，只在定期方可）；恰好此时，按此规定，\Person{艾斯德尔}有三十天不得接近王前。然而，她认为自己应为同胞冒险，即使注定灭亡，也准备为如此崇高的事业赴死。她呼求了天主之后，便进入王的内廷。那野蛮人起初对这非常之举感到惊异，但逐渐被女性的柔媚所打动，最终竟陪同王后赴她所设的宴席。与他同去的还有王的宠臣\Person{哈曼}，但他却是犹太民族死敌。宴饮过后，酒酣耳热之际，\Person{艾斯德尔}俯伏在王膝前，恳求他阻止那威胁她民族的毁灭。王许诺，若她再有所求，必不拒绝。\Person{艾斯德尔}立刻抓住机会，要求处死\Person{哈曼}，作为对她民族的赔补，因\Person{哈曼}曾想摧毁他们。但王不能忘记他的朋友，犹豫片刻，便退出一会儿考虑此事。他返回时，见\Person{哈曼}正抱住王后的膝，便勃然大怒，高喊有人在对王后施暴，遂下令将\Person{哈曼}处死。此时王得知\Person{哈曼}已预备了十字架，要叫\Person{摩爾德開}受刑。于是，\Person{哈曼}被悬在那十字架上，他所有家产都归了\Person{摩爾德開}，而全体犹太人均获释放。\Person{阿塔薛西斯}在位六十二年，由\Person{奥库斯}继位。\par

% structure: p0028 source=h2 target=subsection reason=relative-heading-rank; base=h2

\subsection{第十四章}

在这些事件之后，我理应附上友弟德事迹的记载；因为相传她生活在充军之后，但圣史并未揭示在她那个时代谁是波斯王。然而，圣史将她事迹发生时的那位王称为\Person{拿步高}，而那位拿步高显然不是攻取耶路撒冷的那一位。但我发现充军之后，在波斯人中并没有任何一位名叫拿步高的君王；除非是因为某位王所表现出的愤怒和类似的作为，犹太人便将任何如此行事的王都称作\Person{拿步高}。不过，多数人认为那是\Person{居鲁士}之子\Person{刚比西斯}，理由是：他作为征服者，曾深入\Place{埃及}和\Place{埃塞俄比亚}。但圣史与这种意见相悖；因为记载说，友弟德生活在所论之王的第十二年。然而，刚比西斯执掌最高权力不过八年。因此，如果允许就历史问题作一推测，我倾向于相信她的事迹发生在\Person{奥库斯}王统治时期，他是继第二位\Person{阿尔塔薛西斯}之后登基的。我如此推测是基于这样一个事实：（正如我在世俗历史中所读）他生性残忍，好战。因为他既与邻国作战，又通过战争收复了多年前叛离的埃及。同时，据说他还嘲笑埃及人的神圣礼仪以及被他们奉为神的\Person{阿匹斯}；此事激怒了他的一名太监\Person{巴瓜斯}，此人是埃及人，愤慨于王的行为，后来以杀死王来报仇，认为王侮辱了他所属的民族。如今，默感的历史提到了这位\Person{巴瓜斯}；因为当\Person{敖罗斐乃}奉王命率领军队攻打犹太人时，记载说巴瓜斯也在军中。因此，我拿出这一点来证明我所表达的意见，即那位名为拿步高的王其实就是\Person{奥库斯}，并非没有道理，因为世俗历史学家记载巴瓜斯生活在他的统治时期。但是，任何人都不应感到奇怪，即纯粹的世俗作家并未触及圣经所记载的任何一个要点。天主之神如此眷顾，使圣史严格限于自身的奥秘，不受任何腐败之口或混杂真伪之言的玷污。事实上，圣史与世俗的事务分离，且只能用神圣的语言来表达，显然不应与其他历史等量齐观地混杂在一起。因为若将这部历史与涉及其他事物或追求不同探询的其他历史混在一起，是极不合适的。但现在，我将继续叙述其余部分，并尽可能简短地讲述友弟德所行的事迹。\par

% structure: p0030 source=h2 target=subsection reason=relative-heading-rank; base=h2

\subsection{第十五章}

犹太人，如上所述返回故土后，他们的处境和城邑尚未完全安定，波斯王便向玛待开战，与他们的国王——名叫阿法撒得——交战并获胜。那王被击杀后，他将该民族并入自己的帝国。他又派他任命的统帅敖罗斐乃率领十二万步兵和一万二千骑兵先行出发，对其他民族也如法炮制。敖罗斐乃在蹂躏了基里基雅和阿拉伯之后，用武力攻占了许多城邑，或用恐吓迫使其投降。如今大军已推进到大马士革，使犹太人极度恐惧。但他们既无力抵抗，同时心中又无法接受投降的念头——因为他们先前已从经历中知道了为奴的苦楚——于是成群结队前往圣殿。在那里，他们普遍哀叹，混合着哭泣，恳求天主的援助，说他们已因自己的罪愆受到天主足够的惩罚，请求他宽恕他们这些刚从奴役中被解救出来的遗民。与此同时，敖罗斐乃已接受摩阿布人投降，并将他们收为盟友，一同攻打犹太人。他向他们的首领询问，希伯来人倚靠什么力量，以致他们心中不肯接受投降的念头。一个名叫阿希约尔的人向他陈述了实情，即：犹太人崇拜天主，并由祖先教养敬虔的礼仪，他们从前曾在埃及为奴一段时间，后来蒙天主援助离开那地，并在干涸的海底步行渡过，最终打败了所有敌对的民族，收复了他们祖先居住的领土。此后，他们的事务有各种起伏，或顺利或相反；当他们陷入逆境时，又再次摆脱苦难，发现天主按他们的功过，时而向他们发怒，时而与他们和好，因此当他们犯罪时，就受到敌人攻击或被掳为奴的惩罚，但一旦享有天主的眷顾，他们总是不可征服。所以，如果他们现今无罪，就不可能被征服；但如果情况相反，他们就会轻易被征服。听到这番话，敖罗斐乃因屡战屡胜而得意洋洋，以为一切都得向他屈服，如今却因胜利被说成主要取决于犹太人的罪恶而勃然大怒，下令把阿希约尔推到希伯来人的营中去，让他与他所断言不可征服的人一同灭亡。当时犹太人已退守山地；奉命执行的人来到山脚下，将阿希约尔锁在那里。犹太人察觉后，替他解开锁链，带上山去。他们问明事情缘由，他便向他们解释；他们便平平安安地接待他，等候结果。我还要补充，战后他受了割礼，成了犹太人。好吧，敖罗斐乃看到地势险要，无法到达高处，便用兵士包围群山，并竭尽全力切断希伯来人的所有水源。为此，他们更早地感受到围困的苦楚。因缺水而困乏，他们成群结队地来到他们的首领敖齐雅面前，都倾向于投降。但他回答说，应再等候片刻，期待天主的援助，因此将投降的日期定在第五天之后。\par

% structure: p0032 source=h2 target=subsection reason=relative-heading-rank; base=h2

\subsection{第十六章}

当此事传到\Person{友弟德}（一位极其富有、以美貌著称、但更以德行胜于美貌的寡妇，当时正在营中）耳中时，她认为在同胞的困境中，她应当挺身而出，即使可能招致自身毁灭。于是她妆扮头面，美化容貌，仅带一名侍女，进入敌营。她立即被带到\Person{敖罗斐乃}面前，告诉他说，她同胞的处境已绝望，因此她已为自己的性命逃命作了准备。然后她向将军请求夜间自由出营祈祷的权利。这命令随即传给了哨兵和守门者。但当她通过三天的实践养成了进出营地的习惯，并以此赢得了蛮族人的信任后，\Person{敖罗斐乃}产生了玷污这俘虏的欲望；因为\Person{友弟德}天姿国色，轻易打动了这位波斯人。于是，她被太监\Person{巴哥阿}带到了将军的帐幕；宴会开始，蛮族人饮酒过量，醉得不省人事。随后，在仆人退下后，他未及侵犯那妇人便睡着了。\Person{友弟德}抓住机会，斩下敌人的头颅，带了出去。她被视为照常出营，平安地回到自己同胞那里。次日，希伯来人从高处将\Person{敖罗斐乃}的头颅示众；然后出击，攻向敌营。这时蛮族人聚集在将军的帐幕周围，等待作战信号。当他们发现他被肢解的尸体时，在可耻的惊慌中转身逃窜，在敌人面前溃逃。犹太人则追击逃敌，屠杀数千人后，夺取了营寨和其中的战利品。\Person{友弟德}受到最高的赞扬，据说活了一百零五岁。如果这些事正如我们所信，发生在\Person{奥古斯}王统治的第十二年，那么从\Place{耶路撒冷}重建到那次战争，相隔了二十二年。\Person{奥古斯}总共在位二十三年。他比所有人都残忍，性情更为野蛮。太监\Person{巴哥阿}在他患病时用毒药将他害死。此后，他的儿子\Person{阿尔塞斯}统治了三年，\Person{大流士}统治了四年。\par

% structure: p0034 source=h2 target=subsection reason=relative-heading-rank; base=h2

\subsection{第十七章。}

\Place{马其顿}的\Person{亚历山大}向他开战。他被击败后，王权从波斯人手中被夺走，从\Person{居鲁士}建国算起，已持续了二百五十年。据说征服了几乎所有民族的\Person{亚历山大}访问了\Place{耶路撒冷}的圣殿，并献上礼物；他向他所征服的全境宣告，居住在那里的犹太人可自由返回本国。他在位第十二年末，以及征服\Person{大流士}后七年，死于\Place{巴比伦}。与他一同进行这些重要战争的朋友们将他的帝国瓜分。起初，他们管理所担任的职务，不使用王的称号，而由\Person{亚历山大}的兄弟\Person{阿里德乌斯·腓力}为王，此人性格极其软弱，王权只是名义上和表面上的授予，但实际权力掌握在那些瓜分了军队和行省的人手中。确实，这种状况并未持续很久，所有人都宁愿被称为国王。在\Place{叙利亚}，\Person{塞琉古}是\Person{亚历山大}之后的第一位国王，\Place{波斯}和\Place{巴比伦}也受他统治。那时，犹太人每年向国王缴纳三百塔连得银子的贡金；但他们不是由外邦官员治理，而是由自己的司祭治理。他们按照祖先的习俗生活，直到许多人在长期的和平中再次腐化，开始以叛乱搅乱一切，制造骚乱，同时在情欲、贪婪和权力欲望的影响下觊觎\OriginalTerm{大司祭职}{high-priesthood}。\par

% structure: p0036 source=h2 target=subsection reason=relative-heading-rank; base=h2

\subsection{第十八章。}

首先，在\OriginalTerm{伟大的安提约古}{Antiochus the great}之子\OriginalTerm{色娄苛}{Seleucus}王治下，一个名叫\OriginalTerm{息孟}{Simon}的人向国王诬告了司祭\OriginalTerm{敖尼雅}{Onias}——一位圣洁无瑕的人，试图推翻他，但未成功。然后，过了一段时间，\OriginalTerm{敖尼雅}{Onias}的兄弟\OriginalTerm{雅松}{Jason}前往继承其兄\OriginalTerm{色娄苛}{Seleucus}王位的\OriginalTerm{安提约古}{Antiochus}王那里，许诺增加贡金，条件是将\OriginalTerm{大司祭职}{high-priesthood}转交给他。虽然一个人年复一年地享有大司祭职是不寻常的，而且直到现在仍是不允许的，但国王贪婪之心很容易被说服。于是，\OriginalTerm{敖尼雅}{Onias}被赶下职位，司祭职赐给了\OriginalTerm{雅松}{Jason}。他以极其可耻的方式折磨自己的同胞和国家。然后，当他通过一个名叫\OriginalTerm{默乃劳}{Menelaus}（就是前面提到的那个\OriginalTerm{息孟}{Simon}的兄弟）的人送去了他承诺给国王的钱财时，一条通向野心的道路被打开，\OriginalTerm{默乃劳}{Menelaus}用\OriginalTerm{雅松}{Jason}之前所用的同样手段获得了司祭职。但不久之后，由于他没有提供承诺的金额，他被赶下职位，\OriginalTerm{里息玛苛}{Lysimachus}被取代。于是\OriginalTerm{雅松}{Jason}和\OriginalTerm{默乃劳}{Menelaus}之间发生了可耻的冲突，直到\OriginalTerm{雅松}{Jason}流亡离国。通过这样的例子，人民的道德败坏到如此程度，以至于许多本地人请求\OriginalTerm{安提约古}{Antiochus}允许他们按外邦人的方式生活。当国王批准他们的请求时，所有最卑劣的人都竞相建造庙宇，向偶像献祭，亵渎\OriginalTerm{天主的法律}{law of God}。与此同时，\OriginalTerm{安提约古}{Antiochus}从\OriginalTerm{亚历山大里亚}{Alexandria}返回（他当时正在与埃及王作战，但在\OriginalTerm{保禄}{Paulus}和\OriginalTerm{克拉苏}{Crassus}担任执政官时，因元老院和罗马人民的命令而放弃了），并去了\Place{耶路撒冷}。他发现人民因采纳各种迷信而纷争不已，便摧毁了\OriginalTerm{天主的法律}{law of God}，偏袒那些行不虔之径的人，同时搬走了圣殿的所有装饰品，大肆破坏。这事发生在\OriginalTerm{亚历山大}{Alexander}死后一百五十年，\OriginalTerm{保禄}{Paulus}和\OriginalTerm{克拉苏}{Crassus}如前所述担任执政官，约在\OriginalTerm{安提约古}{Antiochus}开始统治后五年。\par

% structure: p0038 source=h2 target=subsection reason=relative-heading-rank; base=h2

\subsection{第十九章。}

但为了正确保存日期的顺序，并且更清楚地表明这安提约古是谁，我们要列举\OriginalTerm{亚历山大}{Alexander}之后\OriginalTerm{叙利亚}{Syria}诸王的名字和年代。那么，\OriginalTerm{亚历山大}{Alexander}王去世后，正如我们上面所述，他的整个帝国被他的朋友们瓜分，并在一段时间内以国王的名义由他们统治。\OriginalTerm{色娄苛}{Seleucus}在九年后自称为\OriginalTerm{叙利亚}{Syria}王，统治了三十二年。之后是他的儿子\OriginalTerm{安提约古}{Antiochus}，在位二十一年。然后是\OriginalTerm{安提约古}{Antiochus}之子\OriginalTerm{安提约古}{Antiochus}，被称为\OriginalTerm{特乌斯}{Theus}，在位十五年。之后是他的儿子\OriginalTerm{色娄苛}{Seleucus}，被称为\OriginalTerm{卡里尼苛}{Callinicus}，在位二十一年。另一个\OriginalTerm{色娄苛}{Seleucus}，\OriginalTerm{卡里尼苛}{Callinicus}之子，在位三年。他死后，\OriginalTerm{卡里尼苛}{Callinicus}的兄弟\OriginalTerm{安提约古}{Antiochus}统治\OriginalTerm{亚细亚}{Asia}和\OriginalTerm{叙利亚}{Syria}三十七年。这就是\OriginalTerm{路基乌斯·斯奇比欧·亚细亚提库斯}{Lucius Scipio Asiaticus}与之作战的那个\OriginalTerm{安提约古}{Antiochus}；他在战争中失败，被剥夺了部分帝国。他有两个儿子，\OriginalTerm{色娄苛}{Seleucus}和\OriginalTerm{安提约古}{Antiochus}，后者被交给罗马人作人质。于是，\OriginalTerm{伟大的安提约古}{Antiochus the great}死后，他的幼子\OriginalTerm{色娄苛}{Seleucus}获得了王国，在他治下，如前所述，司祭\OriginalTerm{敖尼雅}{Onias}被\OriginalTerm{息孟}{Simon}控告。然后\OriginalTerm{安提约古}{Antiochus}被罗马人释放，代替他作为人质的是当时在位的\OriginalTerm{色娄苛}{Seleucus}之子\OriginalTerm{德默特琉}{Demetrius}。\OriginalTerm{色娄苛}{Seleucus}在统治十二年时去世，他的兄弟\OriginalTerm{安提约古}{Antiochus}，曾为罗马人质，夺取了王国。他在统治开始后五年，正如我们上面所显示的，蹂躏了\Place{耶路撒冷}。因为他必须向罗马人缴纳沉重的贡金，几乎被迫为了应付这笔巨额开支，通过抢劫来筹集资金，并不错过任何掠夺的机会。然后，两年后，\OriginalTerm{犹太人}{Jews}再次遭受与先前类似的灾难，为了防止他们因众多苦难而发动战争，他在堡寨中安置了驻军。接着，为了推翻神圣的法律，他颁布法令，要求所有人抛弃祖先的传统，按外邦人的方式生活。并且不乏有人乐意服从这亵渎的法令。于是出现了可怕的景象：在各城市，街道上公开献祭，而\OriginalTerm{法律}{law}和\OriginalTerm{先知}{prophets}的圣洁书卷被烈火吞噬。\par

% structure: p0040 source=h2 target=subsection reason=relative-heading-rank; base=h2

\subsection{第二十章。}

那时，\OriginalTerm{若望}{John}之子\OriginalTerm{玛塔提雅}{Matthathias}是\OriginalTerm{大司祭}{high-priest}。当他被王的仆人强迫服从法令时，他以非凡的勇气藐视亵渎的法令，并在众人面前杀死了一个公开行亵渎之事的\OriginalTerm{希伯来人}{Hebrew}。这样找到了一位领袖，叛乱立刻发生。\OriginalTerm{玛塔提雅}{Matthathias}离开城镇；许多人聚集到他那里，他组织起一支正规军队。那支队伍中每个人的目标是用武器保卫自己，抵抗亵渎的政府，宁愿战死也不参与不虔的仪式。与此同时，\OriginalTerm{安提约古}{Antiochus}强迫其辖下希腊城市中的\OriginalTerm{犹太人}{Jews}献祭，并以闻所未闻的酷刑折磨拒绝者。此时，发生了那著名而惊人的七兄弟及其母亲的苦难。所有兄弟，当他们被迫违反\OriginalTerm{天主的法律}{law of God}和\OriginalTerm{祖先的习俗}{customs of ancestors}时，宁愿选择死亡。最后，他们的母亲也陪伴他们一同受苦和死亡。\par

% structure: p0042 source=h2 target=subsection reason=relative-heading-rank; base=h2

\subsection{第二十一章。}

在此期间，\Person{玛塔提雅}去世了，他指定自己的儿子\Person{犹大}接替他的位置，作为他所召集的军队的统帅。在他的领导下，与王室军队进行了几次成功的战斗。首先，他摧毁了敌军将领\Person{阿波罗尼}及其全军，此人曾率领大量军队投入战斗。当一位名叫\Person{色龙}的人——当时是\Place{叙利亚}的统治者——听说此事后，他增派了兵力，并因其人数众多而气势汹汹地进攻\Person{犹大}，但战斗打响后，他被击溃并逃窜；损失近八百人后，他返回了\Place{叙利亚}。\Person{安提约基雅}得知此事后，充满愤怒和懊悔，因为他的将领们尽管拥有庞大的军队却被打败，这使他恼怒。于是他召集整个帝国的援军，并向士兵发放饷银，几乎耗尽国库。因为那时他正特别遭受金钱匮乏之苦。原因一方面是犹太人在他统治下每年惯常向他缴纳三百塔冷通银的贡赋，现在却对他叛乱；另一方面，许多希腊城邦和国家因迫害之祸而动荡不安。因为\Person{安提约基雅}甚至没有放过外邦人，他试图劝服他们放弃长期以来持有的迷信，并使他们改宗一种宗教仪式。无疑，那些视一切为不神圣的人很容易被诱导放弃他们古老的崇拜方式，但同时所有人都处于惊慌和灾难之中。因此，税收已停止缴纳。出于这些原因而怒火中烧（因为他，昔日最富有的君王，如今深深感受到因自己邪恶而导致的贫穷），他将军队分给\Person{里息雅斯}，并把\Place{叙利亚}和对抗犹太人的战争托付给他，而他自己则出发前往波斯人那里，向他们征收赋税。于是，\Person{里息雅斯}挑选了\Person{托勒米}、\Person{戈尔基雅}、\Person{多罗}和\Person{尼加诺尔}作为战争将领；并给了他们四万步兵和七千骑兵。起初的进攻使犹太人大为惊慌。然后，当所有人都绝望时，\Person{犹大}鼓励他的战士们勇敢地投入战斗——如果他们信赖天主，一切都会在他们面前退让；因为常常在此之前，少数人对抗多数人却赢得了胜利。宣布了斋戒，举行了祭献，之后他们便下到战场。结果是敌军溃散，\Person{犹大}夺取了他们的营寨，在里面发现了大量金子和\Place{提洛}的财宝。因为来自\Place{叙利亚}的商人，对胜利确信无疑，曾跟随国王军队，希望购买俘虏，现在他们自己却被掠夺了。当使者将这些事报告给\Person{里息雅斯}时，他更加努力地集结军队，一年后再次以庞大的军队进攻犹太人；但被打败后，他撤退到了\Place{安提约基雅}。\par

% structure: p0044 source=h2 target=subsection reason=relative-heading-rank; base=h2

\subsection{第二十二章。}

\Person{犹大}击溃敌人后，返回\Place{耶路撒冷}，一心致力于洁净和重建圣殿，圣殿曾被\Person{安提约基雅}推翻，被外邦人亵渎，呈现出一片凄凉的景象。但由于叙利亚人占据了堡垒，它与圣殿相连，位置却高于圣殿，确实坚不可摧，低处便无法接近，因为频繁从上而下的出击阻止了人们靠近。但\Person{犹大}派出了一支非常强大的队伍来对抗这些袭击者。这样，圣殿的工程得到了保护，圣殿被城墙围起，并指派了武装人员维持永久防御。\Person{里息雅斯}再次率领增援部队回到\Place{犹太}，又一次遭到惨败，损失了自己的军队和援军，这些援军是由各邦派来与他联合作战的。与此同时，\Person{安提约基雅}，正如我们上面所说，已进军\Place{波斯}，试图劫掠该国最富有的城镇\Place{厄里曼}和一座位于那里的充满黄金的圣殿；但由于从四面八方聚集了大量人群保卫该地，他被击退了。此外，他收到了\Person{里息雅斯}作战不利的消息。于是，由于精神痛苦，他患上了身体疾病。但当他内心受折磨时，他想起自己对天主子民施加的苦难，并承认这些灾祸是应得的。然后，几天后，他死了，在位十一年。他把王国留给了他的儿子\Person{安提约基雅}，并赐予他\Person{欧帕托尔}的名字。\par

% structure: p0046 source=h2 target=subsection reason=relative-heading-rank; base=h2

\subsection{第二十三章。}

那时，\Person{犹大}围攻驻扎在堡垒里的叙利亚人。他们因饥荒和缺乏一切而陷入困境，便派使者到国王那里请求援助。于是，\Person{欧帕托尔}率领十万步兵和两万骑兵前来援助，大象在他的阵线前行进，使旁观者极其恐惧。于是犹大放弃围攻，前去迎击国王，并在第一次战斗中击溃了叙利亚人。国王请求和平，但由于他生性背信弃义，不善用和平，背信弃义之后便遭到了报应。因为\Person{色娄苛}的儿子\Person{德默特琉}——我们上面说过他被交给罗马人作为人质——听说了\Person{安提约古}去世的消息，便请求罗马人派他去夺取王位。当这一请求被拒绝后，他秘密地逃离\Place{罗马}，来到\Place{叙利亚}，杀死了安提约古的儿子——后者统治了一年零六个月——夺取了最高权力。正是在他统治期间，犹太人首次请求与罗马人民建立友谊和同盟；这一使团受到友好接待，根据元老院的法令，他们被称为同盟者和朋友。与此同时，德默特琉通过他的将领们对犹大作战。起初，军队由一位名叫\Person{巴基德}的人和一位名叫\Person{阿耳基慕}的犹太人率领；后来被任命为统帅的\Person{尼加诺尔}在战斗中阵亡。然后巴基德和阿耳基慕恢复了权力，扩充了兵力，与犹大作战。叙利亚人在那次战斗中获胜，残酷地滥用他们的胜利。希伯来人选出犹大的兄弟\Person{约纳堂}来接替他。与此同时，阿耳基慕在可怕地蹂躏\Place{耶路撒冷}之后死了；巴基德失去了盟友，便回到国王那里。两年后，巴基德再次向犹太人开战，战败后请求和平。和平在他接受某些条件下被赐予，即他应交出逃兵和俘虏，连同他在战争中夺取的一切。\par

% structure: p0048 source=h2 target=subsection reason=relative-heading-rank; base=h2

\subsection{第二十四章}

当这些事在\Place{犹太}发生时，一个在\Place{洛多斯}受过教育的年轻人，名叫\Person{亚历山大}，自称是\Person{安提约古}的儿子（这是假的），并在\Place{亚历山大里亚}国王\Person{仆托肋米}的力量帮助下，率领一支军队进入\Place{叙利亚}。他在战争中击败了\Person{德默特琉}，并在德默特琉统治十二年后杀死了他。这位亚历山大在向德默特琉开战之前，与\Person{约纳堂}结盟，并赠与他一件紫袍和王室标志。为此，约纳堂曾以辅助部队协助他；在德默特琉失败后，约纳堂最先向他表示祝贺。后来亚历山大也没有违背他所承诺的誓言。因此，在他掌权的五年中，犹太人的事务是和平的。在这种情况下，德默特琉的儿子德默特琉——他在父亲死后逃往\Place{克里特}，在克里特人将领\Person{拉斯德乃}的煽动下——试图通过战争恢复他父亲的王国，但发现自己的力量不足以胜任，便向埃及国王、亚历山大的岳父\Person{仆托肋米非洛默托尔}请求援助，而当时非洛默托尔与他的女婿关系不和。但是，与其说他是被恳求者的哀求所打动，不如说是被夺取叙利亚的希望所驱使，他与德默特琉联合，并将自己已经嫁给亚历山大的女儿许配给他。亚历山大与这两人进行了会战。仆托肋米在战斗中阵亡，但亚历山大被击败；不久后他被杀，他统治了五年，或者如我在许多作者那里所看到的，九年。\par

% structure: p0050 source=h2 target=subsection reason=relative-heading-rank; base=h2

\subsection{第二十五章}

\Person{德默特琉}就这样得到王国后，善待\Person{约纳堂}，与他缔结条约，并恢复犹太人遵守自己的法律。与此同时，曾属于亚历山大党派的\Person{特黎丰}被任命为\Place{叙利亚}总督，以通过战争牵制他。而约纳堂则率军四万，威震一方，下到战场。特黎丰见自己无法匹敌，便伪装渴望和平，并杀死了被接待并邀请与之结盟的\Person{仆托肋买}。约纳堂之后，最高权力被授予他的兄弟\Person{西满}。他隆重地举行了兄弟的葬礼，并建造了那些著名的七座金字塔，工艺极为精美，将兄弟和父亲的遗骸埋葬其中。然后德默特琉更新了与犹太人的条约；考虑到特黎丰给他们造成的损失（因为约纳堂死后，特黎丰通过战争蹂躏了他们的城市和领土），德默特琉永远免除了他们的年贡；因为直到那时，他们一直向叙利亚国王进贡，除非他们武力抵抗。这件事发生在德默特琉王二年；我们之所以提到这一点，是因为直到这一年，我们已经按顺序叙述了亚细亚诸王的时代，以使年代的序列清晰明了。但现在我们将按照事件发生的顺序，叙述犹太人中大司祭或君王的时代，直到基督诞生的时期。\par

% structure: p0052 source=h2 target=subsection reason=relative-heading-rank; base=h2

\subsection{第二十六章}

那麼，在那之后，约纳堂的兄弟西满，如上所述，以大司祭的权柄统治了希伯来人。因为那荣誉既由其同胞，又由罗马人民授予了他。他在德默特琉王第二年开始了对同胞的统治，但八年后，被托勒密的计谋所骗而丧命。他的儿子若翰继位。他因与极强大的民族希尔卡尼人作战有功，而得到希尔卡尼的称号。他统治了二十六年，拥有最高权力。他死后，亚里士多布鲁受任大司祭，是充军之后第一个称王并头戴王冠的人。一年后他去世了。然后他的儿子亚历山大，兼为国王和大司祭，在位二十七年；但在他所作所为中，除了残忍，我没有发现值得一提的事。他留下两个幼子，名叫亚里士多布鲁和希尔卡诺，其妻萨利纳或亚历山大，统治了三年。他死后，两兄弟之间爆发了关于最高权力的可怕冲突。起初，希尔卡诺执掌政权；但后来被他弟弟亚里士多布鲁击败，逃到庞培那里。那位罗马将军，结束了与米特拉达梯的战争，安定了亚美尼亚和本都，实际上是他所到之处的所有民族的征服者，他想要向内陆进军，将附近所有地区并入罗马帝国。因此他调查了战争的原因和取胜的方法。于是他欣然接纳了希尔卡诺，并在其引导下攻击了犹太人；但当城市被攻陷和摧毁时，他饶恕了圣殿。他将亚里士多布鲁戴枷锁送往罗马，将大司祭权归还给希尔卡诺。确定了犹太人应缴纳的贡物后，他派了一个名叫安提帕特的阿什凯隆人做他们的总督。希尔卡诺掌握大权三十四年；但在与帕提亚人作战时被俘。\par

% structure: p0054 source=h2 target=subsection reason=relative-heading-rank; base=h2

\subsection{第二十七章}

然后，一个外邦人，即阿什凯隆人安提帕特之子黑落德，向罗马元老院和人民请求并获得了犹太地的统治权。在他治下，犹太人第一次有了外邦人为王。因为基督的来临已近，按照先知们的预言，他们必须被剥夺自己的统治者，免得他们寻求基督以外的任何事物。在此黑落德统治的第三十三年，基督于萨比努斯和鲁菲努斯执政官任期的十二月二十五日诞生。但我们不敢触及那些记载在\WorkTitle{福音书}中以及后来在\WorkTitle{宗徒大事录}中的事，以免我们紧凑著作的特性在任何程度上减损这些事件的神圣尊严；我将继续叙述余下的事情。主诞生后，黑落德又统治了四年；因为他的整个统治时期共三十七年。之后，分封侯阿尔赫劳来了，统治八年，黑落德统治二十四年。在他治下，即他统治的第十八年，主被钉十字架，当时富菲乌斯·格米努斯和鲁贝利乌斯·格米努斯任执政官；从那时起直到斯提利科任执政官，已过了三百七十二年。\par

% structure: p0056 source=h2 target=subsection reason=relative-heading-rank; base=h2

\subsection{第二十八章}

路加记载了宗徒们的事迹，直到保禄在尼禄皇帝治下被带到罗马。至于尼禄，我不说他是最坏的君王，而是说他堪当被称为所有人甚至野兽中最卑劣的。他是第一个开始迫害的人；而且我不确定他是否也会是最后一个继续迫害的人，如果我们相信许多人所倾向的那样，他将要在假基督出现之前立即显现。我的题材会促使我详细陈述他的恶行，若不与这部著作的目的不符，我本可以进入如此巨大的话题。我只满足于说，他各方面都显得极其可憎和残忍，甚至最终到了杀害自己母亲的地步。此后，他还以庄严盟约的方式娶了一个名叫毕达哥拉斯的人，给皇帝蒙上了新娘面纱，而通常的嫁妆、婚床、婚礼火把以及所有其他仪式都一一具备——这些甚至在妇女身上也不能不使人感到羞耻。至于他其他的行为，我不知道描写它们会激起更大的羞耻还是悲伤。他首先试图消灭基督徒的名号，这与恶行总是与美德为敌，以及一切好人常被恶人视为对他们构成谴责的事实相符。因为那时，我们的神圣宗教在城中已广泛传播。伯多禄在那里执行主教的职务，保禄也在被总督不公正判决后上诉于凯撒而被带到罗马。那时，许多人聚集来听保禄，他们因所认识到的真理以及宗徒们当时常常行的奇迹而归向对天主的敬拜。因为那时发生了伯多禄和保禄与西满之间众所周知的著名交锋。西满用他的巫术飞升到空中，由两个恶魔托着（为了证明他是神），但因宗徒们的祈祷，恶魔被赶走，他在全体百姓眼前坠到地上，摔得粉碎。\par

% structure: p0058 source=h2 target=subsection reason=relative-heading-rank; base=h2

\subsection{第二十九章}

与此同时，基督徒的人数已非常庞大，罗马城被大火焚毁，当时\Person{尼禄}正驻跸于\Place{安提乌姆}。但众议将纵火之责归于皇帝，认为他以此谋求建新城的荣耀。事实上，\Person{尼禄}无论如何也无法摆脱是他下令纵火的指控。于是他转而控告基督徒，最残酷的折磨便加诸无辜者身上。不仅如此，还发明了新颖的死刑：他们被披上野兽的皮，被狗吞食而死；许多人被钉十字架或被火烧死；不少人被留作此用，待白昼结束，被焚烧以为夜间照明。就这样，残酷对待基督徒首次显现出来。此后，他们的宗教也被制定法律禁止；公开颁布的法令宣布身为基督徒为非法。那时，\Person{保禄}和\Person{伯多禄}被判处死刑，前者被刀斩首，后者被钉十字架。这些事在\Place{罗马}发生之际，犹太人因无法忍受在\Person{非斯都·弗洛鲁斯}统治下所受的伤害，开始叛乱。\Person{韦斯巴芗}被\Person{尼禄}派去征讨他们，带着代行执政官的权力，在许多重要战役中击败他们，并迫使他们逃入\Place{耶路撒冷}城内。与此同时，\Person{尼禄}因自知罪孽深重而对自己也心生憎恶，从人间消失，留下不确定是否他对自己施以暴力的疑问：他的尸首确实从未被找到。因此，根据关于他所记载的——\Quote{他的致命伤被医好了（\BibleRef{默 13:3}）}——人们相信，即使他确实用刀了结了自己，他的伤口被治愈了，生命得以保全，他将在世界末日临近时再次被差遣，以施行不法的奥秘。\par

% structure: p0060 source=h2 target=subsection reason=relative-heading-rank; base=h2

\subsection{第三十章}

于是，\Person{加尔巴}夺取了政权；不久，\Person{加尔巴}被杀，\Person{奥托}掌握了政权。随后，来自\Place{高卢}的\Person{维特里乌斯}，倚仗他所统率的军队，进入\Place{罗马城}，杀死\Person{奥托}，篡夺了君权。此后政权传给了\Person{韦斯巴芗}，虽然这是通过邪恶手段达成的，但却有从恶人手中解救国度的好效果。当\Person{韦斯巴芗}围攻\Place{耶路撒冷}时，他取得了皇权；按惯例，军队向他欢呼称帝，将王冠戴在他头上。他立其子\Person{提多}为凯撒，并分配给他部分军队，以及继续围攻\Place{耶路撒冷}的任务。\Person{韦斯巴芗}启程返回\Place{罗马}，受到元老院和民众的热烈欢迎；\Person{维特里乌斯}自杀后，他的统治权得以完全巩固。与此同时，犹太人被严密围困，因不得和谈或投降的机会，终致死于饥荒，街道上处处堆满尸体，因埋葬已无法进行。此外，他们开始吃各种最令人厌恶的东西，甚至不排除人肉，除了那些已经腐烂而无法食用的。于是罗马人冲入疲乏的守城者中。碰巧那时来自乡间和\Place{犹太}其他城镇的全体民众正聚集庆祝逾越节：无疑，因为天主乐意使这渎神的种族在一年中他们曾钉死主的同一时刻被交付于毁灭。法利塞人一时间最勇敢地坚守阵地，捍卫圣殿，最终，他们固执求死，自愿投身火海。据记载，死亡人数为一百一十万，被俘而卖为奴者十万。据说\Person{提多}召集会议后，首先考虑是否应当摧毁这座工艺非凡的圣殿。因为有人认为，这座超越一切人类成就的神圣建筑不应被毁，因为若保存，可成为罗马节制的明证，若毁坏，则将成为罗马残暴的永久证据。但另一方面，其他人包括\Person{提多}本人认为，圣殿尤其应当被推翻，以使犹太教和基督教更彻底地被颠覆；因为这两教虽彼此相反，却出自同一源流；基督徒是从犹太人中兴起的；若根被铲除，枝条也必速速枯萎。就这样，按天主旨意，众人的心被激动，圣殿被毁，距今三百三十一年。这圣殿的最后倾覆，以及犹太人最终的被掳，使他们被逐出本土，散见于世界各地，这每日向世界证明，他们受罚别无他因，只因他们对基督伸出了渎神的手。因为虽然在其他时候，他们常因自己的罪恶而被掳，但从未遭受超过七十年的奴役刑罚。\par

% structure: p0062 source=h2 target=subsection reason=relative-heading-rank; base=h2

\subsection{第三十一章}

然后，过了一段时期，\Person{韦斯帕芗}的儿子\Person{多米提安}迫害了基督徒。此时，他将宗徒暨福音作者\Person{若望}放逐到\Place{帕特摩岛}。在那里，隐秘的奥迹向他启示后，他撰写并发表了神圣\WorkTitle{默示录}一书，该书确实被许多人愚蠢地或不敬地拒绝接受。不久之后，\Person{图拉真}统治下发生了第三次迫害。但他在对基督徒施以酷刑和拷打后发现他们无可死或受罚之罪，便禁止再对他们施加任何残暴。随后，在\Person{哈德良}时期，犹太人试图反叛，并企图洗劫\Place{叙利亚}和\Place{巴勒斯坦}；但在派遣军队攻打他们后，他们被镇压了。此时，\Person{哈德良}以为通过损害该地方就能摧毁基督徒信仰，便在圣殿和主受难的地方设立了邪神的偶像。又因为基督徒主要被认为是犹太人（因为那时\Place{耶路撒冷}教会只有受割礼的司祭），他命令一队士兵持续看守，以阻止所有犹太人接近\Place{耶路撒冷}。但这反而有益于基督徒信仰，因为几乎所有当时仍遵守法律的犹太人都信基督为天主。这无疑是由主的上智眷顾所安排，为使法律的奴役从信仰和教会的自由中被除去。这样，那时来自外邦人的\Person{马尔谷}首先成为\Place{耶路撒冷}的主教。第四次迫害被认为发生在\Person{哈德良}时期，但他后来禁止继续进行，宣称任何人未经明确指控就受审是不公正的。\par

% structure: p0064 source=h2 target=subsection reason=relative-heading-rank; base=h2

\subsection{第三十二章。}

\Person{哈德良}之后，在\Person{安敦尼·庇护}统治下，教会享有和平。随后，在\Person{安敦尼}之子\Person{奥勒留}统治下，第五次迫害开始了。那时，首次在\Place{高卢}地区出现了殉道事件，因为天主的宗教在\Place{阿尔卑斯山}以外被接受得稍晚。第六次基督徒迫害发生在\Person{塞维鲁}皇帝统治下。此时，\Person{奥力振}的父亲\Person{莱昂尼达}倾流了他神圣的殉道之血。然后，在三十八年的间歇期间，基督徒享有和平，除了在那段中间时期，\Person{马克西米努}迫害了一些教会的圣职人员。不久，在\Person{德西乌斯}皇帝统治下，第七次血腥迫害爆发了。接下来，\Person{瓦勒良}成为圣徒的第八位敌人。在他之后，大约五十年后，在\Person{戴克里先}和\Person{马克西米安}皇帝统治下，发生了一场最残酷的迫害，持续十年之久，消耗了天主的子民。此时，几乎整个世界都染上了殉道者的神圣鲜血。事实上，他们争先恐后地投入这些光荣的战斗，当时借着光荣的死亡而殉道比现在通过邪恶的野心谋取主教职位更受热切追求。那时世界从未因战争而如此疲惫，我们赢得胜利的凯旋也从未比我们展示无法被十年屠杀征服时更为伟大。当时殉道者的受难记录也被保存下来；但我不认为适合附加这些，以免超出本书的篇幅限制。\par

% structure: p0066 source=h2 target=subsection reason=relative-heading-rank; base=h2

\subsection{第三十三章。}

好吧，迫害的终结发生在八十八年前，那时皇帝们开始成为基督徒。因为\Person{君士坦丁}当时获得了统治权，他是所有罗马统治者中的第一位基督徒。那时，确实，\Person{君士坦丁}争夺帝位的对手\Person{李锡尼}曾命令他的士兵献祭，并将拒绝献祭者逐出军队。但此事不被视为迫害；它太微不足道，无法对教会造成任何伤害。从那时起，我们一直享有安宁；我也不相信会有进一步的迫害，除了在末世之前假基督将要施行的迫害。因为神圣的话语已经宣告，世界将遭受十次灾祸；既然九次已经受过，剩下的那一次必定是最后一次。在这段时期，基督徒宗教如何盛行，真是令人惊奇。因为\Place{耶路撒冷}曾是一片可怕的废墟，如今却装饰着极为众多而宏伟的圣堂。而且，\Person{君士坦丁}皇帝的母亲\Person{赫肋纳}（她作为奥古斯塔与儿子一同统治）怀着强烈的渴望要见到\Place{耶路撒冷}，她推倒了那里的偶像和庙宇；并且随着时间的推移，借着她皇权的行使，她在主受难、复活和升天的地方建造了圣堂。一个显著的事实是，当主被云彩接到天上时，祂的神圣足迹最后留下的那块地方，无法通过铺砌与街道的其余部分连接起来。因为大地不习惯纯粹的人类接触，拒绝了施加其上的所有铺装材料，并常常将大理石块抛回试图铺设它们的人的脸上。此外，该地方的土地曾被天主踏过的一个持久证据是，足迹仍然可见；尽管每天涌向该地方的人们出于信仰争相带走主脚所踏过的土，但该地方的沙土未受损伤；大地仍保持着与古时相同的面貌，仿佛被印在其上的足迹封印了一般。\par

% structure: p0068 source=h2 target=subsection reason=relative-heading-rank; base=h2

\subsection{第三十四章。}

通过同一位皇后的仁慈努力，那时找到了主的十字架。起初，它当然无法被祝圣，因为犹太人反对，后来它又被毁城的瓦砾覆盖。若非有人怀着如此的信德精神去寻找，它至今也不会显露。于是，赫肋纳首先得知了吾主受难的地点，便带了一队士兵到那里，当地全体居民争先恐后地满足皇后的愿望，命令掘开土地，清理所有邻近的大片废墟。不久，作为她信德与劳苦的报酬，发现了三个十字架（正如当年为主和两个强盗所竖立的）。但随即，更大的困难——辨别主曾悬挂的刑架——扰乱了所有人的心思意念，唯恐因凡人易犯的错误，或许会把属于一个强盗的十字架误祝圣为主十字架。于是他们定计，让一个刚死的人接触这些十字架。这个计划立即执行；恰似天主安排，那时正按惯例举行一个死人的葬礼，众人冲上前从棺架上取下尸体。它无效地接触了前两个十字架，但当它触及基督的十字架时——奇妙之极！全体颤抖着——尸体被震落，在众人眼前站立起来。十字架就这样被找到，并按一切应有的礼仪祝圣了。\par

% structure: p0070 source=h2 target=subsection reason=relative-heading-rank; base=h2

\subsection{第三十五章}

赫肋纳完成了这些事；而在一位基督教君王治下，世界既获得了自由，又在他身上拥有了信德的典范。但从那种平静状态中降临到所有教会身上的危险，比之前所有都可怕得多。因为那时亚略异端爆发，以其灌输的谬误搅乱了整个世界。通过两个亚略——这个不忠信仰的最活跃发起者——皇帝本人也被引入歧途；当他自以为在履行宗教义务时，进行了暴力迫害。主教们被流放；对圣职人员施加残暴；甚至那些与亚略派共融分离的平信徒也受罚。亚略派所宣布的教义如下：圣父为创造世界的目的生了圣子；他凭自己的能力从无中造出了一个新的、第二性体，一个新的、第二位天主；曾有一段时间圣子不存在。为应对这个邪恶，从全世界召集了一次会议在尼西亚举行。三百一十八位主教聚集；信德被完全以书面陈述；亚略异端被谴责；皇帝以一道御旨确认了全部。亚略派于是不敢再对正统信仰有任何企图，混入教会中，仿佛默认了已达成的结论，并不持有不同意见。然而，他们心中仍对公教徒怀有根深蒂固的仇恨，用受唆使的控告人和捏造的罪名攻击那些在信仰问题上无法与之争论的人。\par

% structure: p0072 source=h2 target=subsection reason=relative-heading-rank; base=h2

\subsection{第三十六章}

因此，他们首先攻击并缺席谴责了亚历山大的主教亚大纳削，一位圣人，他曾以执事身份出席尼西亚会议。因为他们向假见证人对他堆积的指控上又加了一条，即他怀着恶意接纳了玛尔采禄和福提诺，这两位曾被会议判决谴责的异端司铎。至于福提诺，他被公正地谴责是毫无疑问的。但就玛尔采禄而言，当时似乎没有发现任何值得谴责的事，而对他的无辜的信念尤其因那一派的\OriginalTerm{心态}{animus}而加强，因为没有人怀疑那些谴责他的法官本人就是异端者。但亚略派与其说想除掉这些人，不如说想除掉亚大纳削本人。因此，他们迫使皇帝走到这一步，即亚大纳削应被流放到高卢。但不久，八十位主教在埃及聚集，宣布亚大纳削被不公正地谴责。此事被提交给君士坦丁：他命令来自全世界的主教在萨尔德斯聚集，并且亚大纳削被谴责的整个程序应由该会议重新审议。与此同时，君士坦丁去世，但在他还是皇帝时召集的会议宣告亚大纳削无罪。玛尔采禄也被恢复了他的主教职，但对西尔米翁主教福提诺的判决并未撤销；因为即使在我们朋友的判断中，他也被视为异端者。然而，就连这一结果也使玛尔采禄沮丧，因为福提诺据知是他年轻时的门徒。但这一点也有助于确保亚大纳削无罪开释，即乌尔撒修和瓦伦斯，亚略派中的首领，当他们在萨尔德斯会议后公开与教会的共融分离时，来到罗马主教尤利乌斯面前，因他们曾谴责无辜者而向他请求宽恕，并公开宣布亚大纳削已被萨尔德斯会议的谕令公正地宣告无罪。\par

% structure: p0074 source=h2 target=subsection reason=relative-heading-rank; base=h2

\subsection{第三十七章}

过了若干时间，\Person{亚大纳削}发觉\Person{马塞鲁}在信仰上绝不纯正，便暂停他的共融。\Person{马塞鲁}有着这样的谦逊，当受如此伟大人物的判断责备时，他自动退让了。但尽管早期是无辜的，后来却显然成了异端者，可以断定，当关于他的判决宣布时，他确实是有罪的。亚略派于是找到这种机会，合谋彻底推翻\Place{萨狄克}会议的法令。因为这一事实似乎给了他们一种表面上的正当性：既然对\Person{亚大纳削}有利的判决是不公正地形成的，正如\Person{马塞鲁}被不当开释一样，现在连\Person{亚大纳削}本人也认为他是异端者。因为\Person{马塞鲁}曾作为\OriginalTerm{撒伯流异端}{Sabellian heresy}的支持者挺身而出。但\Person{福提诺}早已提出了一种新的异端，确实与\Person{撒伯流}在神性位格的合一上不同，却宣称\Person{基督}始于\Person{玛利亚}。因此，亚略派以狡猾的设计，将无罪的与有罪的混淆，并在同一判决中纳入对\Person{福提诺}、\Person{马塞鲁}和\Person{亚大纳削}的谴责。他们无疑这样做，是旨在引导无知者的思想得出结论：那些被认为对\Person{马塞鲁}和\Person{福提诺}表达了有根据意见的人，对\Person{亚大纳削}并非判断错误。然而当时亚略派隐藏了他们的背叛；不敢公开宣扬他们的错谬教义，自称为公教徒。他们认为首要目标应是使\Person{亚大纳削}被赶出教会，他一直是阻挡他们努力的壁垒，他们希望若他被移除，其余人就会转向他们的邪恶意见。跟随亚略派的那些主教部分欣然接受了\Person{亚大纳削}的谴责。另一部分被恐惧和派系所迫，屈服于亚略派的愿望；只有少数人，对他们而言真信仰比任何其他考虑更珍贵，拒绝接受他们不义的判决。其中有\Place{特雷维}的主教\Person{保利诺}。据说，当他面前放置了一封关于此事的信时，他如此写道：他同意谴责\Person{福提诺}和\Person{马塞鲁}，但不赞成谴责\Person{亚大纳削}。\par

% structure: p0076 source=h2 target=subsection reason=relative-heading-rank; base=h2

\subsection{第三十八章}

但那时亚略派见诡计未能成功，便决定诉诸武力。因为对那些得到君主宠信、且他们已通过邪恶奉承完全争取过来的人而言，尝试并实施任何事都是容易的。此外，他们在众议上是不可战胜的；因为几乎整个\Place{潘诺尼亚}两省的主教，以及许多东方主教和整个亚洲的主教，都加入了他们的不忠。但那邪恶团伙的首领是\Place{辛吉杜努姆}的\Person{乌尔萨修}、\Place{穆尔萨}的\Person{瓦伦斯}、\Place{赫拉克利亚}的\Person{狄奥多}、\Place{安提约基雅}的\Person{斯特凡}、\Place{凯撒勒雅}的\Person{阿卡修}、\Place{厄弗所}的\Person{默诺凡}、\Place{劳狄刻雅}的\Person{乔治}，以及\Place{尼禄城}的\Person{纳尔基索}。这些人掌控了皇宫到如此地步，以至于皇帝没有他们的同意就不做任何事。他确实听命于他们所有人，但尤其受\Person{瓦伦斯}的影响。因为在当时，当\Place{穆尔萨}与\Person{马格嫩提乌斯}作战时，\Person{君士坦提乌斯}没有勇气亲自下去观战，而是住在一座位于城外的殉道者教堂里，当地主教\Person{瓦伦斯}与他同在，以鼓舞他的勇气。但\Person{瓦伦斯}狡猾地通过他的代理人安排，使他成为第一个知晓战斗结果的人。他这样做，或是为了赢得王恩，若他能首先传递好消息，或是为了保全自己的性命，因为若结果不幸，他将有逃跑的时间。因此，当国王身边的少数人处于惊恐中，皇帝自己也为焦虑所困时，\Person{瓦伦斯}第一个向他们宣布了敌人的溃败。当\Person{君士坦提乌斯}要求将带来消息的人引见时，\Person{瓦伦斯}为了增加对自己的敬畏，说是一位天使来给他传信。皇帝易于相信，日后习惯公开宣称，他是因\Person{瓦伦斯}的功绩而非军队的勇猛赢得了胜利。\par

% structure: p0078 source=h2 target=subsection reason=relative-heading-rank; base=h2

\subsection{第三十九章}

从这第一个证据表明君王已转向他们一边，亚略派人士鼓起勇气，知道他们可以利用国王的权力，而他们自己的权威几乎无法产生影响。因此，当我们的朋友们不接受他们关于\Person{阿塔纳削}的判决时，皇帝颁布了一道法令，规定凡不签署谴责\Person{阿塔纳削}文件的人将被流放。但是，那时，我们的朋友们在位于\Place{高卢}的城镇\Place{阿尔勒}和\Place{比特雷}召开了主教会议。他们要求在任何人被迫签署反对\Person{阿塔纳削}之前，应先就真正的信仰进行讨论；并坚持认为只有在他们就审判者的人选达成一致后，才能就争议的问题作出决定。但是\Person{瓦伦斯}及其同党不敢就信仰进行讨论，首先想通过武力确保对\Person{阿塔纳削}的谴责。由于双方的冲突，\Person{保利努斯}被流放。与此同时，在\Place{米兰}（当时皇帝所在之地）举行了一次集会；但同样的争论在那里继续，没有丝毫缓和。然后，\Place{韦尔切利}的主教\Person{尤西比乌斯}和\Place{撒丁岛}的\Place{卡拉利斯}的主教\Person{路济弗尔}被流放。然而，\Place{米兰}的司铎\Person{迪奥尼修斯}签署了谴责\Person{阿塔纳削}的文件，条件是在主教们之间就真正的信仰进行调查。但是\Person{瓦伦斯}和\Person{乌尔萨修斯}及其同党，由于害怕人民——他们以非凡的热情维护大公信仰——不敢在公开场合提出他们怪异的教义，而是在皇宫内集会。他们从那里以皇帝的名义发布了一封充满各种邪恶的信函，目的无疑是：如果人民给予好评，他们就在公共权威下提出他们想要的东西；但如果受到反对，所有的怨恨都可能指向国王，而他的错误可以被视为可原谅的，因为他当时只是个望教者，很可能在信仰的奥迹上犯了错误。后来，当这封信在教堂里宣读时，人民表达了对它的厌恶。而\Person{迪奥尼修斯}因为不同意他们，被逐出城市，而\Person{奥克森提乌斯}立即被选为接替他的主教。\Place{罗马}城的主教\Person{利贝里乌斯}和\Place{普瓦捷}的主教\Person{希拉留}也被流放。\Place{图卢兹}的主教\Person{罗达尼乌斯}（他天性较温和，与其说是靠自己的力量，不如说是靠与\Person{希拉留}的共融抵抗亚略派）也受到同样的惩罚。然而，所有这些人都准备暂停与\Person{阿塔纳削}的共融，只是为了在主教们之间就真正的信仰进行调查。但亚略派认为最好将最著名的人物从争论中撤走。因此，我们上面提到的人被流放，那是四十五年前，当\Person{阿尔比提乌斯}和\Person{洛利亚努斯}担任执政官的时候。然而，\Person{利贝里乌斯}不久后因\Place{罗马}的骚乱被允许返回城市。但是众所周知，被流放的人受到全世界的钦佩，并且筹集了大量资金满足他们的需要，同时几乎各省的大公人民代表都去探望他们。\par

% structure: p0080 source=h2 target=subsection reason=relative-heading-rank; base=h2

\subsection{第四十章}

在此期间，亚略派不再像以前那样秘密，而是公开宣告他们怪异的异端教义。此外，他们按自己的观点解释尼西亚会议，并通过在其决议中增加一个字母，给真理蒙上一层模糊。因为原文写的是\OriginalTerm{同质}{Homoousion}，意为\Quote{同质}，他们却坚持写的是\OriginalTerm{相似质}{Homoiousion}，这只是\Quote{相似质}的意思。他们因此承认相似，却取消了同一；因为相似与同一大不相同；打个比方，一幅人体的画像可能像一个人，却不具备人的任何实在。但他们中有些人走得更远，主张\OriginalTerm{不同质}{Anomoiousia}，即不相似的本体。这些争论发展到如此地步，以致整个世界都陷入了这些怪异的错误。因为\Person{瓦伦斯}和\Person{乌尔萨修斯}及其支持者（其名字我们已列出）用这些观点感染了\Place{意大利}、\Place{伊利里亚}和\Place{东方}。\Place{阿尔勒}的主教\Person{萨图尔尼努斯}，一个暴烈好斗的人，同样骚扰了我们的\Place{高卢}地区。还有一种普遍的信念，认为来自\Place{西班牙}的\Person{奥西乌斯}也加入了同样不忠的派别，这似乎更加令人惊奇和难以置信，因为他几乎一生都是我们观点最坚定的拥护者，而且尼西亚会议被认为是在他的倡导下召开的。如果他确实倒戈，原因可能是在他极度年老时（因为当时他已一百多岁，如\Person{圣希拉留}在其书信中所述）陷入了昏聩。当世界被这些事情搅动，教会仿佛因一种疾病而憔悴时，一种虽然不那么激动人心但同样严重的忧虑压在皇帝身上：尽管他偏袒的亚略派看似更强大，但主教们之间关于信仰仍没有一致意见。\par

% structure: p0082 source=h2 target=subsection reason=relative-heading-rank; base=h2

\subsection{第四十一章}

\Emph{于是}，皇帝下令在意大利的\Place{阿里米农}城召开主教会议，并指示总督\Person{托鲁斯}，一旦主教们聚集，不得让他们分离，直到他们就同一信德达成一致，同时许诺他若能成功办成此事，便授予他执政官之职。因此，皇室官员被派往\Place{伊利里亚}、意大利、\Place{阿非利加}和两高卢，四百多名西部主教被召集或强制前往\Place{阿里米农}；皇帝下令为所有这些主教提供食物和住宿。但这在我们地区的人，即阿基坦人、高卢人和不列颠人看来，似乎不妥，因此他们拒绝公共供应，宁愿自费生活。不列颠人中只有三人因自身财力不足，接受了公共资助，而拒绝了其他人提供的捐献；因为他们认为，加重国库负担比加重个人负担更合乎义务。我听说我们的主教\Person{加维迪乌斯}惯常以责备的方式提及此事，但我却倾向于完全不同的判断；我钦佩这些主教没有自己的财产，同时又不接受他人的援助而宁愿接受国库的资助，以便不加重任何人的负担。在这两点上，他们为我们树立了高尚的榜样。其他人没有值得记载的事迹；我回到正题。所有主教聚集之后，如上所述，双方分道扬镳。我方占据了教堂，而亚略派则选择了一座当时特意空置的庙宇作为祈祷之所。但他们的人数不超过八十人；其余都属于我方。经过多次会议，实际上毫无进展，我方坚持信德，对方也未放弃他们的不信。最后，决定派遣十位代表去见皇帝，让他了解各方的信德或意见，并使他明白与异端者无法和好。亚略派同样行事，派遣同等数量的代表，在皇帝面前与我方争论。然而，我方选派的是一些学识浅薄、缺乏谨慎的年轻人；而亚略派方面则派遣了老年人，他们技艺娴熟、才华横溢，并且深受他们旧有的不信教义熏陶；这些人很容易在君主面前占了上风。但我方特别受命不得与亚略派有任何形式的共融，并且要将所有问题完整保留，待主教会议讨论。\par

% structure: p0084 source=h2 target=subsection reason=relative-heading-rank; base=h2

\subsection{第42章}

与此同时，在东方，皇帝效仿西方，下令几乎所有主教在\Place{伊苏里亚}的\Place{塞琉西亚}城聚集。那时，\Person{希拉略}已在\Place{弗里吉亚}度过第四年流放生活，被迫与其他主教一同出席，由副官和总督提供公共交通工具。然而，由于皇帝没有对他下达特别命令，法官们只是遵循普通命令（该命令责令他们将所有主教召集到会议），便将他与其他愿意前往的人一同送去。我猜想，这是天主特别的安排，为的是让这位深受神圣事物教导的人能够出席讨论信德的会议。他到达\Place{塞琉西亚}后，受到极大的欢迎，并吸引了所有人的心和感情。他首先询问高卢人的真实信德，因为当时亚略派散布关于我们的恶言，东方人怀疑我们接受了\Person{撒柏流}的观点，即唯一天主的统一仅仅由三重名字区分。但他在阐明了与\Place{尼西亚}教父们所达成的结论一致的信德后，为西方人作了见证。这样，所有人的心都得到了满足，他被接纳共融，并被接纳为同盟，加入了会议。然后他们开始实际工作，发现邪恶异端的始作俑者，并将其从教会团体中分离出来。这些人包括：\Place{亚历山大里亚}的\Person{格奥尔吉乌}、\Person{阿卡基乌}、\Person{欧多克西乌}、\Person{乌拉尼乌}、\Person{雷昂提乌}、\Person{特奥多西乌}、\Person{埃瓦格里乌}、\Person{特奥多卢斯}。但当会议结束后，任命了一个使团去见皇帝，告知他所发生的事。被谴责的人也去了君主那里，依靠他们的同党势力，并因与君主有共同利益而前往。\par

% structure: p0086 source=h2 target=subsection reason=relative-heading-rank; base=h2

\subsection{第43章}

与此同时，皇帝强迫我方从阿里米诺会议派出的那些代表与异端者共融。同时，他交给他们一份由这些恶人草拟的信仰宣认，这份宣认用欺骗性的措辞表达，看似展示天主教信仰，而其中却暗藏不忠信。因为在外表虚假推理的掩盖下，它废除了\OriginalTerm{本体}{Ousia}一词的使用，声称这词含混不清，且被教父过于仓促采纳，并无圣经权威依据。其目的是使圣子不被相信与圣父同一性体。同一份信仰宣认承认圣子与圣父\Emph{相似}。但措辞中精心准备了欺骗，为的是相似而不相等。于是，代表们被遣送走后，命令执政官不得解散会议，直到所有人通过签署表明同意所起草的信仰宣言；若有人过分顽固地推迟，则应被流放，前提是他们的人数不超过十五人。但当代表们返回时，他们被拒绝共融，尽管他们辩称受到了国王施加的压力。因为当所作的决定被发现后，他们的事务和计划中产生了更大的混乱。随后，我方人员逐渐一批批地——部分出于性格软弱，部分鉴于前往外地的疲惫旅途——投降了对方。这些人在代表返回时已成为两派中较强的一方，并占据了教堂，我方的朋友们被逐出。而一旦我方人员的思想开始倾向那边，他们就成群结队地投奔对方，直到我方朋友的人数减少到二十人。\par

% structure: p0088 source=h2 target=subsection reason=relative-heading-rank; base=h2

\subsection{第四十四章}

但是，这些人虽然人数愈少，却显得愈加强大；其中最坚定的是我们的朋友佛加迪乌斯和通格里主教塞尔瓦提乌斯。由于他们没有屈服于威胁和恐惧，塔乌鲁斯便用恳求攻击他们，并流着泪恳求他们采纳更温和的政策。他指出，主教们被关在一个城里已有七个月——他们因冬季的严寒和实际的匮乏而精疲力竭，回家无望；那么结局会是什么？他敦促他们效法多数，至少以先他们而行的众多人数作为权威。因为佛加迪乌斯公开宣称，他已准备好接受流放和任何可能加于他的惩罚，但不会接受亚略派所起草的那份信仰宣认。这样过去了好几天，都在这种讨论中进行。当他们在和解方面进展甚微时，佛加迪乌斯逐渐开始让步，最后被向他提出的一项建议所克服。因为瓦伦斯和乌尔萨提乌斯断言，当前的信仰宣认是按照天主教教义起草的，由东方人在皇帝授意下提出，不能拒绝而不犯不敬之罪；如果一份令东方人满意的宣认却被西方人拒绝，纷争还有什么可能的结果？最后，如果当前宣认中有任何陈述不如意，他们自己应该添加他们觉得该添加的内容，而他们本人则同意那些可能添加的内容。这友好的表白得到众人欣然接受。我方的人也不再敢于反对，因为他们无论如何渴望结束此事。于是，佛加迪乌斯和塞尔瓦提乌斯起草的宣认开始发表；在这些宣认中，首先谴责了亚略及其整个不忠信的计划，同时宣布天主子与圣父相等，且无开端，即没有时间上的起始。随后，瓦伦斯假装协助我方朋友，附加了一项陈述（其中暗藏诡计）：天主子不像其他受造物一样是受造物；这一声明中的欺骗未被听者察觉。因为在这些话中，虽然否认圣子与其他受造物一样，但仍宣布他是一种受造物，只是优于其他。这样，双方都不能认为完全胜利或完全失败，因为宣认本身有利于亚略派，但随后附加的声明有利于我方朋友。不过，必须排除瓦伦斯所附加的那一项，它在当时未被理解，最终明白时为时已晚。无论如何，会议就这样结束了，这个会议开端良好，而结局可耻。\par

% structure: p0090 source=h2 target=subsection reason=relative-heading-rank; base=h2

\subsection{第四十五章}

于是，亚略派人士事务极为兴旺，万事顺遂，便结队前往皇帝所在的君士坦丁堡。他们在那里遇见了塞琉西亚会议的代表，并以王权强迫他们效法西方人士，接受那异端的信德宣认。许多拒绝者被痛苦的监禁和饥饿折磨，最终屈服，良心被俘。但许多更勇敢抵抗者被剥夺主教职，流放，他人取而代之。于是，最好的司铎或因威胁而恐惧，或被流放，少数人的不忠使所有人退缩。\Person{希拉略}当时也在那里，他跟随塞琉西亚的代表；由于对他没有明确命令，他等候皇帝旨意，看是否可能被命令返回流放地。当他察觉信德已陷入极度危险，因为西方人被欺骗，东方人正被邪恶战胜，他公开提交三份文件，请求国王接见，以便在对手面前辩论信德问题。但亚略派竭力反对。最后，\Person{希拉略}被命令返回高卢，被视为散播纷争、扰乱东方者，而对他的流放判决仍未撤销。当他几乎遍历整个被不忠之恶感染的大地时，他心中充满疑虑，深为负担忧虑所扰。他发觉许多人认为不应与承认里米尼会议的人共融，他认为最佳做法是使所有人悔改和革新。他在高卢多次召开会议，几乎全体主教公开承认所犯错误，他谴责里米尼的决议，重新按原始形式制定教会的信德。然而，阿尔勒主教\Person{撒杜尼努斯}，其实是个极恶之人，品性邪恶败坏，抵制这些正确措施。事实上，此人除了身为异端者的恶名，还被判定犯有无数不可言说的罪行，被逐出教会。如此，失去领袖后，反对\Person{希拉略}的势力被瓦解。佩里格的\Person{帕特努斯}也同样狂热，公然宣扬不忠，被驱逐出司铎职；其余人得到宽恕。众所公认，高卢地区从异端之罪中解脱，仅靠\Person{希拉略}的慈善努力。但当时在安提约基雅的\Person{路济弗尔}持迥异观点。他如此谴责里米尼与会者，甚至与那些在悔改或补赎后接纳他们共融的人断绝共融。此决定对错，我不敢断言。\Person{保利诺}和\Person{罗达尼乌斯}死于夫黎基雅；\Person{希拉略}在返回家乡后第六年去世。\par

% structure: p0092 source=h2 target=subsection reason=relative-heading-rank; base=h2

\subsection{第四十六章}

接着是我们自己时代的岁月，既艰难又危险。在此期间，教会遭受非寻常的邪恶玷污，万事陷入混乱。因为，那时西班牙首次发现了臭名昭著的诺斯替派异端——一种以神秘仪式隐藏自己的致命迷信。此祸患的发源地为东方，尤其埃及，但它如何在那里起源并增长，不易解释。\Person{马尔谷}首先将其引入西班牙，他从埃及出发，出生地为孟斐斯。他的学生是某位\Person{阿加佩}，一个出身不凡的女子，以及一位名为\Person{厄尔丕狄乌}的修辞学家。由这些人，\Person{普里西利安}又受教导，他出身贵族，家资丰厚，勇敢、不安、雄辩、博学多读，极善辩论和讨论——事实上，若他没有以邪恶的研究毁掉卓越的才智，他本是个完全幸福的人。无疑，在他身上可见许多心智和身体的优秀品质。他能长期不眠，忍饥耐渴；他少有聚财之欲，用度极为节俭。但同时他非常虚荣，对世俗知识的自满远超应有；此外，人们相信他自幼便行巫术。他自采纳该有害体系后，以他所拥有的说服和谄媚之术，吸引许多贵族和众多平民接受它。除此之外，喜爱新奇且信德不稳、对一切有淫佚好奇心的女子成群结队地投奔他。他在面容和举止中表现出一种谦卑，从而激发众人对他更大的尊敬和敬意，这助长了此种倾向。现在，这异端的消耗性疾病逐渐蔓延至西班牙大部分地区，甚至一些主教也受其败坏影响。其中，\Person{因斯当提乌}和\Person{萨尔维亚努}不仅表示赞同\Person{普里西利安}的观点，甚至还以一种誓言与他绑定。此情形持续，直至科尔多瓦主教\Person{希基诺}（他住在附近）发现实情，将整件事报告给梅里达司铎\Person{依达西乌}。但他无度地骚扰\Person{因斯当提乌}及其同党，超出事态所需，仿佛给渐长的火焰加了火把，以致他激怒了这些恶人而非压制他们。\par

% structure: p0094 source=h2 target=subsection reason=relative-heading-rank; base=h2

\subsection{第四十七章}

因此，在他们之间发生了许多不值一提的争论之后，在萨拉戈萨召开了一次主教会议，甚至连阿基坦尼亚的主教们也出席了。但异端者不敢将自己置于主教会议的审判之下；然而，判决是在他们缺席的情况下作出的，因斯坦提乌斯和萨尔维亚努斯两位主教，以及平信徒赫尔皮狄乌斯和普里西利安被定罪。还附加规定：若有人接纳被定罪者进入共融，他当明白同样的判决也将加于自身。这项职责被委托给索苏巴主教伊塔基乌斯，他负责将主教的法令告知众人，并特别应将希吉努斯排除在共融之外——此人虽曾率先公开反对异端者，后来却可耻地退却，接纳他们进入共融。与此同时，因斯坦提乌斯和萨尔维亚努斯因司铎的判决被定罪，便在阿尔勒城任命一位主教，即普里西利安——他虽是平信徒，却是这一切麻烦的魁首，并在萨拉戈萨主教会议上与二人一同被定罪。他们这样做是为了增强势力，无疑以为若将一个胆大狡猾的人赋予司铎权威，自己就会处于更安全的境地。但随后伊达基乌斯和伊塔基乌斯更热切地推进其措施，以为祸患可在起头就被压制。然而，他们以不智的谋划，求助于世俗法官，期冀通过其法令和起诉将异端者逐出城市。因此，在多次丢脸的争执之后，经伊达基乌斯的恳求，从当时在位的皇帝格拉提安处获得一道敕令，命令所有异端者不仅离开教堂或城市，而且要被驱逐出他管辖下的所有领土。这道敕令颁布后，\OriginalTerm{诺斯替派}{Gnostics}对自己的事缺乏信心，不敢反对判决，但那些被称为主教的人自行退让，其余的人则因恐惧而四散。\par

% structure: p0096 source=h2 target=subsection reason=relative-heading-rank; base=h2

\subsection{第四十八章。}

于是，因斯坦提乌斯、萨尔维亚努斯和普里西利安启程前往罗马，为要在当时担任罗马主教的达马苏斯面前洗清对他们的指控。他们的旅程穿越阿基坦尼亚的中心地带，在那里受到那些不明真相的人大张旗鼓的接待，于是播下了他们异端的种子。他们尤其以邪恶的教训败坏了厄鲁萨的人民，这些人原本具有良好的宗教倾向。他们被德尔菲努斯逐出波尔多，但在欧克罗提亚的领地稍作停留时，又以他们的错误感染了一些人。随后，他们继续已开始的旅程，陪伴他们的是一群卑劣可耻的人，其中有他们的妻子，甚至还有陌生妇女。在这些妇女中，有欧克罗提亚及其女儿普罗库拉，后者有传言说：她因与普里西利安通奸怀孕，用某些植物堕胎。当他们抵达罗马，希望在达马苏斯面前洗清自己时，甚至未被允许觐见。回到米兰后，他们发现安波罗修同样反对他们。于是他们改变策略，想方设法——既然未能说服当时拥有最高权威的两位主教——通过贿赂和谄媚从皇帝那里获得他们想要的东西。因此，他们买通了公务总管马其顿尼乌斯，获得一道敕令，将先前所颁布的法令践踏于脚下，命令恢复他们的教会。依靠这道敕令，因斯坦提乌斯和普里西利安返回西班牙（因为萨尔维亚努斯已死于罗马城）；于是，他们毫无争斗地收回了曾统治过的那些教会。\par

% structure: p0098 source=h2 target=subsection reason=relative-heading-rank; base=h2

\subsection{第四十九章。}

但伊塔丘斯所缺乏的是力量，而非意愿去抵抗；因为异端者用贿赂收买了总督沃伦提乌斯，从而巩固了他们的势力。此外，伊塔丘斯被这些人以扰乱教会的罪名送上审判，经过一次激烈的起诉后，他被下令当作囚犯带走，恐惧中逃往高卢，投奔了总督格雷戈里。格雷戈里得知所发生的一切后，命令将这场骚乱的制造者带到他面前，并将所有发生的事向皇帝报告，以便断绝异端者一切谄媚或贿赂的门路。但这事徒劳无功，因为由于少数人的放纵和权势，那里一切都可以用钱买到。于是，异端者用诡计，向马其顿尼乌斯献上一大笔钱，确保凭借皇帝的权威，案件的审理从总督手中被夺走，转交给西班牙的副将。那时，西班牙已不再由总督统治，而是由\Quote{长官}派官员将当时住在特里夫斯的伊塔丘斯带回西班牙。然而，他狡猾地逃脱了他们，后来得到普里塔尼乌斯主教的保护，便对他们置之不理。那时，还流传着一个模糊的谣言，说马克西穆斯在不列颠夺取了皇权，并将很快入侵高卢。因此，伊塔丘斯决定，尽管他的处境岌岌可危，还是要等待新皇帝的到来；在此期间，他不采取任何行动。于是，当马克西穆斯作为胜利者进入特里维里城时，伊塔丘斯满怀恶意和控诉，向马克西穆斯倾吐了对普里西利安及其同伙的指控。皇帝受这些陈述的影响，写信给高卢总督和西班牙副将，命令将所有受这可耻异端影响的人带到波尔多的主教会议。于是，因斯但提乌斯和普里西利安被押送到那里；其中，因斯但提乌斯被要求为自己的案件辩护；他被发现无法为自己开脱后，被宣布不配担任主教之职。但普里西利安为了逃避主教的审判，上诉到皇帝。而这之所以被允许，是因为我们这边的人缺乏决心；他们本应判决即使反抗的人，或者如果他们认为自己也可疑，就应保留案件给其他主教审理，而不应将如此明显的罪行案件移交给皇帝。\par

% structure: p0100 source=h2 target=subsection reason=relative-heading-rank; base=h2

\subsection{第五十章}

于是，所有被这诉讼所波及的人都被带到皇帝面前。主教伊达丘斯和伊塔丘斯作为控诉者跟随；我绝不责备他们消灭异端者的热忱，如果他们不是为了取胜而斗争得比适当更激烈的话。而且我的感觉确实是，控诉者与被控诉者一样令我反感。我确实认为伊塔丘斯毫无价值或圣德可言。因为他是一个大胆、多话、厚颜无耻、奢侈浪费的人，过度沉溺于肉欲的快乐。他甚至愚蠢到指控所有那些以阅读为乐或彼此争相守斋的人——无论他们多么圣洁——为普里西利安的朋友或门徒。这个可怜的恶棍甚至公开对马丁提出了一项可耻的异端指控，而马丁当时是一位主教，显然堪与宗徒相比。因为马丁当时住在特里夫斯，不住地恳求伊塔丘斯放弃他的指控，或恳求马克西穆斯不要流那些不幸者的血。他主张，主教们已经判决他们是异端者，将他们逐出教会作为惩罚已经足够；此外，让世俗统治者审理教会案件是一种丑陋且前所未闻的侮辱。事实上，只要马丁还活着，审判就被推迟；而当他即将离开这个世界时，凭借他非凡的影响力，他从马克西穆斯那里得到承诺，不会对罪犯采取残酷的措施。但后来，皇帝被马格努斯和鲁弗斯引入歧途，偏离了马丁所建议的温和路线，将案件委托给总督埃沃迪乌斯，一个严厉苛刻的人。埃沃迪乌斯在两次集会上审判了普里西利安，并判定他有邪恶行为。事实上，普里西利安并没有否认他沉溺于淫荡的教义，惯于在夜间聚集下贱的女人，并赤身裸体地祈祷。因此，埃沃迪乌斯判定他有罪，将他送回监狱，直到他有时机咨询皇帝。于是，事情的详细情况都呈报给皇宫，皇帝下旨处死普里西利安和他的朋友们。\par

% structure: p0102 source=h2 target=subsection reason=relative-heading-rank; base=h2

\subsection{第五十一章}

然而，\Person{依塔纠}见他在最后审判中若再以原告身份提出死刑控诉（因审判必须重审），必会在主教中激起极大反感，遂撤回控诉。但他如此行事之诡计徒劳无功，因为祸害已成。\newline{}果然，\Person{玛克西默}任命了一位名叫\Person{帕特利纠}的、与国库相关的辩护人作为原告。于是，在此人的起诉下，\Person{普里西利安}被判处死刑，与他一同被处死的还有\Person{斐利奇西默}和\Person{阿尔梅纽}，这两人不久前还是神职人员，却加入了\Person{普里西利安}的阵营，背离了天主教徒。\Person{拉特洛尼安}和\Person{欧克罗提亚}也被斩首。\Person{因斯坦提}，如我们上文所述，已被主教们定罪，被流放到不列颠以外的\Place{西利纳岛}。随后，其他人在后续的审判中也受到了追究：\Person{阿撒里沃}和执事\Person{奥勒留}被判处斩首；\Person{提贝里安}被剥夺财产，流放到\Place{西利纳岛}。\Person{特图洛}、\Person{波塔米}和\Person{若望}，因其地位较低且值得怜悯——他们在审判前就已供认自身及其同谋——被判处暂时流放高卢。就这样，这些最不配见天日的人，为了给他人作为可怕的儆戒，或被处死，或被流放。\newline{}\Person{依塔纠}最初凭借上诉权和对公共利益的关注来辩护这一行动，但后来在不断的抨击中终于屈服，将罪责推给那些指导并策动他达成目标的人。然而，在所有这些人中，唯独他被逐出了主教职位。因为\Person{依达纠}虽罪责较轻，但已自愿辞去主教之职：此举明智而可敬，只可惜他后来试图恢复失去的地位，从而败坏了这一步骤的名声。\newline{}普里西利安死后，不仅异端没有被他这位始作俑者所遏制，反而获得力量，传播得更广。他的追随者先前尊他为圣人，此后便开始把他当作殉道者来敬仰。被处决者的遗体被运往西班牙，葬礼极为隆重。甚至，人们认为以\Person{普里西利安}的名义发誓是最高级的宗教行为。但他们与我们之间的争论之战持续不断。这场冲突以可怕的纷争持续了十五年，无论如何都无法平息。如今，一切事物都因不和而动荡混乱，尤其是主教们之间的不和；他们因仇恨、偏袒、恐惧、不忠、嫉妒、结党、贪欲、傲慢、昏睡和懒惰，败坏了一切。总之，大多数人都以疯狂的图谋和顽固执拗对抗少数明智建言者；与此同时，天主的子民和世间贤良都遭到嘲弄和侮辱。\par

\SourceRecord{Translated by Alexander Roberts.；Nicene and Post-Nicene Fathers, Second Series；Vol. 11.；Edited by Philip Schaff and Henry Wace.；Buffalo, NY: Christian Literature Publishing Co.,；1894.；<http://www.newadvent.org/fathers/35052.htm>.}
